\documentclass[
  justified,
  openright,
]{tufte-book}

% --------------------------------------------------
% Engine sanity (XeLaTeX or LuaLaTeX required)
% --------------------------------------------------
\usepackage{iftex}
\ifPDFTeX
  \PackageError{template}{This template requires XeLaTeX or LuaLaTeX}{}
\fi

% --------------------------------------------------
% Core math
% --------------------------------------------------
\usepackage{amsmath,amssymb,amsfonts}

% --------------------------------------------------
% Tufte: subsubsection is forbidden --- remap safely
% --------------------------------------------------
\makeatletter
\let\subsubsection\paragraph
\let\paragraph\paragraph
\makeatother

% --------------------------------------------------
% Fix Pandoc tabularray counter=none bug
% --------------------------------------------------
\makeatletter
\ExplSyntaxOn
\cs_if_exist:NF \c@none { \newcounter{none} }
\ExplSyntaxOff
\makeatother

\usepackage{fontspec}

\IfFontExistsTF{TeX Gyre Pagella}
  {\setmainfont{TeX Gyre Pagella}}
  {\setmainfont{Latin Modern Roman}}

\IfFontExistsTF{TeX Gyre Heros}
  {\setsansfont[Scale=0.92]{TeX Gyre Heros}}
  {\setsansfont{Latin Modern Sans}}

\IfFontExistsTF{TeX Gyre Cursor}
  {\setmonofont[Scale=0.82]{TeX Gyre Cursor}}
  {\setmonofont{Latin Modern Mono}}
% --------------------------------------------------
% Graphics and tables
% --------------------------------------------------
\usepackage{graphicx}
\usepackage{booktabs}
\usepackage{longtable}
\usepackage{array}
\usepackage{tabularray}
\UseTblrLibrary{booktabs}

\setkeys{Gin}{width=\maxwidth,height=\maxheight,keepaspectratio}

% --------------------------------------------------
% Colors
% --------------------------------------------------
\usepackage[table,xcdraw,svgnames]{xcolor}

\definecolor{DarkGray}{RGB}{55,60,70}
\definecolor{CalloutBlue}{HTML}{DCEEFF}
\definecolor{CalloutBorder}{HTML}{5B8BD9}
\definecolor{BrandBlue}{RGB}{18,65,124}
\definecolor{BrandAccent}{RGB}{0,153,204}
\definecolor{LightGray}{RGB}{245,246,248}
\definecolor{MidGray}{RGB}{175,180,190}
\definecolor{WarningAmber}{RGB}{230,100,0}
\definecolor{WarningBg}{RGB}{255,248,235}
\definecolor{CautionYellow}{RGB}{180,140,0}
\definecolor{CautionBg}{RGB}{255,252,230}
\definecolor{TipGreen}{RGB}{36,110,44}
\definecolor{TipBg}{RGB}{237,248,239}
\definecolor{NoteBg}{RGB}{235,243,255}
\definecolor{NoteBlue}{RGB}{18,90,160}
\definecolor{DangerRed}{RGB}{185,28,28}
\definecolor{DangerBg}{RGB}{255,237,237}
\definecolor{TableHeader}{RGB}{18,65,124}
\definecolor{TableRowAlt}{RGB}{235,241,252}
\definecolor{MarginRule}{RGB}{210,215,225}
\definecolor{SpineAccent}{RGB}{0,153,204}
% --------------------------------------------------
% Pandoc syntax highlighting macros
% --------------------------------------------------
\newcommand{\AlertTok}[1]{\textcolor{Red}{#1}}
\newcommand{\AnnotationTok}[1]{\textcolor{Gray}{#1}}
\newcommand{\AttributeTok}[1]{\textcolor{DarkSlateGray}{#1}}
\newcommand{\BaseNTok}[1]{\textcolor{Purple}{#1}}
\newcommand{\BuiltInTok}[1]{#1}
\newcommand{\CharTok}[1]{\textcolor{Teal}{#1}}
\newcommand{\CommentTok}[1]{\textcolor{MidGray}{#1}}
\newcommand{\CommentVarTok}[1]{\textcolor{MidGray}{#1}}
\newcommand{\ConstantTok}[1]{\textcolor{Purple}{#1}}
\newcommand{\ControlFlowTok}[1]{\textcolor{BrandBlue}{#1}}
\newcommand{\DataTypeTok}[1]{\textcolor{Purple}{#1}}
\newcommand{\DecValTok}[1]{\textcolor{Purple}{#1}}
\newcommand{\DocumentationTok}[1]{\textcolor{MidGray}{#1}}
\newcommand{\ErrorTok}[1]{\textcolor{Red}{#1}}
\newcommand{\ExtensionTok}[1]{#1}
\newcommand{\FloatTok}[1]{\textcolor{Purple}{#1}}
\newcommand{\FunctionTok}[1]{\textcolor{BrandBlue}{#1}}
\newcommand{\ImportTok}[1]{\textcolor{BrandAccent}{#1}}
\newcommand{\InformationTok}[1]{\textcolor{BrandAccent}{#1}}
\newcommand{\KeywordTok}[1]{\textcolor{BrandBlue}{#1}}
\newcommand{\NormalTok}[1]{#1}
\newcommand{\OperatorTok}[1]{\textcolor{DarkGray}{#1}}
\newcommand{\OtherTok}[1]{#1}
\newcommand{\PreprocessorTok}[1]{\textcolor{BrandAccent}{#1}}
\newcommand{\RegionMarkerTok}[1]{#1}
\newcommand{\SpecialCharTok}[1]{\textcolor{Teal}{#1}}
\newcommand{\SpecialStringTok}[1]{\textcolor{Teal}{#1}}
\newcommand{\StringTok}[1]{\textcolor{Teal}{#1}}
\newcommand{\VariableTok}[1]{#1}
\newcommand{\VerbatimStringTok}[1]{\textcolor{Teal}{#1}}
\newcommand{\WarningTok}[1]{\textcolor{WarningAmber}{#1}}
% --------------------------------------------------
% Ensure ALL Pandoc *Tok macros are defined (robust)
% --------------------------------------------------
\makeatletter
\@for\tok:=Alert,Annotation,Attribute,BaseN,BuiltIn,Char,Comment,CommentVar,%
Constant,ControlFlow,DataType,DecVal,Documentation,Error,Extension,Float,%
Function,Import,Information,Keyword,Normal,Operator,Other,Preprocessor,%
RegionMarker,SpecialChar,SpecialString,String,Variable,VerbatimString,Warning%
\do{%
  \expandafter\providecommand\csname\tok Tok\endcsname[1]{#1}%
}
\makeatother

% --------------------------------------------------
% Pandoc code block support (Shaded environment)
% --------------------------------------------------
\usepackage{fancyvrb}
\usepackage{framed}

\definecolor{shadecolor}{RGB}{248,248,248}

\newenvironment{Shaded}{%
  \begin{snugshade}
}{%
  \end{snugshade}
}

% --------------------------------------------------
% Pandoc Highlighting environment
% --------------------------------------------------
\DefineVerbatimEnvironment{Highlighting}{Verbatim}{
  commandchars=\\\{\},
  fontsize=\small
}

% --------------------------------------------------
% Tufte: subsubsection is forbidden --- remap safely
% --------------------------------------------------
\makeatletter
\let\subsubsection\paragraph
\let\paragraph\paragraph
\makeatother

% --------------------------------------------------
% Utilities
% --------------------------------------------------
\usepackage{etoolbox}
\usepackage{xparse}
\usepackage{ifthen}
\usepackage{textcase}
\usepackage{titlesec}
% --------------------------------------------------
% Fix Pandoc tabularray counter=none bug
% --------------------------------------------------
\makeatletter
\ExplSyntaxOn
\cs_if_exist:NF \c@none { \newcounter{none} }
\ExplSyntaxOff
\makeatother
% --------------------------------------------------
% Paragraph & section spacing
% --------------------------------------------------
\setlength{\parindent}{1em}
\setlength{\parskip}{0pt}

\titlespacing*{\section}{0pt}{2.0ex plus .2ex}{1.0ex}
\titlespacing*{\subsection}{0pt}{1.5ex plus .2ex}{0.8ex}

% --------------------------------------------------
% Margin notes (tufte)
% --------------------------------------------------
\setlength{\marginparwidth}{2.2in}

\newcommand{\mn}[1]{%
  \marginnote{\raggedright\footnotesize\color{DarkGray}#1}
}

\NewDocumentCommand{\margincallout}{m}{%
  \marginnote{%
    \setlength{\fboxsep}{6pt}%
    \colorbox{CalloutBlue}{%
      \begin{minipage}{\marginparwidth}
      \footnotesize #1
      \end{minipage}
    }
  }
}

% --------------------------------------------------
% Soul / tabularray / marking fixes
% --------------------------------------------------
\usepackage{soul}
\AfterEndEnvironment{longtblr}{\soulsuspend}

\makeatletter
\providecommand{\soulsuspend}{}
\providecommand{\soulsuspend*}{}
\providecommand{\soulsuspend@}{}
\makeatother

\DeclareRobustCommand{\checkbox}{\(\square\)}

\soulregister\checkbox{0}
\soulregister\textbf{1}
\soulregister\emph{1}
\soulregister\href{2}
\soulregister\label{0}
\soulregister\ref{0}
\soulregister\pageref{0}
\soulregister\footnote{1}
\soulregister\cite{0}

% --------------------------------------------------
% Headers / marks (no case mangling)
% --------------------------------------------------
\makeatletter
\renewcommand{\chaptermark}[1]{\markboth{#1}{}}
\renewcommand{\sectionmark}[1]{\markright{#1}}

\pretocmd\markboth{\let\allcapsspacing\relax\let\smallcapsspacing\relax}{}{}
\pretocmd\markright{\let\allcapsspacing\relax\let\smallcapsspacing\relax}{}{}
\makeatother

% --------------------------------------------------
% Blank page helper
% --------------------------------------------------
\newcommand{\blankpage}{%
  \clearpage
  \thispagestyle{empty}
  \null\vfill
  \begin{center}
    \color{MidGray}\sffamily\small\itshape
    (This page intentionally left blank.)
  \end{center}
  \vfill
  \clearpage
}

% --------------------------------------------------
% Fancyhdr
% --------------------------------------------------
\usepackage{fancyhdr}
\pagestyle{fancy}
\fancyhf{}
\fancyhead[LE]{\small\scshape\leftmark}
\fancyhead[RO]{\small\scshape\rightmark}
\fancyfoot[C]{\thepage}
\AtBeginDocument{\pagestyle{fancy}}

% --------------------------------------------------
% Hyperlinks
% --------------------------------------------------
\usepackage{hyperref}
\hypersetup{
  colorlinks=true,
  linkcolor=blue,
  urlcolor=blue,
  citecolor=blue,
  hypertexnames=false
}

% --------------------------------------------------
% Pandoc helpers
% --------------------------------------------------
\providecommand{\tightlist}{%
  \setlength{\itemsep}{0pt}\setlength{\parskip}{0pt}}

% Show only parts and chapters in TOC; number only parts
\setcounter{tocdepth}{1}
\setcounter{secnumdepth}{-1}

% --------------------------------------------------
% Metadata
% --------------------------------------------------



% --------------------------------------------------
% ============================================================
%  DOCUMENT
% ============================================================
\begin{document}
% -- Front matter -------------------------------------------
\frontmatter
\pagestyle{fancy}

% 1) Blank half title page
\thispagestyle{empty}
\null
\clearpage

% 2) Frontispiece (placeholder)
\thispagestyle{empty}
\begin{center}
  \vspace*{0.35\textheight}
  {\sffamily\large Frontispiece}
\end{center}
\clearpage

% 3) Full title page
\begin{titlepage}
    \vspace*{1.6in}
    {\sffamily\fontsize{9}{11}\selectfont\MakeUppercase{Training Manual}}
    \vspace{0.3in}
    {\color{BrandAccent}\rule{4.5in}{2pt}}
    \vspace{0.25in}
    % ADD YOUR TITLE TEXT HERE:
    {\sffamily\bfseries\fontsize{32}{38}\selectfont Coding With //ScreenReaders}
    \vspace{0.2in}
    {\color{BrandAccent}\rule{4.5in}{0.6pt}}
    \vspace{0.5in}
    \vfill
    {\sffamily\fontsize{7.5}{9}\selectfont\color{MidGray}%
      CONFIDENTIAL --- For Internal Use Only%
    }
  \clearpage
  \nopagecolor
\end{titlepage}

% 4) Copyright page (placeholder)
\thispagestyle{empty}
\vspace*{0.25\textheight}
{\small
\noindent Copyright \textcopyright\ \the\year\ Michael Ryan Hunsaker, M.Ed., Ph.D. Coding With ScreenReaders is licensed under a Creative Commons Attribution-NoNCommerical-NoDerivatives 4.0 International: https://creativecommons.org/licenses/by-nc-nd/4.0/
\noindent All rights reserved.
\noindent NO AI TRAINING: Without in any way limiting the author's rights under copyright, any use of this publication to “train” generative artificial intelligence (AI) technologies to generate text is expressly prohibited.
}
\clearpage

% 5) Contents
{
\hypersetup{linkcolor=BrandBlue}
\tableofcontents
\clearpage
}

% 6) Dedication (placeholder)
\chapter*{Dedication}
\thispagestyle{empty}
\begin{center}
  \vspace*{0.3\textheight}
  {\itshape To the countless individuals who have been told they lack the capability to code, yet remain determined to learn and create through accessible coding practices, this book is dedicated.}\\[1em]
\end{center}
\clearpage

% 7) Epigraph (placeholder)
\chapter*{Epigraph}
\thispagestyle{empty}
\begin{flushright}
  \vspace*{0.25\textheight}
  {\itshape The power of the Web is in its universality. Access by everyone regardless of disability is an essential aspect}\\[1em]
  --- Tim Berners-Lee, Inventor of the World Wide Web
\end{flushright}
\clearpage

% 8) Introduction (placeholder)
\newpage

% -- Main content -------------------------------------------
\mainmatter
\pagenumbering{arabic} % Switch to numbers (1, 2, 3...) for the main content


\chapter*{Introduction to Computer Science with JavaScript}

\subsubsection{For Visually Impaired
Students}\label{for-visually-impaired-students}

\subsection{Course Overview}\label{course-overview}

This curriculum introduces computer science concepts using JavaScript,
focusing on Node.js for server-side/terminal execution. All exercises
are text-based and terminal-focused --- no browser required for core
concepts.

\textbf{Target Audience:} Students with no prior programming experience
using screen readers (NVDA, JAWS, VoiceOver, or Windows Narrator).

\begin{center}\rule{0.5\linewidth}{0.5pt}\end{center}

\subsection{Why JavaScript?}\label{why-javascript}

JavaScript is one of the most widely used programming languages in the
world. It powers:

\begin{itemize}
\tightlist
\item
  Web development (front-end and back-end)
\item
  Server applications (Node.js)
\item
  Mobile apps (React Native)
\item
  Desktop apps (Electron)
\end{itemize}

Learning JavaScript opens doors to a vast ecosystem of tools and
opportunities.

\begin{center}\rule{0.5\linewidth}{0.5pt}\end{center}

\subsection{Tools Required}\label{tools-required}

\begin{itemize}
\tightlist
\item
  \textbf{Node.js 18 LTS+} --- \url{https://nodejs.org/}
\item
  \textbf{Visual Studio Code} with the following extensions:

  \begin{itemize}
  \tightlist
  \item
    JavaScript (ES6) code snippets
  \item
    ESLint
  \item
    Prettier
  \end{itemize}
\item
  A screen reader (NVDA, JAWS, or VoiceOver)
\end{itemize}

\subsection{Accessibility Notes}\label{accessibility-notes}

\begin{itemize}
\tightlist
\item
  All program output is text-based (Node.js terminal/REPL).
\item
  JavaScript uses curly braces \texttt{\{\}} for code blocks, like Java.
\item
  VS Code's Screen Reader Optimized Mode works well with JavaScript.
\item
  Use \texttt{Ctrl\ +\ Shift\ +\ P} -\textgreater{} ``Toggle Screen
  Reader Optimized'' in VS Code.
\item
  Use \texttt{F8} / \texttt{Shift+F8} to navigate between errors.
\end{itemize}

\begin{center}\rule{0.5\linewidth}{0.5pt}\end{center}

\subsection{Units}\label{units}

{\def\LTcaptype{none} % do not increment counter
\begin{longtable}[]{@{}ll@{}}
\toprule\noalign{}
Unit & Title \\
\midrule\noalign{}
\endhead
\bottomrule\noalign{}
\endlastfoot
Unit 0 & Setting Up Your JavaScript Development Environment \\
Unit 1 & Foundations --- Getting Started with JavaScript \\
Unit 2 & Numbers, Input, and the Math Object \\
Unit 3 & Control Flow --- Decisions and Loops \\
Unit 4 & Data Structures --- Arrays and Objects \\
Unit 5 & User Input in Node.js \\
Unit 6 & Functions \\
Unit 7 & Classes and Object-Oriented Programming \\
Unit 8 & Files and Error Handling \\
Unit 9 & Modules, Collections, and Testing \\
Unit 10 & Project-Based Learning \\
Unit 11 & Capstone Project \\
\end{longtable}
}

\begin{center}\rule{0.5\linewidth}{0.5pt}\end{center}

\subsection{References}\label{references}

\begin{itemize}
\tightlist
\item
  MDN Web Docs:
  \url{https://developer.mozilla.org/en-US/docs/Web/JavaScript}
\item
  Eloquent JavaScript (free online):
  \url{https://eloquentjavascript.net/}
\item
  Node.js Documentation: \url{https://nodejs.org/docs/}
\end{itemize}

\newpage

\section{Teacher's Guide: Introduction to Computer Science with
JavaScript}\label{teachers-guide-introduction-to-computer-science-with-javascript}

\subsection{For Students with Visual Impairments --- Complete Instructor
Edition}\label{for-students-with-visual-impairments-complete-instructor-edition}

\begin{center}\rule{0.5\linewidth}{0.5pt}\end{center}

\subsection{Table of Contents}\label{table-of-contents}

\begin{enumerate}
\def\labelenumi{\arabic{enumi}.}
\tightlist
\item
  \hyperref[philosophy-and-why-javascript]{Philosophy and Why
  JavaScript}
\item
  \hyperref[javascript-specific-accessibility-challenges]{JavaScript-Specific
  Accessibility Challenges}
\item
  \hyperref[classroom-setup-and-configuration]{Classroom Setup and
  Configuration}
\item
  \hyperref[spoken-javascript-conventions]{Spoken JavaScript
  Conventions}
\item
  \hyperref[unit-01-setup-and-foundations]{Unit 0--1: Setup and
  Foundations}
\item
  \hyperref[unit-2-numbers-input-and-math]{Unit 2: Numbers, Input, and
  Math}
\item
  \hyperref[unit-3-control-flow]{Unit 3: Control Flow}
\item
  \hyperref[unit-4-arrays-and-objects]{Unit 4: Arrays and Objects}
\item
  \hyperref[unit-5-strings]{Unit 5: Strings}
\item
  \hyperref[unit-6-functions]{Unit 6: Functions}
\item
  \hyperref[unit-7-classes-and-oop]{Unit 7: Classes and OOP}
\item
  \hyperref[unit-8-files-modules-and-error-handling]{Unit 8: Files,
  Modules, and Error Handling}
\item
  \hyperref[unit-9-collections-analysis-and-testing]{Unit 9:
  Collections, Analysis, and Testing}
\item
  \hyperref[units-1011-projects-and-capstone]{Units 10--11: Projects and
  Capstone}
\item
  \hyperref[bridging-from-pythonjava-to-javascript]{Bridging from
  Python/Java to JavaScript}
\item
  \hyperref[assessment-strategies]{Assessment Strategies}
\item
  \hyperref[supporting-diverse-needs]{Supporting Diverse Needs}
\item
  \hyperref[alternative-and-supplementary-resources]{Alternative and
  Supplementary Resources}
\item
  \hyperref[quick-reference-sheets]{Quick Reference Sheets}
\end{enumerate}

\begin{center}\rule{0.5\linewidth}{0.5pt}\end{center}

\subsection{1. Philosophy and Why
JavaScript}\label{philosophy-and-why-javascript}

\subsubsection{The Case for Teaching JavaScript to Blind
Students}\label{the-case-for-teaching-javascript-to-blind-students}

JavaScript has a unique profile of benefits and challenges for visually
impaired learners:

\textbf{Universal relevance.} JavaScript is the only language that runs
natively in every web browser. A student who learns JavaScript can
eventually build tools that improve their own accessibility experience
--- custom browser extensions, accessible widgets, automation scripts.
This self-advocacy potential is meaningful.

\textbf{Node.js enables accessible terminal programs.} While JavaScript
began in browsers, Node.js allows running JavaScript programs in the
terminal --- the same accessible text-based environment used for Python.
This curriculum focuses on Node.js to keep all I/O in the terminal where
screen readers excel.

\textbf{Asynchronous programming is a real-world skill.} JavaScript's
\texttt{async}/\texttt{await} model reflects how most modern software
actually works --- programs that wait for I/O, network requests, and
user events. Learning this early prepares students for professional
development.

\textbf{The npm ecosystem.} Over 2 million packages on npm provide
pre-built solutions for almost any problem. This is a double-edged
concern for beginners (dependency complexity), but the ability to say
\texttt{npm\ install} and gain powerful functionality is professionally
relevant.

\textbf{JSON is native.} \texttt{JSON.parse()} and
\texttt{JSON.stringify()} are built into JavaScript with no import
required. For data persistence projects, this reduces friction.

\subsubsection{An Honest Account of
Challenges}\label{an-honest-account-of-challenges}

\textbf{JavaScript's permissiveness is a silent-failure risk.} Unlike
Python (which raises exceptions for type mismatches) and Java (which
rejects them at compile time), JavaScript silently converts types.
\texttt{"5"\ +\ 3} = \texttt{"53"}. \texttt{"5"\ -\ 3} = \texttt{2}.
\texttt{0\ ==\ false} = \texttt{true}. Without visual scanning of
intermediate results, these silent failures can propagate through a
program undetected. This must be addressed through defensive programming
habits taught from day one.

\textbf{Asynchronous programming is conceptually hard.} The
\texttt{async}/\texttt{await} pattern required for user input is a
genuine hurdle for beginners. There is no way to eliminate it from a
Node.js input-based curriculum. The approach in this guide is to
scaffold it carefully and defer full explanation of the event loop to
later units.

\textbf{Multiple valid syntaxes.} Function declarations, function
expressions, arrow functions, methods --- all do essentially the same
thing. \texttt{var}, \texttt{let}, \texttt{const} --- three variable
declarations. \texttt{==} and \texttt{===} --- two equality operators.
This variety, which professionals call ``JavaScript fatigue,'' creates
unnecessary cognitive overhead for beginners. The approach in this
curriculum is to teach one style consistently
(\texttt{const}/\texttt{let}, arrow functions, \texttt{===}) and
acknowledge alternatives exist without teaching them in parallel.

\begin{center}\rule{0.5\linewidth}{0.5pt}\end{center}

\subsection{2. JavaScript-Specific Accessibility
Challenges}\label{javascript-specific-accessibility-challenges}

\subsubsection{Silent Type Coercion}\label{silent-type-coercion}

This is the greatest accessibility risk in JavaScript for blind
students. Unlike Python's \texttt{TypeError} and Java's compile-time
type check, JavaScript silently converts types in ways that can produce
wrong answers with no error:

{\def\LTcaptype{none} % do not increment counter
\begin{longtable}[]{@{}
  >{\raggedright\arraybackslash}p{(\linewidth - 6\tabcolsep) * \real{0.2043}}
  >{\raggedright\arraybackslash}p{(\linewidth - 6\tabcolsep) * \real{0.0753}}
  >{\raggedright\arraybackslash}p{(\linewidth - 6\tabcolsep) * \real{0.2366}}
  >{\raggedright\arraybackslash}p{(\linewidth - 6\tabcolsep) * \real{0.4839}}@{}}
\toprule\noalign{}
\begin{minipage}[b]{\linewidth}\raggedright
Expression
\end{minipage} & \begin{minipage}[b]{\linewidth}\raggedright
Result
\end{minipage} & \begin{minipage}[b]{\linewidth}\raggedright
Expected by Programmer
\end{minipage} & \begin{minipage}[b]{\linewidth}\raggedright
Problem
\end{minipage} \\
\midrule\noalign{}
\endhead
\bottomrule\noalign{}
\endlastfoot
\texttt{"5"\ +\ 3} & \texttt{"53"} & \texttt{8} & \texttt{+}
concatenates when one operand is a string \\
\texttt{"5"\ -\ 3} & \texttt{2} & \texttt{2} & \texttt{-} converts
string to number \\
\texttt{"5"\ *\ "3"} & \texttt{15} & \texttt{15} & Multiplication
converts both \\
\texttt{""\ ==\ false} & \texttt{true} & Maybe? & Loose equality
coercion \\
\texttt{0\ ==\ false} & \texttt{true} & Maybe? & Same \\
\texttt{null\ ==\ undefined} & \texttt{true} & Probably not & \\
\texttt{NaN\ ===\ NaN} & \texttt{false} & \texttt{true} & NaN is not
equal to itself \\
\end{longtable}
}

\textbf{Mitigation:} Require \texttt{Number()} before any arithmetic on
user input. Require \texttt{===} always. Teach \texttt{typeof} and
\texttt{isNaN()} as defensive checks. Make logging intermediate values
an explicit exercise before running any calculation.

\subsubsection{\texorpdfstring{\texttt{undefined} vs.~\texttt{null}
vs.~\texttt{NaN} vs.~\texttt{false} vs.~\texttt{0}
vs.~\texttt{""}}{undefined vs.~null vs.~NaN vs.~false vs.~0 vs.~""}}\label{undefined-vs.-null-vs.-nan-vs.-false-vs.-0-vs.}

All six of these are ``falsy'' in JavaScript. Aurally, they are
indistinguishable unless you call \texttt{typeof} or check explicitly. A
student who hears \texttt{console.log(result)} and the screen reader
announces ``false'' cannot know whether \texttt{result} is
\texttt{false}, \texttt{0}, \texttt{""}, \texttt{null},
\texttt{undefined}, or \texttt{NaN} without further inspection.

\textbf{Teaching approach:} Establish a diagnostic template:

\begin{Shaded}
\begin{Highlighting}[]
\BuiltInTok{console}\OperatorTok{.}\FunctionTok{log}\NormalTok{(}\StringTok{"Value:"}\OperatorTok{,}\NormalTok{ result}\OperatorTok{,} \StringTok{"Type:"}\OperatorTok{,} \KeywordTok{typeof}\NormalTok{ result)}\OperatorTok{;}
\end{Highlighting}
\end{Shaded}

This outputs both the value and its type, providing audio context that
\texttt{console.log(result)} alone does not.

\subsubsection{Callback Hell and Asynchronous
Confusion}\label{callback-hell-and-asynchronous-confusion}

Before \texttt{async}/\texttt{await}, Node.js input used nested
callbacks:

\begin{Shaded}
\begin{Highlighting}[]
\CommentTok{// Confusing nesting {-}{-}{-} avoid in teaching}
\NormalTok{rl}\OperatorTok{.}\FunctionTok{question}\NormalTok{(}\StringTok{"Name? "}\OperatorTok{,}\NormalTok{ (name) }\KeywordTok{=\textgreater{}}\NormalTok{ \{}
\NormalTok{  rl}\OperatorTok{.}\FunctionTok{question}\NormalTok{(}\StringTok{"Age? "}\OperatorTok{,}\NormalTok{ (age) }\KeywordTok{=\textgreater{}}\NormalTok{ \{}
\NormalTok{    rl}\OperatorTok{.}\FunctionTok{question}\NormalTok{(}\StringTok{"City? "}\OperatorTok{,}\NormalTok{ (city) }\KeywordTok{=\textgreater{}}\NormalTok{ \{}
      \BuiltInTok{console}\OperatorTok{.}\FunctionTok{log}\NormalTok{(}\VerbatimStringTok{\textasciigrave{}}\SpecialCharTok{$\{}\NormalTok{name}\SpecialCharTok{\}}\VerbatimStringTok{, }\SpecialCharTok{$\{}\NormalTok{age}\SpecialCharTok{\}}\VerbatimStringTok{, }\SpecialCharTok{$\{}\NormalTok{city}\SpecialCharTok{\}}\VerbatimStringTok{\textasciigrave{}}\NormalTok{)}\OperatorTok{;}
\NormalTok{    \})}\OperatorTok{;}
\NormalTok{  \})}\OperatorTok{;}
\NormalTok{\})}\OperatorTok{;}
\end{Highlighting}
\end{Shaded}

This nesting produces code that is confusing to read AND to hear --- the
indentation deepens with each question. The
\texttt{async}/\texttt{await} pattern linearizes this:

\begin{Shaded}
\begin{Highlighting}[]
\CommentTok{// Clear and linear {-}{-}{-} use this exclusively}
\KeywordTok{const}\NormalTok{ name }\OperatorTok{=} \ControlFlowTok{await} \FunctionTok{askQuestion}\NormalTok{(}\StringTok{"Name? "}\NormalTok{)}\OperatorTok{;}
\KeywordTok{const}\NormalTok{ age }\OperatorTok{=} \ControlFlowTok{await} \FunctionTok{askQuestion}\NormalTok{(}\StringTok{"Age? "}\NormalTok{)}\OperatorTok{;}
\KeywordTok{const}\NormalTok{ city }\OperatorTok{=} \ControlFlowTok{await} \FunctionTok{askQuestion}\NormalTok{(}\StringTok{"City? "}\NormalTok{)}\OperatorTok{;}
\BuiltInTok{console}\OperatorTok{.}\FunctionTok{log}\NormalTok{(}\VerbatimStringTok{\textasciigrave{}}\SpecialCharTok{$\{}\NormalTok{name}\SpecialCharTok{\}}\VerbatimStringTok{, }\SpecialCharTok{$\{}\NormalTok{age}\SpecialCharTok{\}}\VerbatimStringTok{, }\SpecialCharTok{$\{}\NormalTok{city}\SpecialCharTok{\}}\VerbatimStringTok{\textasciigrave{}}\NormalTok{)}\OperatorTok{;}
\end{Highlighting}
\end{Shaded}

Teach \texttt{async}/\texttt{await} from the start and avoid callback
patterns entirely in this course.

\begin{center}\rule{0.5\linewidth}{0.5pt}\end{center}

\subsection{3. Classroom Setup and
Configuration}\label{classroom-setup-and-configuration}

\subsubsection{VS Code for JavaScript}\label{vs-code-for-javascript}

\textbf{Required extensions:}

\begin{itemize}
\tightlist
\item
  \textbf{ESLint} --- catches type issues, undefined variables, and
  style violations before running code
\item
  \textbf{Prettier} --- auto-formats code consistently (run on save with
  \texttt{editor.formatOnSave:\ true})
\item
  \textbf{JavaScript (ES6) code snippets} --- tab-complete common
  patterns
\end{itemize}

\textbf{Recommended ESLint configuration} (\texttt{.eslintrc.json} in
project root):

\begin{Shaded}
\begin{Highlighting}[]
\FunctionTok{\{}
  \DataTypeTok{"env"}\FunctionTok{:} \FunctionTok{\{} \DataTypeTok{"node"}\FunctionTok{:} \KeywordTok{true}\FunctionTok{,} \DataTypeTok{"es2021"}\FunctionTok{:} \KeywordTok{true} \FunctionTok{\},}
  \DataTypeTok{"extends"}\FunctionTok{:} \StringTok{"eslint:recommended"}\FunctionTok{,}
  \DataTypeTok{"rules"}\FunctionTok{:} \FunctionTok{\{}
    \DataTypeTok{"eqeqeq"}\FunctionTok{:} \StringTok{"error"}\FunctionTok{,}
    \DataTypeTok{"no{-}unused{-}vars"}\FunctionTok{:} \StringTok{"warn"}\FunctionTok{,}
    \DataTypeTok{"no{-}undef"}\FunctionTok{:} \StringTok{"error"}\FunctionTok{,}
    \DataTypeTok{"prefer{-}const"}\FunctionTok{:} \StringTok{"warn"}\FunctionTok{,}
    \DataTypeTok{"no{-}console"}\FunctionTok{:} \StringTok{"off"}
  \FunctionTok{\}}
\FunctionTok{\}}
\end{Highlighting}
\end{Shaded}

The \texttt{eqeqeq:\ error} rule makes \texttt{==} a linting error ---
it must be \texttt{===}. This is an automated enforcement of the rule
that prevents the most common silent bugs.

\textbf{Node.js REPL accessibility:}

\begin{Shaded}
\begin{Highlighting}[]
\ExtensionTok{node}                      \CommentTok{\# Start REPL}
\ExtensionTok{.help}                     \CommentTok{\# List REPL commands}
\ExtensionTok{.exit}                     \CommentTok{\# Quit (or Ctrl+D)}
\end{Highlighting}
\end{Shaded}

The Node.js REPL is excellent for early exploration. Results are printed
immediately after each expression, providing audio feedback without the
compile-run-read cycle.

\subsubsection{Alternative JavaScript
Environments}\label{alternative-javascript-environments}

{\def\LTcaptype{none} % do not increment counter
\begin{longtable}[]{@{}
  >{\raggedright\arraybackslash}p{(\linewidth - 6\tabcolsep) * \real{0.2903}}
  >{\raggedright\arraybackslash}p{(\linewidth - 6\tabcolsep) * \real{0.0860}}
  >{\raggedright\arraybackslash}p{(\linewidth - 6\tabcolsep) * \real{0.1398}}
  >{\raggedright\arraybackslash}p{(\linewidth - 6\tabcolsep) * \real{0.4839}}@{}}
\toprule\noalign{}
\begin{minipage}[b]{\linewidth}\raggedright
Environment
\end{minipage} & \begin{minipage}[b]{\linewidth}\raggedright
Platform
\end{minipage} & \begin{minipage}[b]{\linewidth}\raggedright
Accessibility
\end{minipage} & \begin{minipage}[b]{\linewidth}\raggedright
Notes
\end{minipage} \\
\midrule\noalign{}
\endhead
\bottomrule\noalign{}
\endlastfoot
\textbf{Node.js + VS Code} & All & High & Recommended setup \\
\textbf{Node.js + terminal only} & All & High & For students who prefer
minimal UI \\
\textbf{Replit.com} & Web & Medium & No-install; test with student's
setup \\
\textbf{CodeSandbox} & Web & Medium & Cloud editor; variable
accessibility \\
\textbf{GitHub Codespaces} & Web & High & Cloud VS Code; same
accessibility as local \\
\textbf{Glitch.com} & Web & Medium & Project-based; variable
accessibility \\
\textbf{Observable} & Web & Low & Notebook-style; primarily visual \\
\textbf{Deno} & Runtime & High & Alternative to Node.js; similar
accessibility \\
\end{longtable}
}

\begin{center}\rule{0.5\linewidth}{0.5pt}\end{center}

\subsection{4. Spoken JavaScript
Conventions}\label{spoken-javascript-conventions}

{\def\LTcaptype{none} % do not increment counter
\begin{longtable}[]{@{}
  >{\raggedright\arraybackslash}p{(\linewidth - 2\tabcolsep) * \real{0.3333}}
  >{\raggedright\arraybackslash}p{(\linewidth - 2\tabcolsep) * \real{0.6667}}@{}}
\toprule\noalign{}
\begin{minipage}[b]{\linewidth}\raggedright
Symbol/Pattern
\end{minipage} & \begin{minipage}[b]{\linewidth}\raggedright
Spoken Form
\end{minipage} \\
\midrule\noalign{}
\endhead
\bottomrule\noalign{}
\endlastfoot
\texttt{const} & ``const'' \\
\texttt{let} & ``let'' \\
\texttt{var} & ``var'' (avoid; legacy) \\
\texttt{===} & ``triple equals'' or ``strict equal'' \\
\texttt{!==} & ``strict not equal'' \\
\texttt{==} & ``loose equal'' (should hear this rarely) \\
\texttt{=\textgreater{}} & ``arrow'' or ``fat arrow'' \\
\texttt{...arr} & ``spread arr'' \\
\texttt{...args} & ``rest args'' \\
\texttt{async} & ``async'' \\
\texttt{await} & ``await'' \\
\texttt{undefined} & ``undefined'' \\
\texttt{null} & ``null'' \\
\texttt{NaN} & ``NaN'' (say each letter N-A-N) \\
\texttt{typeof} & ``typeof'' \\
\texttt{instanceof} & ``instanceof'' \\
\texttt{new} & ``new'' \\
\texttt{this} & ``this'' \\
\texttt{class} & ``class'' \\
\texttt{extends} & ``extends'' \\
\texttt{super} & ``super'' \\
\texttt{import} & ``import'' \\
\texttt{export} & ``export'' \\
\texttt{require()} & ``require'' \\
\texttt{module.exports} & ``module dot exports'' \\
\texttt{Promise} & ``Promise'' \\
\texttt{async\ function} & ``async function'' \\
\texttt{()\ =\textgreater{}\ \{\}} & ``arrow function'' \\
\texttt{(x)\ =\textgreater{}\ x\ *\ 2} & ``arrow function taking x,
returning x times 2'' \\
\texttt{\#field} & ``hash field'' or ``private field'' \\
\texttt{\textasciigrave{}template\ \$\{var\}\textasciigrave{}} &
``backtick template literal with variable var'' \\
\end{longtable}
}

\textbf{Reading an arrow function:}

\begin{Shaded}
\begin{Highlighting}[]
\KeywordTok{const}\NormalTok{ doubleIfPositive }\OperatorTok{=}\NormalTok{ (n) }\KeywordTok{=\textgreater{}}\NormalTok{ (n }\OperatorTok{\textgreater{}} \DecValTok{0} \OperatorTok{?}\NormalTok{ n }\OperatorTok{*} \DecValTok{2} \OperatorTok{:}\NormalTok{ n)}\OperatorTok{;}
\end{Highlighting}
\end{Shaded}

\emph{``const doubleIfPositive gets arrow function taking n, returning:
if n greater than 0, n times 2, else n''}

\begin{center}\rule{0.5\linewidth}{0.5pt}\end{center}

\subsection{5. Unit 0--1: Setup and
Foundations}\label{unit-01-setup-and-foundations}

\subsubsection{Lesson 0.1: Installing
Node.js}\label{lesson-0.1-installing-node.js}

\paragraph{Instructional Cues}\label{instructional-cues}

\textbf{Verify installation at every step:}

\begin{Shaded}
\begin{Highlighting}[]
\ExtensionTok{node} \AttributeTok{{-}{-}version}    \CommentTok{\# Should print v18.x.x or later}
\ExtensionTok{npm} \AttributeTok{{-}{-}version}     \CommentTok{\# Should print a version number}
\ExtensionTok{npx} \AttributeTok{{-}{-}version}     \CommentTok{\# Should print a version number (npx is bundled with npm)}
\end{Highlighting}
\end{Shaded}

All three must work. If \texttt{npm} is missing, the student may have
installed only the Node.js runtime without npm. Reinstall from
nodejs.org (npm is included in the standard installer).

\textbf{REPL first --- files second:} Begin with \texttt{node} in the
terminal. Let the student type expressions and hear results immediately.
This builds confidence before introducing the file -\textgreater{} run
-\textgreater{} read cycle:

\begin{verbatim}
$ node
> 5 + 3
8
> "Hello " + "World"
'Hello World'
> typeof 42
'number'
> typeof "hello"
'string'
> .exit
\end{verbatim}

\textbf{The npm init conversation:} When initializing a project,
\texttt{npm\ init\ -y} generates \texttt{package.json} automatically.
Walk the student through reading the generated file:

\begin{Shaded}
\begin{Highlighting}[]
\ExtensionTok{npm}\NormalTok{ init }\AttributeTok{{-}y}
\FunctionTok{cat}\NormalTok{ package.json   }\CommentTok{\# or type package.json on Windows}
\end{Highlighting}
\end{Shaded}

\texttt{package.json} is plain JSON text --- fully readable with a
screen reader. Understanding its structure is important for later
lessons about dependencies and scripts.

\paragraph{Non-Obvious Accessibility
Strategies}\label{non-obvious-accessibility-strategies}

\begin{itemize}
\tightlist
\item
  \textbf{Node.js version managers:} \texttt{nvm} (Node Version Manager)
  allows multiple Node.js versions. For a student who may work on
  multiple machines or projects, \texttt{nvm\ use\ 18} is a consistent
  way to set the Node.js version. \texttt{nvm} is command-line-only and
  fully accessible.
\item
  \textbf{Global vs.~local packages:} \texttt{npm\ install\ -g} installs
  globally; \texttt{npm\ install} installs locally (in
  \texttt{node\_modules}). Teach local installation as the default.
  Global packages can interfere across projects.
\end{itemize}

\paragraph{Thought Experiments}\label{thought-experiments}

\begin{enumerate}
\def\labelenumi{\arabic{enumi}.}
\item
  \emph{``JavaScript was created in 10 days in 1995. What marks of that
  hasty design are still visible in modern JavaScript? (Think about
  \texttt{typeof\ null\ ===\ \textquotesingle{}object\textquotesingle{}},
  \texttt{NaN\ !==\ NaN}, \texttt{0\ ==\ false}.) Why haven't these been
  fixed?''}
\item
  \emph{``Node.js uses a single-threaded event loop to handle many
  concurrent operations. A traditional multi-threaded approach uses
  multiple threads running in parallel. For a web server handling
  thousands of requests, what are the tradeoffs of each approach? Which
  tends to be more memory efficient?''}
\item
  \emph{``npm has millions of packages. The average Node.js project
  depends on hundreds of packages, many maintained by single developers
  with no organizational backing. The `left-pad incident' in 2016 broke
  thousands of projects when one 11-line package was unpublished. What
  does this incident reveal about the risks and benefits of the npm
  ecosystem?''}
\item
  \emph{``JavaScript runs in browsers, in Node.js servers, in mobile
  apps via React Native, in desktop apps via Electron, in IoT devices,
  and even in serverless cloud functions. What does this ubiquity tell
  you about the relationship between a language's design and its
  adoption?''}
\end{enumerate}

\begin{center}\rule{0.5\linewidth}{0.5pt}\end{center}

\subsubsection{Lesson 1.3: Variables and Data
Types}\label{lesson-1.3-variables-and-data-types}

\paragraph{Instructional Cues}\label{instructional-cues-1}

\textbf{\texttt{const} vs.~\texttt{let} --- establish the rule on day
one:}

\emph{``Use \texttt{const} for everything. If you get an error saying
you're trying to reassign a \texttt{const}, THEN change it to
\texttt{let}. This approach produces cleaner code and prevents
accidental reassignment bugs.''}

Demonstrate the error:

\begin{Shaded}
\begin{Highlighting}[]
\KeywordTok{const}\NormalTok{ x }\OperatorTok{=} \DecValTok{5}\OperatorTok{;}
\NormalTok{x }\OperatorTok{=} \DecValTok{10}\OperatorTok{;} \CommentTok{// TypeError: Assignment to constant variable.}
\end{Highlighting}
\end{Shaded}

\emph{``This error is helpful. \texttt{const} is protecting you from
changing something you said was fixed.''}

\textbf{Dynamic typing with auditory ambiguity:}

Run this sequence in the REPL and ask the student to predict each output
before pressing Enter:

\begin{Shaded}
\begin{Highlighting}[]
\KeywordTok{typeof} \DecValTok{42}\OperatorTok{;} \CommentTok{// "number"}
\KeywordTok{typeof} \StringTok{"42"}\OperatorTok{;} \CommentTok{// "string"}
\KeywordTok{typeof} \KeywordTok{true}\OperatorTok{;} \CommentTok{// "boolean"}
\KeywordTok{typeof} \KeywordTok{null}\OperatorTok{;} \CommentTok{// "object"  \textless{}{-} the famous bug}
\KeywordTok{typeof} \KeywordTok{undefined}\OperatorTok{;} \CommentTok{// "undefined"}
\KeywordTok{typeof}\NormalTok{ \{\}}\OperatorTok{;} \CommentTok{// "object"}
\KeywordTok{typeof}\NormalTok{ []}\OperatorTok{;} \CommentTok{// "object"  \textless{}{-} arrays are objects!}
\KeywordTok{typeof} \KeywordTok{function}\NormalTok{ () \{\}}\OperatorTok{;} \CommentTok{// "function"}
\end{Highlighting}
\end{Shaded}

\emph{``Notice that both \texttt{\{\}} (object) and \texttt{{[}{]}}
(array) produce `object'. How do you tell them apart?
\texttt{Array.isArray({[}{]})} returns \texttt{true}.
\texttt{Array.isArray(\{\})} returns \texttt{false}.''}

\textbf{Template literals vs.~concatenation:}

\begin{Shaded}
\begin{Highlighting}[]
\KeywordTok{const}\NormalTok{ name }\OperatorTok{=} \StringTok{"Alice"}\OperatorTok{;}
\KeywordTok{const}\NormalTok{ age }\OperatorTok{=} \DecValTok{15}\OperatorTok{;}

\CommentTok{// Concatenation {-}{-}{-} confusing order of operations when read aloud}
\BuiltInTok{console}\OperatorTok{.}\FunctionTok{log}\NormalTok{(}\StringTok{"Hello, "} \OperatorTok{+}\NormalTok{ name }\OperatorTok{+} \StringTok{"! You are "} \OperatorTok{+}\NormalTok{ age }\OperatorTok{+} \StringTok{" years old."}\NormalTok{)}\OperatorTok{;}

\CommentTok{// Template literal {-}{-}{-} linear, clean, easy to hear}
\BuiltInTok{console}\OperatorTok{.}\FunctionTok{log}\NormalTok{(}\VerbatimStringTok{\textasciigrave{}Hello, }\SpecialCharTok{$\{}\NormalTok{name}\SpecialCharTok{\}}\VerbatimStringTok{! You are }\SpecialCharTok{$\{}\NormalTok{age}\SpecialCharTok{\}}\VerbatimStringTok{ years old.\textasciigrave{}}\NormalTok{)}\OperatorTok{;}
\end{Highlighting}
\end{Shaded}

Read both versions aloud with the student. Ask which is easier to follow
aurally. Template literals are almost universally clearer when heard,
because the variable values appear where they belong in the sentence.

\paragraph{Thought Experiments}\label{thought-experiments-1}

\begin{enumerate}
\def\labelenumi{\arabic{enumi}.}
\item
  \emph{``JavaScript has only one number type for integers and decimals.
  Python has separate \texttt{int} and \texttt{float}. Java has
  \texttt{int}, \texttt{long}, \texttt{double}, \texttt{float}. What
  precision problems arise from having a single number type? When does
  \texttt{0.1\ +\ 0.2\ =\ 0.30000000000000004} matter, and when doesn't
  it?''}
\item
  \emph{``JavaScript has six falsy values: \texttt{false}, \texttt{0},
  \texttt{""}, \texttt{null}, \texttt{undefined}, \texttt{NaN}.
  Everything else is truthy --- including \texttt{"false"} (a non-empty
  string), \texttt{{[}{]}} (an empty array), and \texttt{\{\}} (an empty
  object). Design a quiz: ask the student whether ten different
  expressions are truthy or falsy. Which answers are surprising?''}
\item
  \emph{``\texttt{var} has function scope; \texttt{let} and
  \texttt{const} have block scope. What does `scope' mean? How does a
  variable declared inside an \texttt{if} block with \texttt{var}
  vs.~\texttt{let} behave differently? Which behavior is more
  predictable and why?''}
\end{enumerate}

\begin{center}\rule{0.5\linewidth}{0.5pt}\end{center}

\subsection{6. Unit 2: Numbers, Input, and
Math}\label{unit-2-numbers-input-and-math}

\subsubsection{Lesson 2.2: User Input in
Node.js}\label{lesson-2.2-user-input-in-node.js}

\paragraph{Instructional Cues}\label{instructional-cues-2}

\textbf{The input helper module pattern:}

Provide the \texttt{input.js} module in every project skeleton. Students
should not need to re-create this boilerplate:

\begin{Shaded}
\begin{Highlighting}[]
\CommentTok{// input.js}
\KeywordTok{const}\NormalTok{ readline }\OperatorTok{=} \PreprocessorTok{require}\NormalTok{(}\StringTok{"readline"}\NormalTok{)}\OperatorTok{;}
\KeywordTok{const}\NormalTok{ rl }\OperatorTok{=}\NormalTok{ readline}\OperatorTok{.}\FunctionTok{createInterface}\NormalTok{(\{}
  \DataTypeTok{input}\OperatorTok{:} \BuiltInTok{process}\OperatorTok{.}\AttributeTok{stdin}\OperatorTok{,}
  \DataTypeTok{output}\OperatorTok{:} \BuiltInTok{process}\OperatorTok{.}\AttributeTok{stdout}\OperatorTok{,}
\NormalTok{\})}\OperatorTok{;}

\KeywordTok{function} \FunctionTok{askQuestion}\NormalTok{(prompt) \{}
  \ControlFlowTok{return} \KeywordTok{new} \BuiltInTok{Promise}\NormalTok{((resolve) }\KeywordTok{=\textgreater{}}\NormalTok{ rl}\OperatorTok{.}\FunctionTok{question}\NormalTok{(prompt}\OperatorTok{,}\NormalTok{ resolve))}\OperatorTok{;}
\NormalTok{\}}

\KeywordTok{function} \FunctionTok{close}\NormalTok{() \{}
\NormalTok{  rl}\OperatorTok{.}\FunctionTok{close}\NormalTok{()}\OperatorTok{;}
\NormalTok{\}}

\NormalTok{module}\OperatorTok{.}\AttributeTok{exports} \OperatorTok{=}\NormalTok{ \{ askQuestion}\OperatorTok{,}\NormalTok{ close \}}\OperatorTok{;}
\end{Highlighting}
\end{Shaded}

Walk through what each line does at least once so students understand
the structure, but do not require students to write it from scratch
until they're comfortable with async programming.

\textbf{The \texttt{async}/\texttt{await} mental model --- teach this
story:}

\emph{``When your code hits \texttt{await}, imagine it puts up a `please
wait' sign and steps aside. Other things can happen while you wait. When
the answer arrives, your code picks up exactly where it paused. The
\texttt{async} keyword on the function declaration is your promise that
`this function may pause at some point.'\,''}

\begin{Shaded}
\begin{Highlighting}[]
\KeywordTok{async} \KeywordTok{function} \FunctionTok{main}\NormalTok{() \{}
  \BuiltInTok{console}\OperatorTok{.}\FunctionTok{log}\NormalTok{(}\StringTok{"Starting..."}\NormalTok{)}\OperatorTok{;} \CommentTok{// 1: runs immediately}
  \KeywordTok{const}\NormalTok{ name }\OperatorTok{=} \ControlFlowTok{await} \FunctionTok{askQuestion}\NormalTok{(}\StringTok{"Name? "}\NormalTok{)}\OperatorTok{;} \CommentTok{// 2: pauses here until user types}
  \BuiltInTok{console}\OperatorTok{.}\FunctionTok{log}\NormalTok{(}\StringTok{"Got name:"}\OperatorTok{,}\NormalTok{ name)}\OperatorTok{;} \CommentTok{// 3: runs after user answers}
  \KeywordTok{const}\NormalTok{ age }\OperatorTok{=} \ControlFlowTok{await} \FunctionTok{askQuestion}\NormalTok{(}\StringTok{"Age? "}\NormalTok{)}\OperatorTok{;} \CommentTok{// 4: pauses again}
  \BuiltInTok{console}\OperatorTok{.}\FunctionTok{log}\NormalTok{(}\VerbatimStringTok{\textasciigrave{}}\SpecialCharTok{$\{}\NormalTok{name}\SpecialCharTok{\}}\VerbatimStringTok{ is }\SpecialCharTok{$\{}\NormalTok{age}\SpecialCharTok{\}}\VerbatimStringTok{\textasciigrave{}}\NormalTok{)}\OperatorTok{;} \CommentTok{// 5: runs after user answers again}
  \FunctionTok{close}\NormalTok{()}\OperatorTok{;}
\NormalTok{\}}

\FunctionTok{main}\NormalTok{()}\OperatorTok{;} \CommentTok{// Kick off the async function}
\end{Highlighting}
\end{Shaded}

Ask: ``What line is running while the program waits for user input?''
(None --- the function is paused. The event loop can handle other
things.)

\textbf{Why \texttt{main()} must be called at the bottom:}
\texttt{async\ function\ main()} defines the function but doesn't run
it. The \texttt{main()} call at the end actually starts execution.
Without the call, nothing happens. This trips up beginners ---
demonstrate it.

\paragraph{Non-Obvious Accessibility
Strategies}\label{non-obvious-accessibility-strategies-1}

\begin{itemize}
\item
  \textbf{\texttt{rl.close()} is mandatory:} Without
  \texttt{rl.close()}, the Node.js process continues running after
  \texttt{main()} finishes. The terminal appears to hang --- cursor
  blinks but no prompt. For a screen reader user, this is confusing.
  Teach \texttt{close()} as the final step of every program that uses
  input.
\item
  \textbf{Prompt trailing space:} Always include a trailing space in
  prompt strings: \texttt{"Enter\ your\ name:\ "}. Without the space,
  the cursor appears immediately after the colon with no separation,
  which some screen readers render as ``Enter your name:cursor'' --- an
  awkward audio experience.
\item
  \textbf{Input validation loop:} Without validation, a blank Enter or a
  non-numeric string for a number causes silent failure. Provide this
  reusable pattern:
\end{itemize}

\begin{Shaded}
\begin{Highlighting}[]
\KeywordTok{async} \KeywordTok{function} \FunctionTok{getNumber}\NormalTok{(prompt) \{}
  \ControlFlowTok{while}\NormalTok{ (}\KeywordTok{true}\NormalTok{) \{}
    \KeywordTok{const}\NormalTok{ input }\OperatorTok{=} \ControlFlowTok{await} \FunctionTok{askQuestion}\NormalTok{(prompt)}\OperatorTok{;}
    \KeywordTok{const}\NormalTok{ num }\OperatorTok{=} \BuiltInTok{Number}\NormalTok{(input}\OperatorTok{.}\FunctionTok{trim}\NormalTok{())}\OperatorTok{;}
    \ControlFlowTok{if}\NormalTok{ (}\OperatorTok{!}\PreprocessorTok{isNaN}\NormalTok{(num)) }\ControlFlowTok{return}\NormalTok{ num}\OperatorTok{;}
    \BuiltInTok{console}\OperatorTok{.}\FunctionTok{log}\NormalTok{(}\StringTok{"Invalid input: please enter a number."}\NormalTok{)}\OperatorTok{;}
\NormalTok{  \}}
\NormalTok{\}}
\end{Highlighting}
\end{Shaded}

Teach this pattern and require students to use it whenever they need
numeric input.

\paragraph{Thought Experiments}\label{thought-experiments-2}

\begin{enumerate}
\def\labelenumi{\arabic{enumi}.}
\item
  \emph{``The \texttt{readline} module uses an `event listener' pattern
  --- \texttt{rl.question(prompt,\ callback)} means `when the user
  answers, call this function.' The \texttt{async}/\texttt{await}
  pattern wraps this in a Promise. What problem does the Promise solve
  that the raw callback doesn't? (Think about reading the code
  linearly.)''}
\item
  \emph{``Node.js uses a single-threaded event loop. While your program
  is waiting for user input (\texttt{await\ askQuestion}), the thread is
  free. What else could the thread be doing? In a real application, this
  might mean handling a different request or updating state. What does
  this tell you about why \texttt{async}/\texttt{await} is important for
  scalable programs?''}
\item
  \emph{``Synchronous input (blocking until the user responds) and
  asynchronous input (doing other things while waiting) represent
  fundamentally different programming models. For a single-user
  command-line program, which matters? For a server handling 10,000
  simultaneous users, which matters more? Why?''}
\end{enumerate}

\begin{center}\rule{0.5\linewidth}{0.5pt}\end{center}

\subsection{7. Unit 3: Control Flow}\label{unit-3-control-flow}

\subsubsection{Lessons 3.1--3.3}\label{lessons-3.13.3}

\paragraph{Instructional Cues}\label{instructional-cues-3}

\textbf{\texttt{===} enforcement --- make it automatic:}

With ESLint configured to flag \texttt{==}, the student will hear a
warning whenever they use loose equality. Walk through the first warning
together:

\begin{Shaded}
\begin{Highlighting}[]
\ControlFlowTok{if}\NormalTok{ (userInput }\OperatorTok{==} \DecValTok{5}\NormalTok{) \{   }\CommentTok{// ESLint: Expected \textquotesingle{}===\textquotesingle{} and instead saw \textquotesingle{}==\textquotesingle{}}
\end{Highlighting}
\end{Shaded}

\emph{``ESLint is protecting you from a common silent bug. Change this
to \texttt{===}. The only exception is checking for null-or-undefined:
\texttt{x\ ==\ null} is sometimes used as a shorthand for
\texttt{x\ ===\ null\ \textbar{}\textbar{}\ x\ ===\ undefined} --- but
avoid this until you understand it.''}

\textbf{Truthy/falsy in input validation:}

\begin{Shaded}
\begin{Highlighting}[]
\KeywordTok{const}\NormalTok{ name }\OperatorTok{=} \ControlFlowTok{await} \FunctionTok{askQuestion}\NormalTok{(}\StringTok{"Enter your name: "}\NormalTok{)}\OperatorTok{;}
\ControlFlowTok{if}\NormalTok{ (}\OperatorTok{!}\NormalTok{name}\OperatorTok{.}\FunctionTok{trim}\NormalTok{()) \{}
  \BuiltInTok{console}\OperatorTok{.}\FunctionTok{log}\NormalTok{(}\StringTok{"Name cannot be empty."}\NormalTok{)}\OperatorTok{;}
  \CommentTok{// re{-}prompt...}
\NormalTok{\}}
\end{Highlighting}
\end{Shaded}

\texttt{!name.trim()} is true when the trimmed string is empty
(\texttt{""} is falsy). This is idiomatic JavaScript. Read it aloud:
``if not name dot trim'' --- ask: ``What does `not an empty string'
mean? Why is an empty string falsy?''

\textbf{Ternary operator --- limit its use:}

The ternary \texttt{condition\ ?\ valueIfTrue\ :\ valueIfFalse} is
concise but can be hard to follow when read aloud:

\begin{Shaded}
\begin{Highlighting}[]
\KeywordTok{const}\NormalTok{ status }\OperatorTok{=}\NormalTok{ score }\OperatorTok{\textgreater{}=} \DecValTok{60} \OperatorTok{?} \StringTok{"pass"} \OperatorTok{:} \StringTok{"fail"}\OperatorTok{;}
\end{Highlighting}
\end{Shaded}

Read: \emph{``const status gets: if score greater-than-or-equal 60,
quote pass, else quote fail''}

Teach it but establish this rule: ``If the ternary is easy to read aloud
in one breath, use it. If it requires re-reading, use if-else.''

\paragraph{Thought Experiments}\label{thought-experiments-3}

\begin{enumerate}
\def\labelenumi{\arabic{enumi}.}
\item
  \emph{``JavaScript's \texttt{==} performs type coercion before
  comparing. TypeScript (a typed superset of JavaScript) was partly
  created to address this. What does TypeScript's existence tell you
  about the real-world cost of JavaScript's permissiveness? What did
  developers gain from TypeScript that JavaScript didn't provide?''}
\item
  \emph{``Consider: \texttt{if\ (userName)}
  vs.~\texttt{if\ (userName\ !==\ \textquotesingle{}\textquotesingle{}\ \&\&\ userName\ !==\ null\ \&\&\ userName\ !==\ undefined)}.
  The first uses JavaScript's truthiness; the second is explicit. In a
  code review, which would you prefer seeing? What are the readability
  and correctness tradeoffs?''}
\item
  \emph{``JavaScript's optional chaining (\texttt{?.}) and nullish
  coalescing (\texttt{??}) were added to the language to handle
  null/undefined more gracefully. \texttt{user?.address?.city} returns
  \texttt{undefined} instead of throwing if \texttt{user} or
  \texttt{address} is null. Why did this feature take decades to add to
  a language created in 1995?''}
\item
  \emph{``The \texttt{switch} statement falls through by default in
  JavaScript. Modern JavaScript has been moving away from
  \texttt{switch} toward mapping objects for dispatch tables.
  \texttt{const\ actions\ =\ \{\ \textquotesingle{}add\textquotesingle{}:\ addFunc,\ \textquotesingle{}delete\textquotesingle{}:\ deleteFunc\ \};\ actions{[}command{]}?.();}.
  What are the accessibility implications of this pattern vs.~a
  \texttt{switch} statement when read with a screen reader?''}
\end{enumerate}

\begin{center}\rule{0.5\linewidth}{0.5pt}\end{center}

\subsection{8. Unit 4: Arrays and
Objects}\label{unit-4-arrays-and-objects}

\subsubsection{Lesson 4.1: Arrays and Functional
Methods}\label{lesson-4.1-arrays-and-functional-methods}

\paragraph{Instructional Cues}\label{instructional-cues-4}

\textbf{Introduce \texttt{map}, \texttt{filter}, \texttt{reduce} as a
trio:}

Don't teach them one per lesson --- teach the relationship between them
first:

\emph{``These three methods each take a function and apply it to every
element in an array. \texttt{map} transforms elements. \texttt{filter}
keeps or removes elements. \texttt{reduce} combines elements into one
value. Together they cover most array processing.''}

Then give one shared example that uses all three:

\begin{Shaded}
\begin{Highlighting}[]
\KeywordTok{const}\NormalTok{ students }\OperatorTok{=}\NormalTok{ [}
\NormalTok{  \{ }\DataTypeTok{name}\OperatorTok{:} \StringTok{"Alice"}\OperatorTok{,} \DataTypeTok{score}\OperatorTok{:} \DecValTok{88}\NormalTok{ \}}\OperatorTok{,}
\NormalTok{  \{ }\DataTypeTok{name}\OperatorTok{:} \StringTok{"Bob"}\OperatorTok{,} \DataTypeTok{score}\OperatorTok{:} \DecValTok{62}\NormalTok{ \}}\OperatorTok{,}
\NormalTok{  \{ }\DataTypeTok{name}\OperatorTok{:} \StringTok{"Charlie"}\OperatorTok{,} \DataTypeTok{score}\OperatorTok{:} \DecValTok{95}\NormalTok{ \}}\OperatorTok{,}
\NormalTok{  \{ }\DataTypeTok{name}\OperatorTok{:} \StringTok{"Diana"}\OperatorTok{,} \DataTypeTok{score}\OperatorTok{:} \DecValTok{71}\NormalTok{ \}}\OperatorTok{,}
\NormalTok{]}\OperatorTok{;}

\CommentTok{// 1. Filter passing students (score \textgreater{}= 70)}
\KeywordTok{const}\NormalTok{ passing }\OperatorTok{=}\NormalTok{ students}\OperatorTok{.}\FunctionTok{filter}\NormalTok{((s) }\KeywordTok{=\textgreater{}}\NormalTok{ s}\OperatorTok{.}\AttributeTok{score} \OperatorTok{\textgreater{}=} \DecValTok{70}\NormalTok{)}\OperatorTok{;}

\CommentTok{// 2. Map to get just names}
\KeywordTok{const}\NormalTok{ passingNames }\OperatorTok{=}\NormalTok{ passing}\OperatorTok{.}\FunctionTok{map}\NormalTok{((s) }\KeywordTok{=\textgreater{}}\NormalTok{ s}\OperatorTok{.}\AttributeTok{name}\NormalTok{)}\OperatorTok{;}

\CommentTok{// 3. Reduce to build a report string}
\KeywordTok{const}\NormalTok{ report }\OperatorTok{=}\NormalTok{ passingNames}\OperatorTok{.}\FunctionTok{reduce}\NormalTok{(}
\NormalTok{  (acc}\OperatorTok{,}\NormalTok{ name) }\KeywordTok{=\textgreater{}}\NormalTok{ acc }\OperatorTok{+}\NormalTok{ name }\OperatorTok{+} \StringTok{", "}\OperatorTok{,}
  \StringTok{"Passing: "}\OperatorTok{,}
\NormalTok{)}\OperatorTok{;}

\BuiltInTok{console}\OperatorTok{.}\FunctionTok{log}\NormalTok{(report)}\OperatorTok{;} \CommentTok{// "Passing: Alice, Charlie, Diana, "}
\BuiltInTok{console}\OperatorTok{.}\FunctionTok{log}\NormalTok{(}\StringTok{"Count:"}\OperatorTok{,}\NormalTok{ passing}\OperatorTok{.}\AttributeTok{length}\NormalTok{)}\OperatorTok{;}
\end{Highlighting}
\end{Shaded}

\textbf{Accessible printing for arrays of objects:}

\begin{Shaded}
\begin{Highlighting}[]
\CommentTok{// Accessible: numbered, labeled fields}
\NormalTok{students}\OperatorTok{.}\FunctionTok{forEach}\NormalTok{((student}\OperatorTok{,}\NormalTok{ index) }\KeywordTok{=\textgreater{}}\NormalTok{ \{}
  \BuiltInTok{console}\OperatorTok{.}\FunctionTok{log}\NormalTok{(}\VerbatimStringTok{\textasciigrave{}}\SpecialCharTok{$\{}\NormalTok{index }\OperatorTok{+} \DecValTok{1}\SpecialCharTok{\}}\VerbatimStringTok{. }\SpecialCharTok{$\{}\NormalTok{student}\OperatorTok{.}\AttributeTok{name}\SpecialCharTok{\}}\VerbatimStringTok{: }\SpecialCharTok{$\{}\NormalTok{student}\OperatorTok{.}\AttributeTok{score}\SpecialCharTok{\}}\VerbatimStringTok{\textasciigrave{}}\NormalTok{)}\OperatorTok{;}
\NormalTok{\})}\OperatorTok{;}

\CommentTok{// Less accessible: raw object output}
\BuiltInTok{console}\OperatorTok{.}\FunctionTok{log}\NormalTok{(students)}\OperatorTok{;} \CommentTok{// Produces cluttered bracket notation when read aloud}
\end{Highlighting}
\end{Shaded}

\textbf{Chaining --- build up incrementally:}

Show chaining as a composition of steps, each of which you can trace
separately:

\begin{Shaded}
\begin{Highlighting}[]
\KeywordTok{const}\NormalTok{ result }\OperatorTok{=}\NormalTok{ numbers}
  \OperatorTok{.}\FunctionTok{filter}\NormalTok{((n) }\KeywordTok{=\textgreater{}}\NormalTok{ n }\OperatorTok{\%} \DecValTok{2} \OperatorTok{===} \DecValTok{0}\NormalTok{) }\CommentTok{// Step 1: keep evens}
  \OperatorTok{.}\FunctionTok{map}\NormalTok{((n) }\KeywordTok{=\textgreater{}}\NormalTok{ n }\OperatorTok{**} \DecValTok{2}\NormalTok{) }\CommentTok{// Step 2: square each}
  \OperatorTok{.}\FunctionTok{reduce}\NormalTok{((sum}\OperatorTok{,}\NormalTok{ n) }\KeywordTok{=\textgreater{}}\NormalTok{ sum }\OperatorTok{+}\NormalTok{ n}\OperatorTok{,} \DecValTok{0}\NormalTok{)}\OperatorTok{;} \CommentTok{// Step 3: sum}

\CommentTok{// DEBUG: break the chain to trace each step}
\KeywordTok{const}\NormalTok{ step1 }\OperatorTok{=}\NormalTok{ numbers}\OperatorTok{.}\FunctionTok{filter}\NormalTok{((n) }\KeywordTok{=\textgreater{}}\NormalTok{ n }\OperatorTok{\%} \DecValTok{2} \OperatorTok{===} \DecValTok{0}\NormalTok{)}\OperatorTok{;}
\BuiltInTok{console}\OperatorTok{.}\FunctionTok{log}\NormalTok{(}\StringTok{"After filter:"}\OperatorTok{,}\NormalTok{ step1)}\OperatorTok{;}
\KeywordTok{const}\NormalTok{ step2 }\OperatorTok{=}\NormalTok{ step1}\OperatorTok{.}\FunctionTok{map}\NormalTok{((n) }\KeywordTok{=\textgreater{}}\NormalTok{ n }\OperatorTok{**} \DecValTok{2}\NormalTok{)}\OperatorTok{;}
\BuiltInTok{console}\OperatorTok{.}\FunctionTok{log}\NormalTok{(}\StringTok{"After map:"}\OperatorTok{,}\NormalTok{ step2)}\OperatorTok{;}
\KeywordTok{const}\NormalTok{ step3 }\OperatorTok{=}\NormalTok{ step2}\OperatorTok{.}\FunctionTok{reduce}\NormalTok{((sum}\OperatorTok{,}\NormalTok{ n) }\KeywordTok{=\textgreater{}}\NormalTok{ sum }\OperatorTok{+}\NormalTok{ n}\OperatorTok{,} \DecValTok{0}\NormalTok{)}\OperatorTok{;}
\BuiltInTok{console}\OperatorTok{.}\FunctionTok{log}\NormalTok{(}\StringTok{"Final sum:"}\OperatorTok{,}\NormalTok{ step3)}\OperatorTok{;}
\end{Highlighting}
\end{Shaded}

The debug version (step-by-step) is more accessible during development.
Consolidate to the chain after confirming correctness.

\paragraph{Thought Experiments}\label{thought-experiments-4}

\begin{enumerate}
\def\labelenumi{\arabic{enumi}.}
\item
  \emph{``\texttt{map}, \texttt{filter}, and \texttt{reduce} come from
  functional programming --- a paradigm based on mathematical functions
  that have no side effects. JavaScript arrays also have
  \texttt{forEach} which is for side effects only. What is the
  philosophical distinction between \texttt{map} and \texttt{forEach}?
  Why does this distinction matter for reasoning about code
  correctness?''}
\item
  \emph{``The \texttt{sort} method in JavaScript sorts in-place and
  mutates the original array. \texttt{filter}, \texttt{map}, and
  \texttt{reduce} return new arrays, leaving the original unchanged. Why
  the asymmetry? What bugs could arise from \texttt{sort}'s mutation?
  How would you sort without mutating:
  \texttt{{[}...arr{]}.sort(...)}?''}
\item
  \_``Functional composition:
  \texttt{const\ double\ =\ n\ =\textgreater{}\ n\ \_\ 2;\ const\ addOne\ =\ n\ =\textgreater{}\ n\ +\ 1;\ const\ doubleAndAdd\ =\ n\ =\textgreater{}\ addOne(double(n))}.
  How does function composition scale? What is the advantage of
  composing small, focused functions over writing one large function
  that does everything?''*
\item
  \emph{``The spread operator \texttt{{[}...arr{]}} creates a shallow
  copy of an array. For an array of primitives \texttt{{[}1,\ 2,\ 3{]}},
  the spread copy is fully independent. For an array of objects
  \texttt{{[}\{name:\ \textquotesingle{}A\textquotesingle{}\},\ \{name:\ \textquotesingle{}B\textquotesingle{}\}{]}},
  the objects inside are still shared. What does `shallow copy' mean?
  When does this distinction cause bugs?''}
\end{enumerate}

\begin{center}\rule{0.5\linewidth}{0.5pt}\end{center}

\subsubsection{Lesson 4.2: Objects}\label{lesson-4.2-objects}

\paragraph{Instructional Cues}\label{instructional-cues-5}

\textbf{Destructuring --- introduce before objects are complex:}

\begin{Shaded}
\begin{Highlighting}[]
\KeywordTok{const}\NormalTok{ person }\OperatorTok{=}\NormalTok{ \{ }\DataTypeTok{name}\OperatorTok{:} \StringTok{"Alice"}\OperatorTok{,} \DataTypeTok{age}\OperatorTok{:} \DecValTok{15}\OperatorTok{,} \DataTypeTok{city}\OperatorTok{:} \StringTok{"Portland"}\NormalTok{ \}}\OperatorTok{;}

\CommentTok{// Without destructuring:}
\KeywordTok{const}\NormalTok{ name }\OperatorTok{=}\NormalTok{ person}\OperatorTok{.}\AttributeTok{name}\OperatorTok{;}
\KeywordTok{const}\NormalTok{ age }\OperatorTok{=}\NormalTok{ person}\OperatorTok{.}\AttributeTok{age}\OperatorTok{;}

\CommentTok{// With destructuring:}
\KeywordTok{const}\NormalTok{ \{ name}\OperatorTok{,}\NormalTok{ age \} }\OperatorTok{=}\NormalTok{ person}\OperatorTok{;}

\CommentTok{// In function parameter:}
\KeywordTok{function} \FunctionTok{introduce}\NormalTok{(\{ name}\OperatorTok{,}\NormalTok{ age \}) \{}
  \BuiltInTok{console}\OperatorTok{.}\FunctionTok{log}\NormalTok{(}\VerbatimStringTok{\textasciigrave{}Hi, I\textquotesingle{}m }\SpecialCharTok{$\{}\NormalTok{name}\SpecialCharTok{\}}\VerbatimStringTok{, }\SpecialCharTok{$\{}\NormalTok{age}\SpecialCharTok{\}}\VerbatimStringTok{ years old.\textasciigrave{}}\NormalTok{)}\OperatorTok{;}
\NormalTok{\}}
\FunctionTok{introduce}\NormalTok{(person)}\OperatorTok{;}
\end{Highlighting}
\end{Shaded}

Destructuring is ubiquitous in modern JavaScript. Teaching it early
reduces the verbosity of object access code, which in turn makes
programs cleaner to read and hear.

\textbf{\texttt{Object.keys()}, \texttt{Object.values()},
\texttt{Object.entries()} as navigation tools:}

For a student who cannot visually inspect an object in a debugger panel,
these methods provide audio-accessible object introspection:

\begin{Shaded}
\begin{Highlighting}[]
\KeywordTok{const}\NormalTok{ config }\OperatorTok{=}\NormalTok{ \{ }\DataTypeTok{theme}\OperatorTok{:} \StringTok{"dark"}\OperatorTok{,} \DataTypeTok{fontSize}\OperatorTok{:} \DecValTok{14}\OperatorTok{,} \DataTypeTok{notifications}\OperatorTok{:} \KeywordTok{true}\NormalTok{ \}}\OperatorTok{;}

\CommentTok{// What properties does this object have?}
\BuiltInTok{console}\OperatorTok{.}\FunctionTok{log}\NormalTok{(}\BuiltInTok{Object}\OperatorTok{.}\FunctionTok{keys}\NormalTok{(config)}\OperatorTok{.}\FunctionTok{join}\NormalTok{(}\StringTok{", "}\NormalTok{))}\OperatorTok{;}

\CommentTok{// What are all the values?}
\BuiltInTok{Object}\OperatorTok{.}\FunctionTok{values}\NormalTok{(config)}\OperatorTok{.}\FunctionTok{forEach}\NormalTok{((v) }\KeywordTok{=\textgreater{}} \BuiltInTok{console}\OperatorTok{.}\FunctionTok{log}\NormalTok{(v))}\OperatorTok{;}

\CommentTok{// Both together:}
\BuiltInTok{Object}\OperatorTok{.}\FunctionTok{entries}\NormalTok{(config)}\OperatorTok{.}\FunctionTok{forEach}\NormalTok{(([key}\OperatorTok{,}\NormalTok{ value]) }\KeywordTok{=\textgreater{}}\NormalTok{ \{}
  \BuiltInTok{console}\OperatorTok{.}\FunctionTok{log}\NormalTok{(}\VerbatimStringTok{\textasciigrave{}}\SpecialCharTok{$\{}\NormalTok{key}\SpecialCharTok{\}}\VerbatimStringTok{: }\SpecialCharTok{$\{}\NormalTok{value}\SpecialCharTok{\}}\VerbatimStringTok{\textasciigrave{}}\NormalTok{)}\OperatorTok{;}
\NormalTok{\})}\OperatorTok{;}
\end{Highlighting}
\end{Shaded}

\textbf{Optional chaining for safe nested access:}

\begin{Shaded}
\begin{Highlighting}[]
\KeywordTok{const}\NormalTok{ user }\OperatorTok{=}\NormalTok{ \{ }\DataTypeTok{profile}\OperatorTok{:}\NormalTok{ \{ }\DataTypeTok{address}\OperatorTok{:} \KeywordTok{null}\NormalTok{ \} \}}\OperatorTok{;}

\CommentTok{// Without optional chaining {-}{-}{-} crashes if address is null}
\KeywordTok{const}\NormalTok{ city }\OperatorTok{=}\NormalTok{ user}\OperatorTok{.}\AttributeTok{profile}\OperatorTok{.}\AttributeTok{address}\OperatorTok{.}\AttributeTok{city}\OperatorTok{;} \CommentTok{// TypeError!}

\CommentTok{// With optional chaining {-}{-}{-} returns undefined instead}
\KeywordTok{const}\NormalTok{ city }\OperatorTok{=}\NormalTok{ user}\OperatorTok{?.}\AttributeTok{profile}\OperatorTok{?.}\AttributeTok{address}\OperatorTok{?.}\AttributeTok{city}\OperatorTok{;} \CommentTok{// undefined, no crash}

\CommentTok{// With nullish coalescing for a default}
\KeywordTok{const}\NormalTok{ city }\OperatorTok{=}\NormalTok{ user}\OperatorTok{?.}\AttributeTok{profile}\OperatorTok{?.}\AttributeTok{address}\OperatorTok{?.}\AttributeTok{city} \OperatorTok{??} \StringTok{"Unknown city"}\OperatorTok{;}
\end{Highlighting}
\end{Shaded}

This pattern is essential for working with real-world data where fields
may be missing.

\paragraph{Thought Experiments}\label{thought-experiments-5}

\begin{enumerate}
\def\labelenumi{\arabic{enumi}.}
\item
  \emph{``JavaScript objects can have properties added or removed at any
  time.
  \texttt{user.newField\ =\ \textquotesingle{}value\textquotesingle{}}
  just works. Java objects have a fixed set of fields defined at compile
  time. Which approach is better for representing data from an API that
  might add new fields in future versions? Which is better for ensuring
  all required fields are present?''}
\item
  \emph{``\texttt{Object.freeze(obj)} makes an object immutable --- you
  can't add, remove, or modify properties. This mimics the behavior of a
  Java \texttt{final} field or Python's \texttt{frozenset}. When would
  immutable objects be useful? What patterns become possible when you
  can guarantee an object won't change?''}
\item
  \emph{``Prototype-based inheritance means JavaScript objects can
  inherit from other objects directly, without classes. Classes
  (introduced in ES6) are `syntactic sugar' over this prototype system.
  What does this mean? What can you do with prototypes that you cannot
  easily express with class syntax?''}
\item
  \emph{``JSON (JavaScript Object Notation) is the universal data
  exchange format. It looks almost like JavaScript object literals but
  has stricter rules: all keys must be double-quoted strings, no
  functions, no comments. Why must JSON be stricter than JavaScript
  object literals? What security problems would arise if JSON allowed
  functions?''}
\end{enumerate}

\begin{center}\rule{0.5\linewidth}{0.5pt}\end{center}

\subsection{9. Unit 5: Strings}\label{unit-5-strings}

\paragraph{Instructional Cues}\label{instructional-cues-6}

\textbf{Template literals for multiline strings:}

\begin{Shaded}
\begin{Highlighting}[]
\KeywordTok{const}\NormalTok{ menu }\OperatorTok{=} \VerbatimStringTok{\textasciigrave{}}
\VerbatimStringTok{{-}{-}{-} Main Menu {-}{-}{-}}
\VerbatimStringTok{1. Add item}
\VerbatimStringTok{2. View all}
\VerbatimStringTok{3. Search}
\VerbatimStringTok{4. Delete}
\VerbatimStringTok{5. Quit}
\VerbatimStringTok{{-}{-}{-}{-}{-}{-}{-}{-}{-}{-}{-}{-}{-}{-}{-}{-}{-}\textasciigrave{}}\OperatorTok{;}

\BuiltInTok{console}\OperatorTok{.}\FunctionTok{log}\NormalTok{(menu)}\OperatorTok{;}
\end{Highlighting}
\end{Shaded}

Template literals with newlines are much cleaner for building multiline
output than repeated \texttt{console.log()} calls. Teach this pattern
for menus and reports.

\textbf{Regular expression basics --- connect to real need:}

Before introducing regex syntax, establish the problem it solves:

\emph{``You have user input that might have extra spaces, uppercase, or
punctuation. You want just the words. Without regex, you'd call
\texttt{.trim().toLowerCase().replace(\textquotesingle{},\textquotesingle{},\ \textquotesingle{}\textquotesingle{}).replace(\textquotesingle{}.\textquotesingle{},\ \textquotesingle{}\textquotesingle{})...}
for every punctuation mark. Regex lets you express `any non-letter
character' once.''}

\begin{Shaded}
\begin{Highlighting}[]
\KeywordTok{const}\NormalTok{ messy }\OperatorTok{=} \StringTok{"  Hello, World!   How are you?  "}\OperatorTok{;}
\KeywordTok{const}\NormalTok{ words }\OperatorTok{=}\NormalTok{ messy}
  \OperatorTok{.}\FunctionTok{trim}\NormalTok{()}
  \OperatorTok{.}\FunctionTok{toLowerCase}\NormalTok{()}
  \OperatorTok{.}\FunctionTok{split}\NormalTok{(}\SpecialStringTok{/}\SpecialCharTok{\textbackslash{}W+}\SpecialStringTok{/}\NormalTok{)}
  \OperatorTok{.}\FunctionTok{filter}\NormalTok{((w) }\KeywordTok{=\textgreater{}}\NormalTok{ w}\OperatorTok{.}\AttributeTok{length} \OperatorTok{\textgreater{}} \DecValTok{0}\NormalTok{)}\OperatorTok{;}
\BuiltInTok{console}\OperatorTok{.}\FunctionTok{log}\NormalTok{(words)}\OperatorTok{;} \CommentTok{// [\textquotesingle{}hello\textquotesingle{}, \textquotesingle{}world\textquotesingle{}, \textquotesingle{}how\textquotesingle{}, \textquotesingle{}are\textquotesingle{}, \textquotesingle{}you\textquotesingle{}]}
\end{Highlighting}
\end{Shaded}

Teach \texttt{\textbackslash{}w+} (word characters),
\texttt{\textbackslash{}d+} (digits), \texttt{\textbackslash{}s+}
(whitespace), and the \texttt{g} and \texttt{i} flags. Don't attempt to
teach full regex --- a basic vocabulary is enough for data cleaning
tasks.

\paragraph{Thought Experiments}\label{thought-experiments-6}

\begin{enumerate}
\def\labelenumi{\arabic{enumi}.}
\item
  \emph{``Template literals use backticks \texttt{\textasciigrave{}}
  while regular strings use \texttt{"} or \texttt{\textquotesingle{}}.
  Backticks on a keyboard are typically in the upper-left corner. For
  touch typists, backticks are less natural than quotes. What keyboard
  shortcuts or text expansions might a blind programmer use to reduce
  the friction of frequently-needed but awkward characters?''}
\item
  \emph{``String immutability in JavaScript (same as Python) means every
  \texttt{.toUpperCase()}, \texttt{.replace()}, and \texttt{.slice()}
  creates a new string. For a loop that builds a long string
  incrementally, what is the efficient alternative? (Hint: consider
  \texttt{Array.push()} + \texttt{.join()} as a form of
  StringBuilder.)''}
\end{enumerate}

\begin{center}\rule{0.5\linewidth}{0.5pt}\end{center}

\subsection{10. Unit 6: Functions}\label{unit-6-functions}

\subsubsection{Extended Instructional
Cues}\label{extended-instructional-cues}

\textbf{The four function forms --- teach one style, acknowledge the
rest:}

\begin{Shaded}
\begin{Highlighting}[]
\CommentTok{// 1. Function declaration (hoisted {-}{-}{-} can be called before declaration)}
\KeywordTok{function} \FunctionTok{add}\NormalTok{(a}\OperatorTok{,}\NormalTok{ b) \{}
  \ControlFlowTok{return}\NormalTok{ a }\OperatorTok{+}\NormalTok{ b}\OperatorTok{;}
\NormalTok{\}}

\CommentTok{// 2. Function expression (not hoisted)}
\KeywordTok{const}\NormalTok{ multiply }\OperatorTok{=} \KeywordTok{function}\NormalTok{ (a}\OperatorTok{,}\NormalTok{ b) \{}
  \ControlFlowTok{return}\NormalTok{ a }\OperatorTok{*}\NormalTok{ b}\OperatorTok{;}
\NormalTok{\}}\OperatorTok{;}

\CommentTok{// 3. Arrow function {-}{-}{-} SHORT BODY (implicit return)}
\KeywordTok{const}\NormalTok{ square }\OperatorTok{=}\NormalTok{ (n) }\KeywordTok{=\textgreater{}}\NormalTok{ n }\OperatorTok{*}\NormalTok{ n}\OperatorTok{;}

\CommentTok{// 4. Arrow function {-}{-}{-} LONG BODY (explicit return required)}
\KeywordTok{const}\NormalTok{ absolute }\OperatorTok{=}\NormalTok{ (n) }\KeywordTok{=\textgreater{}}\NormalTok{ \{}
  \ControlFlowTok{if}\NormalTok{ (n }\OperatorTok{\textless{}} \DecValTok{0}\NormalTok{) }\ControlFlowTok{return} \OperatorTok{{-}}\NormalTok{n}\OperatorTok{;}
  \ControlFlowTok{return}\NormalTok{ n}\OperatorTok{;}
\NormalTok{\}}\OperatorTok{;}
\end{Highlighting}
\end{Shaded}

\textbf{Teach style 3 (short arrow) and style 4 (long arrow) as the
primary forms.} Acknowledge 1 and 2 exist and how to read them. Do not
attempt to teach all four simultaneously.

\textbf{The implicit return trap --- a critical accessibility issue:}

\begin{Shaded}
\begin{Highlighting}[]
\CommentTok{// These DO return a value:}
\KeywordTok{const}\NormalTok{ double }\OperatorTok{=}\NormalTok{ (n) }\KeywordTok{=\textgreater{}}\NormalTok{ n }\OperatorTok{*} \DecValTok{2}\OperatorTok{;} \CommentTok{// [x] implicit return}
\KeywordTok{const}\NormalTok{ greet }\OperatorTok{=}\NormalTok{ (name) }\KeywordTok{=\textgreater{}} \VerbatimStringTok{\textasciigrave{}Hi }\SpecialCharTok{$\{}\NormalTok{name}\SpecialCharTok{\}}\VerbatimStringTok{\textasciigrave{}}\OperatorTok{;} \CommentTok{// [x] implicit return}

\CommentTok{// This does NOT return a value (the curly brace requires explicit return):}
\KeywordTok{const}\NormalTok{ broken }\OperatorTok{=}\NormalTok{ (n) }\KeywordTok{=\textgreater{}}\NormalTok{ \{}
\NormalTok{  n }\OperatorTok{*} \DecValTok{2}\OperatorTok{;} \CommentTok{// No return statement! Returns undefined.}
\NormalTok{\}}\OperatorTok{;}

\CommentTok{// Fixed:}
\KeywordTok{const}\NormalTok{ fixed }\OperatorTok{=}\NormalTok{ (n) }\KeywordTok{=\textgreater{}}\NormalTok{ \{}
  \ControlFlowTok{return}\NormalTok{ n }\OperatorTok{*} \DecValTok{2}\OperatorTok{;}
\NormalTok{\}}\OperatorTok{;}
\end{Highlighting}
\end{Shaded}

When heard aloud, \texttt{n\ =\textgreater{}\ \{\ n\ *\ 2;\ \}} and
\texttt{n\ =\textgreater{}\ \{\ return\ n\ *\ 2;\ \}} sound different
--- the \texttt{return} keyword is present in one and absent in the
other. Teach students to listen for \texttt{return} inside block-bodied
arrow functions.

\textbf{Higher-order functions --- connect to previous array methods:}

\begin{Shaded}
\begin{Highlighting}[]
\CommentTok{// map() takes a function as argument}
\KeywordTok{const}\NormalTok{ doubled }\OperatorTok{=}\NormalTok{ [}\DecValTok{1}\OperatorTok{,} \DecValTok{2}\OperatorTok{,} \DecValTok{3}\NormalTok{]}\OperatorTok{.}\FunctionTok{map}\NormalTok{((n) }\KeywordTok{=\textgreater{}}\NormalTok{ n }\OperatorTok{*} \DecValTok{2}\NormalTok{)}\OperatorTok{;}

\CommentTok{// We can also pass a named function}
\KeywordTok{const}\NormalTok{ double }\OperatorTok{=}\NormalTok{ (n) }\KeywordTok{=\textgreater{}}\NormalTok{ n }\OperatorTok{*} \DecValTok{2}\OperatorTok{;}
\KeywordTok{const}\NormalTok{ doubled2 }\OperatorTok{=}\NormalTok{ [}\DecValTok{1}\OperatorTok{,} \DecValTok{2}\OperatorTok{,} \DecValTok{3}\NormalTok{]}\OperatorTok{.}\FunctionTok{map}\NormalTok{(double)}\OperatorTok{;}

\CommentTok{// This works because map expects a function {-}{-}{-} it doesn\textquotesingle{}t care if it\textquotesingle{}s}
\CommentTok{// named or anonymous, arrow or traditional}
\end{Highlighting}
\end{Shaded}

\emph{``Every time you used \texttt{array.map(...)}, you were passing a
function as an argument. Higher-order functions --- functions that take
other functions --- are normal JavaScript.''}

\textbf{Closures --- establish the mental model before the code:}

\emph{``A closure is a function that remembers the environment where it
was created, even after that environment is gone. Imagine a letter in an
envelope. The letter is the function. The envelope is the environment
--- it contains variables the function needs. Even if the person who
wrote the letter is gone, the letter still has access to everything in
its envelope.''}

\begin{Shaded}
\begin{Highlighting}[]
\KeywordTok{function} \FunctionTok{makeCounter}\NormalTok{(startValue }\OperatorTok{=} \DecValTok{0}\NormalTok{) \{}
  \KeywordTok{let}\NormalTok{ count }\OperatorTok{=}\NormalTok{ startValue}\OperatorTok{;} \CommentTok{// This is the "envelope"}

  \ControlFlowTok{return}\NormalTok{ \{}
    \DataTypeTok{increment}\OperatorTok{:}\NormalTok{ () }\KeywordTok{=\textgreater{}} \OperatorTok{++}\NormalTok{count}\OperatorTok{,}
    \DataTypeTok{decrement}\OperatorTok{:}\NormalTok{ () }\KeywordTok{=\textgreater{}} \OperatorTok{{-}{-}}\NormalTok{count}\OperatorTok{,}
    \DataTypeTok{reset}\OperatorTok{:}\NormalTok{ () }\KeywordTok{=\textgreater{}}\NormalTok{ \{}
\NormalTok{      count }\OperatorTok{=}\NormalTok{ startValue}\OperatorTok{;}
\NormalTok{    \}}\OperatorTok{,}
    \DataTypeTok{getCount}\OperatorTok{:}\NormalTok{ () }\KeywordTok{=\textgreater{}}\NormalTok{ count}\OperatorTok{,}
\NormalTok{  \}}\OperatorTok{;}
\NormalTok{\}}

\KeywordTok{const}\NormalTok{ counter }\OperatorTok{=} \FunctionTok{makeCounter}\NormalTok{(}\DecValTok{10}\NormalTok{)}\OperatorTok{;}
\NormalTok{counter}\OperatorTok{.}\FunctionTok{increment}\NormalTok{()}\OperatorTok{;}
\NormalTok{counter}\OperatorTok{.}\FunctionTok{increment}\NormalTok{()}\OperatorTok{;}
\NormalTok{counter}\OperatorTok{.}\FunctionTok{increment}\NormalTok{()}\OperatorTok{;}
\BuiltInTok{console}\OperatorTok{.}\FunctionTok{log}\NormalTok{(counter}\OperatorTok{.}\FunctionTok{getCount}\NormalTok{())}\OperatorTok{;} \CommentTok{// 13 {-}{-}{-} count persists!}
\end{Highlighting}
\end{Shaded}

\paragraph{Thought Experiments}\label{thought-experiments-7}

\begin{enumerate}
\def\labelenumi{\arabic{enumi}.}
\item
  \emph{``Arrow functions don't have their own \texttt{this} --- they
  inherit \texttt{this} from the surrounding scope. Regular functions
  create their own \texttt{this}. This was a major source of bugs in
  pre-ES6 JavaScript. How many of the following do you think involved
  confusing \texttt{this}? (a) Typos, (b) Variable scope errors, (c)
  Async timing bugs, (d) \texttt{this} binding bugs. In JavaScript
  specifically, (d) is disproportionately common --- why?''}
\item
  \emph{``Pure functions always produce the same output for the same
  input and have no side effects. \texttt{Math.sin(x)} is pure.
  \texttt{console.log(x)} is not (it has a side effect). Why are pure
  functions easier to test? Why are pure functions easier to reason
  about in a screen-reader debugging session?''}
\item
  \emph{``Memoization caches the results of function calls so repeated
  calls with the same arguments return the cached result. Write a
  description of a memoized \texttt{fibonacci(n)} function. What is
  saved? What is the tradeoff? When would memoization be harmful rather
  than helpful?''}
\item
  \emph{``JavaScript is a multi-paradigm language: you can write
  imperative code (explicit step-by-step), object-oriented code (classes
  and methods), or functional code (pure functions and immutable data).
  For the accessible data manager project, which paradigm produces the
  most maintainable code? The most testable? The most readable when
  heard aloud?''}
\end{enumerate}

\begin{center}\rule{0.5\linewidth}{0.5pt}\end{center}

\subsection{11. Unit 7: Classes and OOP}\label{unit-7-classes-and-oop}

\paragraph{Instructional Cues}\label{instructional-cues-7}

\textbf{\texttt{toString()} is mandatory:}

\begin{Shaded}
\begin{Highlighting}[]
\KeywordTok{class}\NormalTok{ BankAccount \{}
\NormalTok{  \#balance}\OperatorTok{;}

  \FunctionTok{constructor}\NormalTok{(owner}\OperatorTok{,}\NormalTok{ initialBalance }\OperatorTok{=} \DecValTok{0}\NormalTok{) \{}
    \KeywordTok{this}\OperatorTok{.}\AttributeTok{owner} \OperatorTok{=}\NormalTok{ owner}\OperatorTok{;}
    \KeywordTok{this}\OperatorTok{.}\NormalTok{\#}\AttributeTok{balance} \OperatorTok{=}\NormalTok{ initialBalance}\OperatorTok{;}
\NormalTok{  \}}

  \FunctionTok{toString}\NormalTok{() \{}
    \ControlFlowTok{return} \VerbatimStringTok{\textasciigrave{}BankAccount[owner=}\SpecialCharTok{$\{}\KeywordTok{this}\OperatorTok{.}\AttributeTok{owner}\SpecialCharTok{\}}\VerbatimStringTok{, balance=$}\SpecialCharTok{$\{}\KeywordTok{this}\OperatorTok{.}\NormalTok{\#}\AttributeTok{balance}\OperatorTok{.}\FunctionTok{toFixed}\NormalTok{(}\DecValTok{2}\NormalTok{)}\SpecialCharTok{\}}\VerbatimStringTok{]\textasciigrave{}}\OperatorTok{;}
\NormalTok{  \}}
\NormalTok{\}}

\KeywordTok{const}\NormalTok{ account }\OperatorTok{=} \KeywordTok{new} \FunctionTok{BankAccount}\NormalTok{(}\StringTok{"Alice"}\OperatorTok{,} \DecValTok{1000}\NormalTok{)}\OperatorTok{;}
\BuiltInTok{console}\OperatorTok{.}\FunctionTok{log}\NormalTok{(account}\OperatorTok{.}\FunctionTok{toString}\NormalTok{())}\OperatorTok{;} \CommentTok{// BankAccount[owner=Alice, balance=$1000.00]}
\BuiltInTok{console}\OperatorTok{.}\FunctionTok{log}\NormalTok{(}\VerbatimStringTok{\textasciigrave{}}\SpecialCharTok{$\{}\NormalTok{account}\SpecialCharTok{\}}\VerbatimStringTok{\textasciigrave{}}\NormalTok{)}\OperatorTok{;} \CommentTok{// Also calls toString() implicitly}
\end{Highlighting}
\end{Shaded}

Require \texttt{toString()} in every class. Explain: ``When your screen
reader reads \texttt{console.log(myAccount)}, it reads
\texttt{toString()}. Without it, you hear \texttt{{[}object\ Object{]}}
--- completely useless. Define \texttt{toString()} to hear meaningful
information.''

\textbf{Private fields \texttt{\#} --- establish as the default:}

\begin{Shaded}
\begin{Highlighting}[]
\KeywordTok{class}\NormalTok{ Dog \{}
\NormalTok{  \#name}\OperatorTok{;} \CommentTok{// Private {-}{-}{-} cannot be accessed outside the class}
\NormalTok{  \#age}\OperatorTok{;}

  \FunctionTok{constructor}\NormalTok{(name}\OperatorTok{,}\NormalTok{ age) \{}
    \KeywordTok{this}\OperatorTok{.}\NormalTok{\#}\AttributeTok{name} \OperatorTok{=}\NormalTok{ name}\OperatorTok{;}
    \KeywordTok{this}\OperatorTok{.}\NormalTok{\#}\AttributeTok{age} \OperatorTok{=}\NormalTok{ age}\OperatorTok{;}
\NormalTok{  \}}

  \KeywordTok{get} \FunctionTok{name}\NormalTok{() \{}
    \ControlFlowTok{return} \KeywordTok{this}\OperatorTok{.}\NormalTok{\#}\AttributeTok{name}\OperatorTok{;}
\NormalTok{  \} }\CommentTok{// Getter {-}{-}{-} controlled access}
  \KeywordTok{get} \FunctionTok{age}\NormalTok{() \{}
    \ControlFlowTok{return} \KeywordTok{this}\OperatorTok{.}\NormalTok{\#}\AttributeTok{age}\OperatorTok{;}
\NormalTok{  \}}

  \KeywordTok{set} \FunctionTok{age}\NormalTok{(value) \{}
    \ControlFlowTok{if}\NormalTok{ (value }\OperatorTok{\textless{}} \DecValTok{0}\NormalTok{) }\ControlFlowTok{throw} \KeywordTok{new} \BuiltInTok{Error}\NormalTok{(}\StringTok{"Age cannot be negative"}\NormalTok{)}\OperatorTok{;}
    \KeywordTok{this}\OperatorTok{.}\NormalTok{\#}\AttributeTok{age} \OperatorTok{=}\NormalTok{ value}\OperatorTok{;}
\NormalTok{  \}}

  \FunctionTok{bark}\NormalTok{() \{}
    \BuiltInTok{console}\OperatorTok{.}\FunctionTok{log}\NormalTok{(}\VerbatimStringTok{\textasciigrave{}}\SpecialCharTok{$\{}\KeywordTok{this}\OperatorTok{.}\NormalTok{\#}\AttributeTok{name}\SpecialCharTok{\}}\VerbatimStringTok{ says: Woof!\textasciigrave{}}\NormalTok{)}\OperatorTok{;}
\NormalTok{  \}}

  \FunctionTok{toString}\NormalTok{() \{}
    \ControlFlowTok{return} \VerbatimStringTok{\textasciigrave{}Dog[name=}\SpecialCharTok{$\{}\KeywordTok{this}\OperatorTok{.}\NormalTok{\#}\AttributeTok{name}\SpecialCharTok{\}}\VerbatimStringTok{, age=}\SpecialCharTok{$\{}\KeywordTok{this}\OperatorTok{.}\NormalTok{\#}\AttributeTok{age}\SpecialCharTok{\}}\VerbatimStringTok{]\textasciigrave{}}\OperatorTok{;}
\NormalTok{  \}}
\NormalTok{\}}
\end{Highlighting}
\end{Shaded}

\emph{``Private fields starting with \texttt{\#} are truly private ---
not just conventional privacy. \texttt{myDog.\#name} outside the class
throws a SyntaxError. This is real encapsulation.''}

\paragraph{Thought Experiments}\label{thought-experiments-8}

\begin{enumerate}
\def\labelenumi{\arabic{enumi}.}
\item
  \emph{``JavaScript classes are `syntactic sugar' over prototype-based
  inheritance. TypeScript adds static typing on top of JavaScript
  classes. With three layers (prototypes -\textgreater{} ES6 classes
  -\textgreater{} TypeScript classes), is JavaScript's OOP model more or
  less trustworthy than Java's, which has only one layer? What does
  `syntactic sugar' imply about the behavior of the underlying
  mechanism?''}
\item
  \emph{``The \texttt{class} keyword was added to JavaScript in 2015, 20
  years after the language was created. In those 20 years, JavaScript
  developers built massive applications without classes using only
  prototypes and closures. What does this tell you about the role of
  syntax vs.~capability in programming languages?''}
\item
  \emph{``Mixins are a way to compose behavior from multiple sources
  without inheritance. Instead of \texttt{class\ A\ extends\ B}, you
  copy methods from B onto A's prototype. JavaScript supports this; Java
  does not (without interfaces). When would mixins be more flexible than
  single inheritance? What problems could they cause?''}
\end{enumerate}

\begin{center}\rule{0.5\linewidth}{0.5pt}\end{center}

\subsection{12. Unit 8: Files, Modules, and Error
Handling}\label{unit-8-files-modules-and-error-handling}

\subsubsection{Lesson 8.1: Files}\label{lesson-8.1-files}

\paragraph{Instructional Cues}\label{instructional-cues-8}

\textbf{\texttt{\_\_dirname} and path construction:}

\begin{Shaded}
\begin{Highlighting}[]
\KeywordTok{const}\NormalTok{ path }\OperatorTok{=} \PreprocessorTok{require}\NormalTok{(}\StringTok{"path"}\NormalTok{)}\OperatorTok{;}
\KeywordTok{const}\NormalTok{ fs }\OperatorTok{=} \PreprocessorTok{require}\NormalTok{(}\StringTok{"fs"}\NormalTok{)}\OperatorTok{;}

\CommentTok{// \_\_dirname is the directory of the current file}
\KeywordTok{const}\NormalTok{ dataFile }\OperatorTok{=}\NormalTok{ path}\OperatorTok{.}\FunctionTok{join}\NormalTok{(}\BuiltInTok{\_\_dirname}\OperatorTok{,} \StringTok{"data"}\OperatorTok{,} \StringTok{"contacts.json"}\NormalTok{)}\OperatorTok{;}

\CommentTok{// Path joining prevents slash vs. backslash issues on Windows/macOS/Linux}
\BuiltInTok{console}\OperatorTok{.}\FunctionTok{log}\NormalTok{(}\StringTok{"Data file location:"}\OperatorTok{,}\NormalTok{ dataFile)}\OperatorTok{;}
\end{Highlighting}
\end{Shaded}

For a student who cannot visually navigate a file tree,
\texttt{console.log(\_\_dirname)} and \texttt{console.log(dataFile)} are
essential diagnostic tools. Teach them to print file paths as the first
debugging step when a file isn't found.

\textbf{\texttt{fs.existsSync()} before file operations:}

\begin{Shaded}
\begin{Highlighting}[]
\ControlFlowTok{if}\NormalTok{ (fs}\OperatorTok{.}\FunctionTok{existsSync}\NormalTok{(dataFile)) \{}
  \KeywordTok{const}\NormalTok{ data }\OperatorTok{=}\NormalTok{ fs}\OperatorTok{.}\FunctionTok{readFileSync}\NormalTok{(dataFile}\OperatorTok{,} \StringTok{"utf8"}\NormalTok{)}\OperatorTok{;}
  \ControlFlowTok{return} \BuiltInTok{JSON}\OperatorTok{.}\FunctionTok{parse}\NormalTok{(data)}\OperatorTok{;}
\NormalTok{\} }\ControlFlowTok{else}\NormalTok{ \{}
  \BuiltInTok{console}\OperatorTok{.}\FunctionTok{log}\NormalTok{(}\StringTok{"No data file found. Starting with empty data."}\NormalTok{)}\OperatorTok{;}
  \ControlFlowTok{return}\NormalTok{ []}\OperatorTok{;}
\NormalTok{\}}
\end{Highlighting}
\end{Shaded}

Always check existence before reading. The alternative --- catching
\texttt{ENOENT} in a try-catch --- works but puts normal ``first run''
behavior in the error path.

\subsubsection{Lesson 8.3: CommonJS
Modules}\label{lesson-8.3-commonjs-modules}

\textbf{Module design for accessibility:}

A well-designed module has a clear exports interface. Teach students to
think of their module's exports as its ``public API'' --- what other
files can use:

\begin{Shaded}
\begin{Highlighting}[]
\CommentTok{// contacts.js {-}{-}{-} exports only what other files need}
\KeywordTok{const}\NormalTok{ path }\OperatorTok{=} \PreprocessorTok{require}\NormalTok{(}\StringTok{"path"}\NormalTok{)}\OperatorTok{;}
\KeywordTok{const}\NormalTok{ fs }\OperatorTok{=} \PreprocessorTok{require}\NormalTok{(}\StringTok{"fs"}\NormalTok{)}\OperatorTok{;}

\KeywordTok{const}\NormalTok{ DATA\_FILE }\OperatorTok{=}\NormalTok{ path}\OperatorTok{.}\FunctionTok{join}\NormalTok{(}\BuiltInTok{\_\_dirname}\OperatorTok{,} \StringTok{".."}\OperatorTok{,} \StringTok{"data"}\OperatorTok{,} \StringTok{"contacts.json"}\NormalTok{)}\OperatorTok{;}

\KeywordTok{function} \FunctionTok{loadContacts}\NormalTok{() \{}
  \CommentTok{/* ... */}
\NormalTok{\}}
\KeywordTok{function} \FunctionTok{saveContacts}\NormalTok{(contacts) \{}
  \CommentTok{/* ... */}
\NormalTok{\}}
\KeywordTok{function} \FunctionTok{findContact}\NormalTok{(contacts}\OperatorTok{,}\NormalTok{ name) \{}
  \CommentTok{/* ... */}
\NormalTok{\}}

\CommentTok{// Export only the public interface}
\NormalTok{module}\OperatorTok{.}\AttributeTok{exports} \OperatorTok{=}\NormalTok{ \{ loadContacts}\OperatorTok{,}\NormalTok{ saveContacts}\OperatorTok{,}\NormalTok{ findContact \}}\OperatorTok{;}
\end{Highlighting}
\end{Shaded}

\emph{``Everything not in \texttt{module.exports} is private to this
file. Other files can use \texttt{loadContacts} and
\texttt{saveContacts}, but they cannot directly access
\texttt{DATA\_FILE}. This is encapsulation at the module level.''}

\paragraph{Thought Experiments}\label{thought-experiments-9}

\begin{enumerate}
\def\labelenumi{\arabic{enumi}.}
\item
  \emph{``Node.js's CommonJS (\texttt{require}) and the ES Module
  standard (\texttt{import}) are two different module systems in
  JavaScript. This creates confusion --- some packages use one, some use
  the other. TypeScript can use both. Why did two module systems emerge?
  What does this history say about how standards develop in open-source
  ecosystems?''}
\item
  \emph{``Synchronous file I/O (\texttt{readFileSync}) blocks execution
  until the file is read. Asynchronous file I/O (\texttt{readFile} or
  \texttt{fs.promises.readFile}) allows other code to run while waiting.
  For a startup script that reads configuration before doing anything
  else, which makes more sense? For a web server that handles thousands
  of requests, which matters more?''}
\item
  \emph{``Custom error classes allow code to throw and catch specific,
  named errors:
  \texttt{throw\ new\ ValidationError(\textquotesingle{}email\textquotesingle{},\ \textquotesingle{}Invalid\ format\textquotesingle{})}.
  What is the advantage of this over throwing
  \texttt{new\ Error(\textquotesingle{}Validation\ failed:\ email\ field\ invalid\ format\textquotesingle{})}?
  Think about code that catches and handles the error.''}
\end{enumerate}

\begin{center}\rule{0.5\linewidth}{0.5pt}\end{center}

\subsection{13. Unit 9: Collections, Analysis, and
Testing}\label{unit-9-collections-analysis-and-testing}

\subsubsection{Jest Testing}\label{jest-testing}

\paragraph{Instructional Cues}\label{instructional-cues-9}

\textbf{Jest output is accessible:}

Run \texttt{npx\ jest\ -\/-verbose} for maximum accessibility. The
\texttt{-\/-verbose} flag produces:

\begin{verbatim}
PASS src/math.test.js
  add function
    [x] adds two positive numbers (3ms)
    [x] adds negative numbers (1ms)
    [x] handles zero (0ms)
  getGrade function
    [x] returns A for score >= 90 (0ms)
    [x] returns F for score < 60 (0ms)

Test Suites: 1 passed, 1 total
Tests:       5 passed, 5 total
\end{verbatim}

Each test name is announced. Failures include file path, test name,
expected vs.~received values, and a diff. All of this is plain text ---
fully screen-reader-accessible.

\textbf{Add to \texttt{package.json} scripts:}

\begin{Shaded}
\begin{Highlighting}[]
\StringTok{"scripts"}\ErrorTok{:} \FunctionTok{\{}
    \DataTypeTok{"test"}\FunctionTok{:} \StringTok{"jest {-}{-}verbose"}\FunctionTok{,}
    \DataTypeTok{"test:watch"}\FunctionTok{:} \StringTok{"jest {-}{-}verbose {-}{-}watch"}
\FunctionTok{\}}
\end{Highlighting}
\end{Shaded}

\texttt{npm\ test} runs all tests. \texttt{npm\ run\ test:watch} reruns
tests whenever a file changes --- useful during active development.

\textbf{Test naming convention:}

\begin{Shaded}
\begin{Highlighting}[]
\FunctionTok{describe}\NormalTok{(}\StringTok{"BankAccount"}\OperatorTok{,}\NormalTok{ () }\KeywordTok{=\textgreater{}}\NormalTok{ \{}
  \FunctionTok{describe}\NormalTok{(}\StringTok{"deposit"}\OperatorTok{,}\NormalTok{ () }\KeywordTok{=\textgreater{}}\NormalTok{ \{}
    \FunctionTok{test}\NormalTok{(}\StringTok{"increases balance by deposited amount"}\OperatorTok{,}\NormalTok{ () }\KeywordTok{=\textgreater{}}\NormalTok{ \{}
      \CommentTok{/* ... */}
\NormalTok{    \})}\OperatorTok{;}
    \FunctionTok{test}\NormalTok{(}\StringTok{"rejects negative amounts"}\OperatorTok{,}\NormalTok{ () }\KeywordTok{=\textgreater{}}\NormalTok{ \{}
      \CommentTok{/* ... */}
\NormalTok{    \})}\OperatorTok{;}
    \FunctionTok{test}\NormalTok{(}\StringTok{"rejects zero amount"}\OperatorTok{,}\NormalTok{ () }\KeywordTok{=\textgreater{}}\NormalTok{ \{}
      \CommentTok{/* ... */}
\NormalTok{    \})}\OperatorTok{;}
\NormalTok{  \})}\OperatorTok{;}
  \FunctionTok{describe}\NormalTok{(}\StringTok{"withdraw"}\OperatorTok{,}\NormalTok{ () }\KeywordTok{=\textgreater{}}\NormalTok{ \{}
    \FunctionTok{test}\NormalTok{(}\StringTok{"decreases balance by withdrawn amount"}\OperatorTok{,}\NormalTok{ () }\KeywordTok{=\textgreater{}}\NormalTok{ \{}
      \CommentTok{/* ... */}
\NormalTok{    \})}\OperatorTok{;}
    \FunctionTok{test}\NormalTok{(}\StringTok{"throws InsufficientFundsError when balance is too low"}\OperatorTok{,}\NormalTok{ () }\KeywordTok{=\textgreater{}}\NormalTok{ \{}
      \CommentTok{/* ... */}
\NormalTok{    \})}\OperatorTok{;}
\NormalTok{  \})}\OperatorTok{;}
\NormalTok{\})}\OperatorTok{;}
\end{Highlighting}
\end{Shaded}

The nested \texttt{describe} / \texttt{test} structure, when announced
by Jest with \texttt{-\/-verbose}, reads as complete sentences:
``BankAccount \textgreater{} deposit \textgreater{} increases balance by
deposited amount.''

\paragraph{Thought Experiments}\label{thought-experiments-10}

\begin{enumerate}
\def\labelenumi{\arabic{enumi}.}
\item
  \emph{``Jest's \texttt{expect(value).toEqual(expected)} performs deep
  comparison --- it checks that two objects have the same structure and
  values, not that they are the same object in memory. Why is deep
  comparison usually what you want when testing? When would you
  specifically want reference equality (\texttt{toBe})?''}
\item
  \emph{``Snapshot testing in Jest captures the output of a component or
  function the first time, then checks future runs against that
  snapshot. What are the advantages? What are the risks? When might a
  snapshot become `stale' and mislead you into thinking code is correct
  when it isn't?''}
\item
  \emph{``Code coverage measures what percentage of your code is
  executed by tests. 100\% coverage is a goal but not a guarantee of
  correctness. Design a function with 100\% test coverage that still has
  a bug. (Hint: think about boundary conditions --- what if the test
  covers the case where x=5 but not x=0?)''}
\end{enumerate}

\begin{center}\rule{0.5\linewidth}{0.5pt}\end{center}

\subsection{14. Units 10--11: Projects and
Capstone}\label{units-1011-projects-and-capstone}

\subsubsection{JavaScript Project
Structure}\label{javascript-project-structure}

\begin{verbatim}
my-data-manager/
|--- src/
|   |--- models/
|   |   +--- contact.js
|   |--- storage/
|   |   +--- dataManager.js
|   |--- ui/
|   |   +--- consoleUI.js
|   +--- utils/
|       |--- input.js         (the readline helper)
|       +--- validation.js    (reusable validation functions)
|--- data/
|   +--- contacts.json
|--- tests/
|   |--- models/
|   |   +--- contact.test.js
|   +--- storage/
|       +--- dataManager.test.js
|--- .eslintrc.json
|--- package.json
+--- index.js
\end{verbatim}

\textbf{\texttt{package.json} scripts --- teach as accessibility tools:}

\begin{Shaded}
\begin{Highlighting}[]
\FunctionTok{\{}
  \DataTypeTok{"scripts"}\FunctionTok{:} \FunctionTok{\{}
    \DataTypeTok{"start"}\FunctionTok{:} \StringTok{"node index.js"}\FunctionTok{,}
    \DataTypeTok{"test"}\FunctionTok{:} \StringTok{"jest {-}{-}verbose"}\FunctionTok{,}
    \DataTypeTok{"test:watch"}\FunctionTok{:} \StringTok{"jest {-}{-}verbose {-}{-}watch"}\FunctionTok{,}
    \DataTypeTok{"lint"}\FunctionTok{:} \StringTok{"eslint src/"}\FunctionTok{,}
    \DataTypeTok{"lint:fix"}\FunctionTok{:} \StringTok{"eslint src/ {-}{-}fix"}
  \FunctionTok{\}}
\FunctionTok{\}}
\end{Highlighting}
\end{Shaded}

\texttt{npm\ start} --- consistent command regardless of entry point
filename. \texttt{npm\ test} --- runs all tests. \texttt{npm\ run\ lint}
--- checks for errors without running the program.

\textbf{\texttt{npm\ run\ lint} as a pre-run check:} Teach students to
run \texttt{npm\ run\ lint} before running the program. ESLint catches
undefined variables, unused variables, and \texttt{==} before the
program runs --- much faster than finding runtime errors.

\subsubsection{Accessible Output Design for
Capstone}\label{accessible-output-design-for-capstone}

Require students to design their program's output to be navigable and
self-describing:

\begin{Shaded}
\begin{Highlighting}[]
\CommentTok{// Navigation landmark prints}
\KeywordTok{function} \FunctionTok{printSeparator}\NormalTok{(title) \{}
  \KeywordTok{const}\NormalTok{ line }\OperatorTok{=} \StringTok{"{-}"}\OperatorTok{.}\FunctionTok{repeat}\NormalTok{(}\DecValTok{40}\NormalTok{)}\OperatorTok{;}
  \BuiltInTok{console}\OperatorTok{.}\FunctionTok{log}\NormalTok{(line)}\OperatorTok{;}
  \ControlFlowTok{if}\NormalTok{ (title) }\BuiltInTok{console}\OperatorTok{.}\FunctionTok{log}\NormalTok{(}\VerbatimStringTok{\textasciigrave{}  }\SpecialCharTok{$\{}\NormalTok{title}\OperatorTok{.}\FunctionTok{toUpperCase}\NormalTok{()}\SpecialCharTok{\}}\VerbatimStringTok{\textasciigrave{}}\NormalTok{)}\OperatorTok{;}
  \BuiltInTok{console}\OperatorTok{.}\FunctionTok{log}\NormalTok{(line)}\OperatorTok{;}
\NormalTok{\}}

\CommentTok{// Numbered list with clear field labels}
\KeywordTok{function} \FunctionTok{printContactList}\NormalTok{(contacts) \{}
  \ControlFlowTok{if}\NormalTok{ (contacts}\OperatorTok{.}\AttributeTok{length} \OperatorTok{===} \DecValTok{0}\NormalTok{) \{}
    \BuiltInTok{console}\OperatorTok{.}\FunctionTok{log}\NormalTok{(}\StringTok{"No contacts found."}\NormalTok{)}\OperatorTok{;}
    \ControlFlowTok{return}\OperatorTok{;}
\NormalTok{  \}}
  \FunctionTok{printSeparator}\NormalTok{(}\StringTok{"contacts"}\NormalTok{)}\OperatorTok{;}
\NormalTok{  contacts}\OperatorTok{.}\FunctionTok{forEach}\NormalTok{((c}\OperatorTok{,}\NormalTok{ i) }\KeywordTok{=\textgreater{}}\NormalTok{ \{}
    \BuiltInTok{console}\OperatorTok{.}\FunctionTok{log}\NormalTok{(}\VerbatimStringTok{\textasciigrave{}}\SpecialCharTok{$\{}\NormalTok{i }\OperatorTok{+} \DecValTok{1}\SpecialCharTok{\}}\VerbatimStringTok{. }\SpecialCharTok{$\{}\NormalTok{c}\OperatorTok{.}\AttributeTok{name}\SpecialCharTok{\}}\VerbatimStringTok{\textasciigrave{}}\NormalTok{)}\OperatorTok{;}
    \BuiltInTok{console}\OperatorTok{.}\FunctionTok{log}\NormalTok{(}\VerbatimStringTok{\textasciigrave{}   Phone: }\SpecialCharTok{$\{}\NormalTok{c}\OperatorTok{.}\AttributeTok{phone} \OperatorTok{||} \StringTok{"not provided"}\SpecialCharTok{\}}\VerbatimStringTok{\textasciigrave{}}\NormalTok{)}\OperatorTok{;}
    \BuiltInTok{console}\OperatorTok{.}\FunctionTok{log}\NormalTok{(}\VerbatimStringTok{\textasciigrave{}   Email: }\SpecialCharTok{$\{}\NormalTok{c}\OperatorTok{.}\AttributeTok{email} \OperatorTok{||} \StringTok{"not provided"}\SpecialCharTok{\}}\VerbatimStringTok{\textasciigrave{}}\NormalTok{)}\OperatorTok{;}
\NormalTok{  \})}\OperatorTok{;}
  \BuiltInTok{console}\OperatorTok{.}\FunctionTok{log}\NormalTok{(}\VerbatimStringTok{\textasciigrave{}Total: }\SpecialCharTok{$\{}\NormalTok{contacts}\OperatorTok{.}\AttributeTok{length}\SpecialCharTok{\}}\VerbatimStringTok{ contact(s)\textasciigrave{}}\NormalTok{)}\OperatorTok{;}
\NormalTok{\}}

\CommentTok{// Confirmation messages {-}{-}{-} every action must produce one}
\KeywordTok{function} \FunctionTok{confirmAction}\NormalTok{(action}\OperatorTok{,}\NormalTok{ details) \{}
  \BuiltInTok{console}\OperatorTok{.}\FunctionTok{log}\NormalTok{(}\VerbatimStringTok{\textasciigrave{}[x] }\SpecialCharTok{$\{}\NormalTok{action}\SpecialCharTok{\}}\VerbatimStringTok{: }\SpecialCharTok{$\{}\NormalTok{details}\SpecialCharTok{\}}\VerbatimStringTok{\textasciigrave{}}\NormalTok{)}\OperatorTok{;}
\NormalTok{\}}
\end{Highlighting}
\end{Shaded}

\paragraph{Thought Experiments}\label{thought-experiments-11}

\begin{enumerate}
\def\labelenumi{\arabic{enumi}.}
\item
  \emph{``Your capstone is a command-line program. Eventually, it could
  become a web application. What would need to change? What would
  transfer directly? Write a two-column list: `Can reuse as-is' and
  `Must rewrite.'\,''}
\item
  \emph{``The accessible design principle `no visual-only cues for
  critical information' means every status, error, or confirmation must
  be communicated through text, not color or icon alone. Where in your
  program have you relied on visual formatting? How would you replace
  it?''}
\item
  \emph{``Open-source your capstone project. Write a README that
  explains: what the program does, how to install it, how to run it, and
  specifically how to use it with a screen reader. What information
  would a sighted user find obvious but a blind user needs explicitly
  stated?''}
\end{enumerate}

\begin{center}\rule{0.5\linewidth}{0.5pt}\end{center}

\subsection{15. Bridging from Python/Java to
JavaScript}\label{bridging-from-pythonjava-to-javascript}

\subsubsection{From Python to
JavaScript}\label{from-python-to-javascript}

{\def\LTcaptype{none} % do not increment counter
\begin{longtable}[]{@{}
  >{\raggedright\arraybackslash}p{(\linewidth - 4\tabcolsep) * \real{0.2745}}
  >{\raggedright\arraybackslash}p{(\linewidth - 4\tabcolsep) * \real{0.4020}}
  >{\raggedright\arraybackslash}p{(\linewidth - 4\tabcolsep) * \real{0.3235}}@{}}
\toprule\noalign{}
\begin{minipage}[b]{\linewidth}\raggedright
Python
\end{minipage} & \begin{minipage}[b]{\linewidth}\raggedright
JavaScript
\end{minipage} & \begin{minipage}[b]{\linewidth}\raggedright
Key Difference
\end{minipage} \\
\midrule\noalign{}
\endhead
\bottomrule\noalign{}
\endlastfoot
\texttt{print("hello")} & \texttt{console.log("hello")} & Method
vs.~function \\
\texttt{input("prompt")} & \texttt{await\ askQuestion("prompt")} & Async
in JavaScript \\
\texttt{f"Hello\ \{name\}"} &
\texttt{\textasciigrave{}Hello\ \$\{name\}\textasciigrave{}} & f-string
vs.~template literal \\
\texttt{def\ func(x):} &
\texttt{const\ func\ =\ x\ =\textgreater{}\ ...} & Arrow function
style \\
\texttt{True}/\texttt{False} & \texttt{true}/\texttt{false} & Lowercase
in JavaScript \\
\texttt{None} & \texttt{null} and \texttt{undefined} & JS has two
``nothing'' values \\
\texttt{my\_list} & \texttt{myList} & JavaScript uses camelCase \\
\texttt{len(x)} & \texttt{x.length} & Property, not function \\
\texttt{x\ is\ None} & \texttt{x\ ===\ null} or
\texttt{x\ ===\ undefined} & Different null checks \\
\texttt{import\ json} &
\texttt{const\ fs\ =\ require(\textquotesingle{}fs\textquotesingle{})} &
Built-in \texttt{JSON} object, no import \\
\texttt{if\ \_\_name\_\_\ ==\ \textquotesingle{}\_\_main\_\_\textquotesingle{}:}
& No direct equivalent (just call \texttt{main()}) & \\
\end{longtable}
}

\textbf{Thought experiments for Python-\textgreater JS students:}

\begin{enumerate}
\def\labelenumi{\arabic{enumi}.}
\item
  \emph{``In Python, you use spaces to define blocks (indentation). In
  JavaScript, you use \texttt{\{\}}. Python's blocks are invisible ---
  they're just the alignment of characters. JavaScript's blocks are
  explicitly marked. Which feels more natural? Which provides clearer
  information when read aloud?''}
\item
  \emph{``Python's \texttt{None} and JavaScript's
  \texttt{null}/\texttt{undefined} both represent `nothing,' but they
  behave differently. Python has one `nothing'; JavaScript has two. When
  is \texttt{null} appropriate? When is \texttt{undefined} appropriate?
  Under what circumstances would you encounter one vs.~the other?''}
\item
  \emph{``Python's type system and JavaScript's both use dynamic typing,
  but JavaScript's is more permissive --- it coerces types rather than
  raising errors. After using both, which approach do you find more
  helpful for learning? Which is more appropriate for production
  software?''}
\end{enumerate}

\begin{center}\rule{0.5\linewidth}{0.5pt}\end{center}

\subsection{16. Assessment Strategies}\label{assessment-strategies}

\subsubsection{Alternative Assessments for
JavaScript}\label{alternative-assessments-for-javascript}

\textbf{Console output analysis:} Provide a JavaScript program and 3--4
sample inputs. Ask the student to predict what the program prints for
each input before running it. This tests reading comprehension without
requiring writing.

\textbf{Linting-first review:} Submit code that is free of all ESLint
warnings. A student who achieves this has internalized JavaScript's best
practices at the syntax level.

\textbf{Refactoring exercises:} Give a working but poorly structured
program. Ask the student to refactor it to use functions, modules, or
classes appropriately. Assess the quality of the architecture, not just
whether it runs.

\textbf{Test-first implementation:} Give the student a set of failing
tests. Ask them to write the implementation that makes the tests pass.
This tests code comprehension and problem-solving without a blank-page
start.

\begin{center}\rule{0.5\linewidth}{0.5pt}\end{center}

\subsection{17. Supporting Diverse
Needs}\label{supporting-diverse-needs}

\subsubsection{For Students Who Are New to Async
Programming}\label{for-students-who-are-new-to-async-programming}

Async/await is conceptually hard. For students who struggle:

\begin{enumerate}
\def\labelenumi{\arabic{enumi}.}
\tightlist
\item
  Provide the complete \texttt{input.js} helper and instruct them to use
  it without modifying it.
\item
  For the first several lessons, treat
  \texttt{async\ function\ main()\ \{\ ...\ \}\ main();} as boilerplate,
  not content.
\item
  Introduce the conceptual explanation of async/await only after
  students are comfortable with the rest of the language.
\item
  Consider deferring async until after basic functions are mastered.
\end{enumerate}

\subsubsection{For Students with Additional Cognitive
Considerations}\label{for-students-with-additional-cognitive-considerations}

\begin{itemize}
\tightlist
\item
  ESLint provides immediate text feedback for common errors --- use it
  as a learning tool. Hearing ``no-unused-vars'' and then explaining
  what that means is a teaching moment triggered by the student's own
  code.
\item
  JavaScript's many syntactic variations can be overwhelming. Enforce a
  single consistent style (arrow functions, \texttt{const}/\texttt{let},
  \texttt{===}) and defer alternatives.
\item
  Use the REPL liberally for short experiments so the student gets
  immediate feedback without file management overhead.
\end{itemize}

\begin{center}\rule{0.5\linewidth}{0.5pt}\end{center}

\subsection{18. Alternative and Supplementary
Resources}\label{alternative-and-supplementary-resources}

\subsubsection{Textbooks and Reading}\label{textbooks-and-reading}

{\def\LTcaptype{none} % do not increment counter
\begin{longtable}[]{@{}
  >{\raggedright\arraybackslash}p{(\linewidth - 6\tabcolsep) * \real{0.3529}}
  >{\raggedright\arraybackslash}p{(\linewidth - 6\tabcolsep) * \real{0.1618}}
  >{\raggedright\arraybackslash}p{(\linewidth - 6\tabcolsep) * \real{0.0956}}
  >{\raggedright\arraybackslash}p{(\linewidth - 6\tabcolsep) * \real{0.3897}}@{}}
\toprule\noalign{}
\begin{minipage}[b]{\linewidth}\raggedright
Resource
\end{minipage} & \begin{minipage}[b]{\linewidth}\raggedright
Level
\end{minipage} & \begin{minipage}[b]{\linewidth}\raggedright
Format
\end{minipage} & \begin{minipage}[b]{\linewidth}\raggedright
Accessibility
\end{minipage} \\
\midrule\noalign{}
\endhead
\bottomrule\noalign{}
\endlastfoot
\textbf{Eloquent JavaScript} (Haverbeke, free online) &
Beginner--Intermediate & Free HTML & Excellent screen-reader-friendly
HTML \\
\textbf{You Don't Know JS} (Simpson, free on GitHub) & Intermediate &
Free HTML/PDF & In-depth on closures, scope, types \\
\textbf{JavaScript: The Good Parts} (Crockford) & Intermediate & Print &
Short; focuses on the reliable subset of JS \\
\textbf{JavaScript: The Definitive Guide} (Flanagan) & Reference &
Print/ebook & Comprehensive; good for looking up specifics \\
\textbf{MDN Web Docs} (developer.mozilla.org) & Reference & Free HTML &
Excellent HTML accessibility; definitive JS reference \\
\textbf{Node.js Documentation} (nodejs.org/docs) & Reference & Free HTML
& Official Node.js API reference \\
\textbf{JavaScript.info} & Beginner--Advanced & Free HTML & Very
well-organized; plain HTML; comprehensive \\
\end{longtable}
}

\textbf{javascript.info is especially recommended} --- it's a
comprehensive free tutorial organized from basics to advanced topics,
all in clean, accessible HTML. It's often more readable than MDN for
learning (vs.~reference), and screen readers navigate it very well.

\subsubsection{Online Courses and
Platforms}\label{online-courses-and-platforms}

{\def\LTcaptype{none} % do not increment counter
\begin{longtable}[]{@{}
  >{\raggedright\arraybackslash}p{(\linewidth - 6\tabcolsep) * \real{0.1984}}
  >{\raggedright\arraybackslash}p{(\linewidth - 6\tabcolsep) * \real{0.3095}}
  >{\raggedright\arraybackslash}p{(\linewidth - 6\tabcolsep) * \real{0.1032}}
  >{\raggedright\arraybackslash}p{(\linewidth - 6\tabcolsep) * \real{0.3889}}@{}}
\toprule\noalign{}
\begin{minipage}[b]{\linewidth}\raggedright
Platform
\end{minipage} & \begin{minipage}[b]{\linewidth}\raggedright
Course
\end{minipage} & \begin{minipage}[b]{\linewidth}\raggedright
Accessibility
\end{minipage} & \begin{minipage}[b]{\linewidth}\raggedright
Notes
\end{minipage} \\
\midrule\noalign{}
\endhead
\bottomrule\noalign{}
\endlastfoot
\textbf{freeCodeCamp} & JavaScript Algorithms \& Data Structures &
Medium & Certification path; some UI inaccessible \\
\textbf{Codecademy} & Learn JavaScript & Partial & Editor accessible;
some UI elements not \\
\textbf{Exercism.io} & JavaScript Track & High & Terminal-based; mentor
feedback \\
\textbf{JavaScript30} (Wesbos) & 30 vanilla JS projects & Medium &
Video-based; transcripts available \\
\textbf{edX} & Various JavaScript courses & High & Video transcripts;
accessible exercises \\
\textbf{Scrimba} & Learn JavaScript & Low & Interactive video; poor
screen reader support \\
\textbf{Khan Academy} & JavaScript course & Medium & Browser-based
editor; variable accessibility \\
\textbf{NodeSchool.io} & Node.js workshops & High & Command-line
workshopper format; fully accessible \\
\end{longtable}
}

\textbf{NodeSchool.io is especially recommended} --- workshops are run
entirely in the terminal via a command-line tool. Students work through
exercises, get feedback from the tool, and progress at their own pace.
The terminal-first design makes it highly accessible.

\subsubsection{Video and Audio
Resources}\label{video-and-audio-resources}

{\def\LTcaptype{none} % do not increment counter
\begin{longtable}[]{@{}
  >{\raggedright\arraybackslash}p{(\linewidth - 4\tabcolsep) * \real{0.4240}}
  >{\raggedright\arraybackslash}p{(\linewidth - 4\tabcolsep) * \real{0.3360}}
  >{\raggedright\arraybackslash}p{(\linewidth - 4\tabcolsep) * \real{0.2400}}@{}}
\toprule\noalign{}
\begin{minipage}[b]{\linewidth}\raggedright
Resource
\end{minipage} & \begin{minipage}[b]{\linewidth}\raggedright
Description
\end{minipage} & \begin{minipage}[b]{\linewidth}\raggedright
Accessibility
\end{minipage} \\
\midrule\noalign{}
\endhead
\bottomrule\noalign{}
\endlastfoot
\textbf{Traversy Media} (YouTube) & JavaScript crash courses and
projects & Captions available \\
\textbf{The Net Ninja} (YouTube) & In-depth JavaScript series & Captions
available \\
\textbf{Academind} (YouTube) & Deep dives into JS/Node & Captions
available \\
\textbf{Kevin Powell} (YouTube) & JavaScript and web fundamentals &
Captions; clear explanations \\
\textbf{CS50W} (Harvard) & Web programming with Python and JavaScript &
Full transcripts; high quality \\
\textbf{JavaScript: Understanding the Weird Parts} (Udemy) & Deep JS
internals & Captions \\
\end{longtable}
}

\subsubsection{JavaScript-Specific Accessibility
Tools}\label{javascript-specific-accessibility-tools}

{\def\LTcaptype{none} % do not increment counter
\begin{longtable}[]{@{}
  >{\raggedright\arraybackslash}p{(\linewidth - 4\tabcolsep) * \real{0.1778}}
  >{\raggedright\arraybackslash}p{(\linewidth - 4\tabcolsep) * \real{0.3667}}
  >{\raggedright\arraybackslash}p{(\linewidth - 4\tabcolsep) * \real{0.4556}}@{}}
\toprule\noalign{}
\begin{minipage}[b]{\linewidth}\raggedright
Tool
\end{minipage} & \begin{minipage}[b]{\linewidth}\raggedright
Purpose
\end{minipage} & \begin{minipage}[b]{\linewidth}\raggedright
Notes
\end{minipage} \\
\midrule\noalign{}
\endhead
\bottomrule\noalign{}
\endlastfoot
\textbf{ESLint} & Code quality and error prevention & Command-line;
fully accessible \\
\textbf{Prettier} & Consistent code formatting & Run on save; no screen
interaction needed \\
\textbf{Jest} & Testing framework & Accessible output;
\texttt{-\/-verbose} flag \\
\textbf{jshint.com} & Online JS linter & Plain HTML; basic linting \\
\textbf{node-inspect} & Node.js debugger & Terminal-based debugging \\
\textbf{0x} & Node.js profiler & Terminal output \\
\textbf{nvm} & Node version manager & Command-line; accessible \\
\end{longtable}
}

\subsubsection{Communities}\label{communities}

{\def\LTcaptype{none} % do not increment counter
\begin{longtable}[]{@{}
  >{\raggedright\arraybackslash}p{(\linewidth - 2\tabcolsep) * \real{0.4545}}
  >{\raggedright\arraybackslash}p{(\linewidth - 2\tabcolsep) * \real{0.5455}}@{}}
\toprule\noalign{}
\begin{minipage}[b]{\linewidth}\raggedright
Community
\end{minipage} & \begin{minipage}[b]{\linewidth}\raggedright
Relevance
\end{minipage} \\
\midrule\noalign{}
\endhead
\bottomrule\noalign{}
\endlastfoot
\textbf{Node.js Community} (nodejs.org/community) & Official Node.js
community resources \\
\textbf{JavaScript Weekly} (newsletter) & Weekly digest of JS news;
plain HTML email \\
\textbf{Stack Overflow JavaScript} & Q\&A community; plain HTML;
excellent accessibility \\
\textbf{Dev.to} & Developer blog platform; reasonable accessibility \\
\textbf{GitHub Discussions} on popular JS projects & Community support;
text-based \\
\textbf{NFB Jernigan Institute} & Resources for blind STEM students \\
\textbf{SIGACCESS} & Research on accessible computing including
programming \\
\end{longtable}
}

\begin{center}\rule{0.5\linewidth}{0.5pt}\end{center}

\subsection{19. Quick Reference Sheets}\label{quick-reference-sheets}

\subsubsection{JavaScript vs Python vs Java at a
Glance}\label{javascript-vs-python-vs-java-at-a-glance}

{\def\LTcaptype{none} % do not increment counter
\begin{longtable}[]{@{}
  >{\raggedright\arraybackslash}p{(\linewidth - 6\tabcolsep) * \real{0.1628}}
  >{\raggedright\arraybackslash}p{(\linewidth - 6\tabcolsep) * \real{0.1860}}
  >{\raggedright\arraybackslash}p{(\linewidth - 6\tabcolsep) * \real{0.3488}}
  >{\raggedright\arraybackslash}p{(\linewidth - 6\tabcolsep) * \real{0.3023}}@{}}
\toprule\noalign{}
\begin{minipage}[b]{\linewidth}\raggedright
Concept
\end{minipage} & \begin{minipage}[b]{\linewidth}\raggedright
Python
\end{minipage} & \begin{minipage}[b]{\linewidth}\raggedright
Java
\end{minipage} & \begin{minipage}[b]{\linewidth}\raggedright
JavaScript
\end{minipage} \\
\midrule\noalign{}
\endhead
\bottomrule\noalign{}
\endlastfoot
Print & \texttt{print("x")} & \texttt{System.out.println("x")} &
\texttt{console.log("x")} \\
Variable & \texttt{x\ =\ 5} & \texttt{int\ x\ =\ 5;} &
\texttt{const\ x\ =\ 5;} \\
String format & \texttt{f"Hi\ \{name\}"} &
\texttt{String.format("Hi\ \%s",\ name)} &
\texttt{\textasciigrave{}Hi\ \$\{name\}\textasciigrave{}} \\
User input & \texttt{input("?\ ")} & \texttt{scanner.nextLine()} &
\texttt{await\ askQuestion("?\ ")} \\
Array/List & \texttt{{[}1,\ 2,\ 3{]}} &
\texttt{int{[}{]}\ arr\ =\ \{1,2,3\}} & \texttt{{[}1,\ 2,\ 3{]}} \\
Dictionary/Map & \texttt{\{"a":\ 1\}} &
\texttt{new\ HashMap\textless{}\textgreater{}()} &
\texttt{\{\ a:\ 1\ \}} \\
Function & \texttt{def\ f(x):} & \texttt{public\ int\ f(int\ x)\ \{} &
\texttt{const\ f\ =\ x\ =\textgreater{}\ ...} \\
Class & \texttt{class\ Dog:} & \texttt{public\ class\ Dog\ \{} &
\texttt{class\ Dog\ \{} \\
Exception & \texttt{try/except} & \texttt{try/catch} &
\texttt{try/catch} \\
Null & \texttt{None} & \texttt{null} &
\texttt{null}/\texttt{undefined} \\
Boolean & \texttt{True}/\texttt{False} & \texttt{true}/\texttt{false} &
\texttt{true}/\texttt{false} \\
Length & \texttt{len(x)} & \texttt{x.length} & \texttt{x.length} \\
For each & \texttt{for\ x\ in\ list:} &
\texttt{for\ (T\ x\ :\ list)\ \{} &
\texttt{for\ (const\ x\ of\ list)\ \{} \\
Import & \texttt{import\ math} & \texttt{import\ java.util.*;} &
\texttt{const\ x\ =\ require(\textquotesingle{}./x\textquotesingle{})} \\
\end{longtable}
}

\subsubsection{Accessible Output
Patterns}\label{accessible-output-patterns}

\begin{Shaded}
\begin{Highlighting}[]
\CommentTok{// Standard separator}
\BuiltInTok{console}\OperatorTok{.}\FunctionTok{log}\NormalTok{(}\StringTok{"{-}"}\OperatorTok{.}\FunctionTok{repeat}\NormalTok{(}\DecValTok{40}\NormalTok{))}\OperatorTok{;}

\CommentTok{// Section header}
\BuiltInTok{console}\OperatorTok{.}\FunctionTok{log}\NormalTok{(}\VerbatimStringTok{\textasciigrave{}}\SpecialCharTok{\textbackslash{}n}\VerbatimStringTok{=== }\SpecialCharTok{$\{}\NormalTok{title}\OperatorTok{.}\FunctionTok{toUpperCase}\NormalTok{()}\SpecialCharTok{\}}\VerbatimStringTok{ ===\textasciigrave{}}\NormalTok{)}\OperatorTok{;}

\CommentTok{// Numbered list with details}
\NormalTok{items}\OperatorTok{.}\FunctionTok{forEach}\NormalTok{((item}\OperatorTok{,}\NormalTok{ i) }\KeywordTok{=\textgreater{}}\NormalTok{ \{}
  \BuiltInTok{console}\OperatorTok{.}\FunctionTok{log}\NormalTok{(}\VerbatimStringTok{\textasciigrave{}}\SpecialCharTok{$\{}\NormalTok{i }\OperatorTok{+} \DecValTok{1}\SpecialCharTok{\}}\VerbatimStringTok{. }\SpecialCharTok{$\{}\NormalTok{item}\OperatorTok{.}\AttributeTok{name}\SpecialCharTok{\}}\VerbatimStringTok{\textasciigrave{}}\NormalTok{)}\OperatorTok{;}
  \BuiltInTok{console}\OperatorTok{.}\FunctionTok{log}\NormalTok{(}\VerbatimStringTok{\textasciigrave{}   }\SpecialCharTok{$\{}\NormalTok{item}\OperatorTok{.}\AttributeTok{detail}\SpecialCharTok{\}}\VerbatimStringTok{\textasciigrave{}}\NormalTok{)}\OperatorTok{;}
\NormalTok{\})}\OperatorTok{;}
\BuiltInTok{console}\OperatorTok{.}\FunctionTok{log}\NormalTok{(}\VerbatimStringTok{\textasciigrave{}Total: }\SpecialCharTok{$\{}\NormalTok{items}\OperatorTok{.}\AttributeTok{length}\SpecialCharTok{\}}\VerbatimStringTok{\textasciigrave{}}\NormalTok{)}\OperatorTok{;}

\CommentTok{// Menu}
\KeywordTok{const}\NormalTok{ MENU }\OperatorTok{=} \VerbatimStringTok{\textasciigrave{}}
\VerbatimStringTok{{-}{-}{-} Main Menu {-}{-}{-}}
\VerbatimStringTok{1. Add item}
\VerbatimStringTok{2. View all items}
\VerbatimStringTok{3. Search}
\VerbatimStringTok{4. Edit item}
\VerbatimStringTok{5. Delete item}
\VerbatimStringTok{6. Save}
\VerbatimStringTok{7. Quit}
\VerbatimStringTok{{-}{-}{-}{-}{-}{-}{-}{-}{-}{-}{-}{-}{-}{-}{-}{-}{-}\textasciigrave{}}\OperatorTok{;}

\CommentTok{// Confirmation}
\BuiltInTok{console}\OperatorTok{.}\FunctionTok{log}\NormalTok{(}\VerbatimStringTok{\textasciigrave{}[x] Added: "}\SpecialCharTok{$\{}\NormalTok{name}\SpecialCharTok{\}}\VerbatimStringTok{"\textasciigrave{}}\NormalTok{)}\OperatorTok{;}
\BuiltInTok{console}\OperatorTok{.}\FunctionTok{log}\NormalTok{(}\VerbatimStringTok{\textasciigrave{}[x] Deleted: "}\SpecialCharTok{$\{}\NormalTok{name}\SpecialCharTok{\}}\VerbatimStringTok{"\textasciigrave{}}\NormalTok{)}\OperatorTok{;}
\BuiltInTok{console}\OperatorTok{.}\FunctionTok{log}\NormalTok{(}\VerbatimStringTok{\textasciigrave{}[x] Saved to }\SpecialCharTok{$\{}\NormalTok{filename}\SpecialCharTok{\}}\VerbatimStringTok{\textasciigrave{}}\NormalTok{)}\OperatorTok{;}

\CommentTok{// Error messages}
\BuiltInTok{console}\OperatorTok{.}\FunctionTok{log}\NormalTok{(}\VerbatimStringTok{\textasciigrave{}Error: Name cannot be empty. Please try again.\textasciigrave{}}\NormalTok{)}\OperatorTok{;}
\BuiltInTok{console}\OperatorTok{.}\FunctionTok{log}\NormalTok{(}\VerbatimStringTok{\textasciigrave{}Error: "}\SpecialCharTok{$\{}\NormalTok{filename}\SpecialCharTok{\}}\VerbatimStringTok{" not found. Starting with empty data.\textasciigrave{}}\NormalTok{)}\OperatorTok{;}
\BuiltInTok{console}\OperatorTok{.}\FunctionTok{log}\NormalTok{(}\VerbatimStringTok{\textasciigrave{}Error: Task number must be between 1 and }\SpecialCharTok{$\{}\NormalTok{count}\SpecialCharTok{\}}\VerbatimStringTok{.\textasciigrave{}}\NormalTok{)}\OperatorTok{;}
\end{Highlighting}
\end{Shaded}

\subsubsection{Common JavaScript Error Quick
Reference}\label{common-javascript-error-quick-reference}

{\def\LTcaptype{none} % do not increment counter
\begin{longtable}[]{@{}
  >{\raggedright\arraybackslash}p{(\linewidth - 4\tabcolsep) * \real{0.3680}}
  >{\raggedright\arraybackslash}p{(\linewidth - 4\tabcolsep) * \real{0.2640}}
  >{\raggedright\arraybackslash}p{(\linewidth - 4\tabcolsep) * \real{0.3680}}@{}}
\toprule\noalign{}
\begin{minipage}[b]{\linewidth}\raggedright
Error
\end{minipage} & \begin{minipage}[b]{\linewidth}\raggedright
Common Cause
\end{minipage} & \begin{minipage}[b]{\linewidth}\raggedright
First Check
\end{minipage} \\
\midrule\noalign{}
\endhead
\bottomrule\noalign{}
\endlastfoot
\texttt{ReferenceError:\ x\ is\ not\ defined} & Typo or not declared &
Spelling; scope \\
\texttt{TypeError:\ Cannot\ read\ property\ of\ undefined} & Accessing
property on \texttt{undefined} & Which variable is undefined; optional
chaining \\
\texttt{TypeError:\ x\ is\ not\ a\ function} & Calling non-function &
\texttt{typeof\ x} before calling \\
\texttt{SyntaxError:\ Unexpected\ token} & Syntax mistake & Line number;
check brackets, commas \\
\texttt{SyntaxError:\ Cannot\ use\ import\ statement} & ES module syntax
in CommonJS & Use \texttt{require()} or add
\texttt{"type":\ "module"} \\
\texttt{Error:\ ENOENT:\ no\ such\ file\ or\ directory} & File not found
& \texttt{console.log(\_\_dirname)} and check path \\
\texttt{Error:\ EACCES:\ permission\ denied} & No write permission &
Check file permissions \\
\texttt{RangeError:\ Maximum\ call\ stack\ size\ exceeded} & Infinite
recursion & Check recursive function's base case \\
\texttt{UnhandledPromiseRejectionWarning} & Missing \texttt{await} or
\texttt{.catch()} & Add \texttt{try/catch} around \texttt{await}
calls \\
\texttt{NaN} (not an error, but wrong value) & Arithmetic on non-number
& \texttt{console.log(typeof\ variable)} before math \\
\end{longtable}
}

\subsubsection{Input Helper Module (Copy into every
project)}\label{input-helper-module-copy-into-every-project}

\begin{Shaded}
\begin{Highlighting}[]
\CommentTok{// utils/input.js}
\KeywordTok{const}\NormalTok{ readline }\OperatorTok{=} \PreprocessorTok{require}\NormalTok{(}\StringTok{"readline"}\NormalTok{)}\OperatorTok{;}

\KeywordTok{const}\NormalTok{ rl }\OperatorTok{=}\NormalTok{ readline}\OperatorTok{.}\FunctionTok{createInterface}\NormalTok{(\{}
  \DataTypeTok{input}\OperatorTok{:} \BuiltInTok{process}\OperatorTok{.}\AttributeTok{stdin}\OperatorTok{,}
  \DataTypeTok{output}\OperatorTok{:} \BuiltInTok{process}\OperatorTok{.}\AttributeTok{stdout}\OperatorTok{,}
\NormalTok{\})}\OperatorTok{;}

\CommentTok{/**}
\CommentTok{ * Prompts the user and returns their response as a string.}
\CommentTok{ * }\AnnotationTok{@param}\CommentTok{ }\CommentVarTok{\{string\}}\CommentTok{ prompt {-} The question to display}
\CommentTok{ * }\AnnotationTok{@returns}\CommentTok{ \{Promise}\KeywordTok{\textless{}string\textgreater{}}\CommentTok{\} The user\textquotesingle{}s trimmed response}
\CommentTok{ */}
\KeywordTok{function} \FunctionTok{askQuestion}\NormalTok{(prompt) \{}
  \ControlFlowTok{return} \KeywordTok{new} \BuiltInTok{Promise}\NormalTok{((resolve) }\KeywordTok{=\textgreater{}}\NormalTok{ \{}
\NormalTok{    rl}\OperatorTok{.}\FunctionTok{question}\NormalTok{(prompt}\OperatorTok{,}\NormalTok{ (answer) }\KeywordTok{=\textgreater{}} \FunctionTok{resolve}\NormalTok{(answer))}\OperatorTok{;}
\NormalTok{  \})}\OperatorTok{;}
\NormalTok{\}}

\CommentTok{/**}
\CommentTok{ * Prompts until a valid number is entered.}
\CommentTok{ * }\AnnotationTok{@param}\CommentTok{ }\CommentVarTok{\{string\}}\CommentTok{ prompt {-} The question to display}
\CommentTok{ * }\AnnotationTok{@returns}\CommentTok{ \{Promise}\KeywordTok{\textless{}number\textgreater{}}\CommentTok{\} A valid number}
\CommentTok{ */}
\KeywordTok{async} \KeywordTok{function} \FunctionTok{askNumber}\NormalTok{(prompt) \{}
  \ControlFlowTok{while}\NormalTok{ (}\KeywordTok{true}\NormalTok{) \{}
    \KeywordTok{const}\NormalTok{ input }\OperatorTok{=} \ControlFlowTok{await} \FunctionTok{askQuestion}\NormalTok{(prompt)}\OperatorTok{;}
    \KeywordTok{const}\NormalTok{ num }\OperatorTok{=} \BuiltInTok{Number}\NormalTok{(input}\OperatorTok{.}\FunctionTok{trim}\NormalTok{())}\OperatorTok{;}
    \ControlFlowTok{if}\NormalTok{ (}\OperatorTok{!}\PreprocessorTok{isNaN}\NormalTok{(num)) }\ControlFlowTok{return}\NormalTok{ num}\OperatorTok{;}
    \BuiltInTok{console}\OperatorTok{.}\FunctionTok{log}\NormalTok{(}\StringTok{"Invalid input: please enter a number."}\NormalTok{)}\OperatorTok{;}
\NormalTok{  \}}
\NormalTok{\}}

\CommentTok{/**}
\CommentTok{ * Closes the readline interface. Call at the end of every program.}
\CommentTok{ */}
\KeywordTok{function} \FunctionTok{close}\NormalTok{() \{}
\NormalTok{  rl}\OperatorTok{.}\FunctionTok{close}\NormalTok{()}\OperatorTok{;}
\NormalTok{\}}

\NormalTok{module}\OperatorTok{.}\AttributeTok{exports} \OperatorTok{=}\NormalTok{ \{ askQuestion}\OperatorTok{,}\NormalTok{ askNumber}\OperatorTok{,}\NormalTok{ close \}}\OperatorTok{;}
\end{Highlighting}
\end{Shaded}

\newpage

\section{Unit 0: Setting Up Your JavaScript Development
Environment}\label{unit-0-setting-up-your-javascript-development-environment}

\begin{center}\rule{0.5\linewidth}{0.5pt}\end{center}

\subsection{Lesson 0.1: Installing
Node.js}\label{lesson-0.1-installing-node.js-1}

Node.js lets you run JavaScript on your computer (outside a browser)
using the terminal.

\subsubsection{Installing Node.js}\label{installing-node.js}

Download \textbf{Node.js 18 LTS} from \url{https://nodejs.org/}.

\paragraph{Windows}\label{windows}

\begin{enumerate}
\def\labelenumi{\arabic{enumi}.}
\tightlist
\item
  Download the \texttt{.msi} installer.
\item
  \textbf{CRUCIAL STEP:} Ensure ``Add to PATH'' is checked during
  installation.
\item
  Activate ``Install'' and ``Finish''.
\end{enumerate}

\paragraph{macOS}\label{macos}

\begin{enumerate}
\def\labelenumi{\arabic{enumi}.}
\tightlist
\item
  Download the \texttt{.pkg} installer and follow the prompts.
\item
  Or use Homebrew: \texttt{brew\ install\ node}
\end{enumerate}

\paragraph{Linux (Debian/Ubuntu)}\label{linux-debianubuntu}

\begin{Shaded}
\begin{Highlighting}[]
\ExtensionTok{curl} \AttributeTok{{-}fsSL}\NormalTok{ https://deb.nodesource.com/setup\_18.x }\KeywordTok{|} \FunctionTok{sudo} \AttributeTok{{-}E}\NormalTok{ bash }\AttributeTok{{-}}
\FunctionTok{sudo}\NormalTok{ apt install nodejs}
\end{Highlighting}
\end{Shaded}

\subsubsection{Verify Installation}\label{verify-installation}

\begin{Shaded}
\begin{Highlighting}[]
\ExtensionTok{node} \AttributeTok{{-}{-}version}     \CommentTok{\# Should show v18.x.x or later}
\ExtensionTok{npm} \AttributeTok{{-}{-}version}      \CommentTok{\# Should show npm version}
\end{Highlighting}
\end{Shaded}

\subsubsection{Skills Checklist}\label{skills-checklist}

\begin{itemize}
\tightlist
\item[$\square$]
  Running \texttt{node\ -\/-version} shows the version number.
\item[$\square$]
  Running \texttt{npm\ -\/-version} shows the npm version.
\end{itemize}

\begin{center}\rule{0.5\linewidth}{0.5pt}\end{center}

\subsection{Lesson 0.2: Setting Up VS
Code}\label{lesson-0.2-setting-up-vs-code}

\begin{enumerate}
\def\labelenumi{\arabic{enumi}.}
\tightlist
\item
  Install VS Code from \url{https://code.visualstudio.com/}.
\item
  Press \texttt{Ctrl\ +\ Shift\ +\ X} to open Extensions. Install:

  \begin{itemize}
  \tightlist
  \item
    \textbf{JavaScript (ES6) code snippets} --- Handy code shortcuts.
  \item
    \textbf{ESLint} --- Catches common errors.
  \end{itemize}
\item
  Enable Screen Reader Mode: \texttt{Ctrl\ +\ Shift\ +\ P}
  -\textgreater{} ``Toggle Screen Reader Optimized''.
\end{enumerate}

\subsubsection{VS Code Shortcuts for
JavaScript}\label{vs-code-shortcuts-for-javascript}

{\def\LTcaptype{none} % do not increment counter
\begin{longtable}[]{@{}
  >{\raggedright\arraybackslash}p{(\linewidth - 2\tabcolsep) * \real{0.3966}}
  >{\raggedright\arraybackslash}p{(\linewidth - 2\tabcolsep) * \real{0.6034}}@{}}
\toprule\noalign{}
\begin{minipage}[b]{\linewidth}\raggedright
Action
\end{minipage} & \begin{minipage}[b]{\linewidth}\raggedright
Shortcut
\end{minipage} \\
\midrule\noalign{}
\endhead
\bottomrule\noalign{}
\endlastfoot
Run file (via terminal) &
\texttt{Ctrl\ +\ \textbackslash{}\textasciigrave{}\textasciigrave{}\ then}node
filename.js\texttt{\textbar{}\ \textbar{}\ Command\ Palette\ \ \ \ \ \ \ \ \ \textbar{}}Ctrl
+ Shift +
P\texttt{\textbar{}\ \textbar{}\ Problems\ View\ \ \ \ \ \ \ \ \ \ \ \textbar{}}Ctrl
+ Shift +
M\texttt{\textbar{}\ \textbar{}\ Next\ Error\ \ \ \ \ \ \ \ \ \ \ \ \ \ \textbar{}}F8\texttt{\textbar{}\ \textbar{}\ Format\ Code\ \ \ \ \ \ \ \ \ \ \ \ \ \textbar{}}Shift
+ Alt +
F\texttt{\textbar{}\ \textbar{}\ Go\ to\ Definition\ \ \ \ \ \ \ \ \textbar{}}F12` \\
\end{longtable}
}

\begin{center}\rule{0.5\linewidth}{0.5pt}\end{center}

\subsection{Lesson 0.3: Your First JavaScript
File}\label{lesson-0.3-your-first-javascript-file}

\subsubsection{Hello World}\label{hello-world}

Create \texttt{hello.js}:

\begin{Shaded}
\begin{Highlighting}[]
\BuiltInTok{console}\OperatorTok{.}\FunctionTok{log}\NormalTok{(}\StringTok{"Hello, World!"}\NormalTok{)}\OperatorTok{;}
\BuiltInTok{console}\OperatorTok{.}\FunctionTok{log}\NormalTok{(}\StringTok{"Welcome to JavaScript!"}\NormalTok{)}\OperatorTok{;}
\end{Highlighting}
\end{Shaded}

Run it:

\begin{Shaded}
\begin{Highlighting}[]
\ExtensionTok{node}\NormalTok{ hello.js}
\end{Highlighting}
\end{Shaded}

\subsubsection{The Node.js REPL}\label{the-node.js-repl}

For quick experiments, open an interactive session:

\begin{Shaded}
\begin{Highlighting}[]
\ExtensionTok{node}
\end{Highlighting}
\end{Shaded}

Type JavaScript directly and press Enter to execute. Type \texttt{.exit}
to quit.

\subsubsection{Key Differences from
Python}\label{key-differences-from-python}

{\def\LTcaptype{none} % do not increment counter
\begin{longtable}[]{@{}
  >{\raggedright\arraybackslash}p{(\linewidth - 2\tabcolsep) * \real{0.3714}}
  >{\raggedright\arraybackslash}p{(\linewidth - 2\tabcolsep) * \real{0.6286}}@{}}
\toprule\noalign{}
\begin{minipage}[b]{\linewidth}\raggedright
Python
\end{minipage} & \begin{minipage}[b]{\linewidth}\raggedright
JavaScript
\end{minipage} \\
\midrule\noalign{}
\endhead
\bottomrule\noalign{}
\endlastfoot
\texttt{print()} & \texttt{console.log()} \\
\texttt{\#} for comments & \texttt{//} for single-line, \texttt{/*\ */}
for multi-line \\
No semicolons required & Semicolons optional but recommended \\
Indentation defines blocks & Curly braces \texttt{\{\}} define blocks \\
\texttt{def} for functions & \texttt{function} or
\texttt{=\textgreater{}} for functions \\
\texttt{True}/\texttt{False} & \texttt{true}/\texttt{false}
(lowercase) \\
\texttt{None} & \texttt{null} and \texttt{undefined} \\
\texttt{len(x)} & \texttt{x.length} \\
\end{longtable}
}

\subsubsection{Skills Checklist}\label{skills-checklist-1}

\begin{itemize}
\tightlist
\item[$\square$]
  I can create a \texttt{.js} file and run it with
  \texttt{node\ filename.js}.
\item[$\square$]
  I can use the Node.js REPL for quick experiments.
\end{itemize}

\begin{center}\rule{0.5\linewidth}{0.5pt}\end{center}

\section{Unit 1: Foundations --- Getting Started with
JavaScript}\label{unit-1-foundations-getting-started-with-javascript}

\subsection{Overview}\label{overview}

This unit covers JavaScript fundamentals: variables, data types, and
basic output.

\begin{center}\rule{0.5\linewidth}{0.5pt}\end{center}

\subsection{Lesson 1.1: Welcome to Computer
Science}\label{lesson-1.1-welcome-to-computer-science}

\emph{(See Python Unit 1.1 for conceptual content --- hardware,
software, algorithms, and problem-solving are language-agnostic.)}

\subsubsection{JavaScript's Role}\label{javascripts-role}

JavaScript was created by Brendan Eich at Netscape in 1995. Originally
just for web browsers, it now runs everywhere thanks to Node.js (2009).

\begin{center}\rule{0.5\linewidth}{0.5pt}\end{center}

\subsection{Lesson 1.2: Variables and
Output}\label{lesson-1.2-variables-and-output}

\subsubsection{\texorpdfstring{\texttt{console.log()}}{console.log()}}\label{console.log}

\begin{Shaded}
\begin{Highlighting}[]
\BuiltInTok{console}\OperatorTok{.}\FunctionTok{log}\NormalTok{(}\StringTok{"Hello, World!"}\NormalTok{)}\OperatorTok{;}
\BuiltInTok{console}\OperatorTok{.}\FunctionTok{log}\NormalTok{(}\DecValTok{42}\NormalTok{)}\OperatorTok{;}
\BuiltInTok{console}\OperatorTok{.}\FunctionTok{log}\NormalTok{(}\KeywordTok{true}\NormalTok{)}\OperatorTok{;}
\BuiltInTok{console}\OperatorTok{.}\FunctionTok{log}\NormalTok{(}\StringTok{"Value:"}\OperatorTok{,} \DecValTok{42}\OperatorTok{,} \StringTok{"and"}\OperatorTok{,} \StringTok{"text"}\NormalTok{)}\OperatorTok{;}
\end{Highlighting}
\end{Shaded}

\subsubsection{Declaring Variables}\label{declaring-variables}

JavaScript has three ways to declare variables:

{\def\LTcaptype{none} % do not increment counter
\begin{longtable}[]{@{}
  >{\raggedright\arraybackslash}p{(\linewidth - 6\tabcolsep) * \real{0.1129}}
  >{\raggedright\arraybackslash}p{(\linewidth - 6\tabcolsep) * \real{0.1290}}
  >{\raggedright\arraybackslash}p{(\linewidth - 6\tabcolsep) * \real{0.1935}}
  >{\raggedright\arraybackslash}p{(\linewidth - 6\tabcolsep) * \real{0.5645}}@{}}
\toprule\noalign{}
\begin{minipage}[b]{\linewidth}\raggedright
Keyword
\end{minipage} & \begin{minipage}[b]{\linewidth}\raggedright
Scope
\end{minipage} & \begin{minipage}[b]{\linewidth}\raggedright
Reassignable
\end{minipage} & \begin{minipage}[b]{\linewidth}\raggedright
Use When
\end{minipage} \\
\midrule\noalign{}
\endhead
\bottomrule\noalign{}
\endlastfoot
\texttt{const} & Block & No & Value won't change (default choice) \\
\texttt{let} & Block & Yes & Value will change \\
\texttt{var} & Function & Yes & Legacy --- avoid in modern code \\
\end{longtable}
}

\begin{Shaded}
\begin{Highlighting}[]
\KeywordTok{const} \ConstantTok{PI} \OperatorTok{=} \FloatTok{3.14159}\OperatorTok{;} \CommentTok{// Cannot be reassigned}
\KeywordTok{let}\NormalTok{ score }\OperatorTok{=} \DecValTok{100}\OperatorTok{;} \CommentTok{// Can change}
\NormalTok{score }\OperatorTok{=} \DecValTok{150}\OperatorTok{;} \CommentTok{// OK}

\CommentTok{// var is function{-}scoped (legacy)}
\KeywordTok{var}\NormalTok{ oldWay }\OperatorTok{=} \StringTok{"avoid this"}\OperatorTok{;}
\end{Highlighting}
\end{Shaded}

\textbf{Best practice:} Use \texttt{const} by default; use \texttt{let}
only when you need to reassign.

\subsubsection{Template Literals}\label{template-literals}

\begin{Shaded}
\begin{Highlighting}[]
\KeywordTok{const}\NormalTok{ name }\OperatorTok{=} \StringTok{"Alice"}\OperatorTok{;}
\KeywordTok{const}\NormalTok{ age }\OperatorTok{=} \DecValTok{15}\OperatorTok{;}

\CommentTok{// Template literal (backticks {-}{-}{-} preferred)}
\BuiltInTok{console}\OperatorTok{.}\FunctionTok{log}\NormalTok{(}\VerbatimStringTok{\textasciigrave{}Hello, }\SpecialCharTok{$\{}\NormalTok{name}\SpecialCharTok{\}}\VerbatimStringTok{! You are }\SpecialCharTok{$\{}\NormalTok{age}\SpecialCharTok{\}}\VerbatimStringTok{ years old.\textasciigrave{}}\NormalTok{)}\OperatorTok{;}

\CommentTok{// Concatenation (older style)}
\BuiltInTok{console}\OperatorTok{.}\FunctionTok{log}\NormalTok{(}\StringTok{"Hello, "} \OperatorTok{+}\NormalTok{ name }\OperatorTok{+} \StringTok{"! You are "} \OperatorTok{+}\NormalTok{ age }\OperatorTok{+} \StringTok{" years old."}\NormalTok{)}\OperatorTok{;}
\end{Highlighting}
\end{Shaded}

\subsubsection{Skills Checklist}\label{skills-checklist-2}

\begin{itemize}
\tightlist
\item[$\square$]
  I can use \texttt{console.log()} to display output.
\item[$\square$]
  I can declare variables using \texttt{const} and \texttt{let}.
\item[$\square$]
  I can use template literals with \texttt{\$\{\}}.
\end{itemize}

\begin{center}\rule{0.5\linewidth}{0.5pt}\end{center}

\subsection{Lesson 1.3: Data Types}\label{lesson-1.3-data-types}

\subsubsection{Primitive Types}\label{primitive-types}

\begin{Shaded}
\begin{Highlighting}[]
\CommentTok{// Number (integers and floats use the same type)}
\KeywordTok{let}\NormalTok{ integer }\OperatorTok{=} \DecValTok{42}\OperatorTok{;}
\KeywordTok{let}\NormalTok{ decimal }\OperatorTok{=} \FloatTok{3.14}\OperatorTok{;}
\KeywordTok{let}\NormalTok{ negative }\OperatorTok{=} \OperatorTok{{-}}\DecValTok{10}\OperatorTok{;}

\CommentTok{// String}
\KeywordTok{let}\NormalTok{ name }\OperatorTok{=} \StringTok{"Alice"}\OperatorTok{;}
\KeywordTok{let}\NormalTok{ greeting }\OperatorTok{=} \StringTok{"Hello"}\OperatorTok{;}
\KeywordTok{let}\NormalTok{ multiLine }\OperatorTok{=} \VerbatimStringTok{\textasciigrave{}Line 1}
\VerbatimStringTok{Line 2\textasciigrave{}}\OperatorTok{;}

\CommentTok{// Boolean}
\KeywordTok{let}\NormalTok{ isRaining }\OperatorTok{=} \KeywordTok{false}\OperatorTok{;}
\KeywordTok{let}\NormalTok{ isLoggedIn }\OperatorTok{=} \KeywordTok{true}\OperatorTok{;}

\CommentTok{// null (intentional absence of value)}
\KeywordTok{let}\NormalTok{ emptyValue }\OperatorTok{=} \KeywordTok{null}\OperatorTok{;}

\CommentTok{// undefined (value not assigned)}
\KeywordTok{let}\NormalTok{ notAssigned}\OperatorTok{;}
\BuiltInTok{console}\OperatorTok{.}\FunctionTok{log}\NormalTok{(notAssigned)}\OperatorTok{;} \CommentTok{// undefined}

\CommentTok{// Check type}
\BuiltInTok{console}\OperatorTok{.}\FunctionTok{log}\NormalTok{(}\KeywordTok{typeof} \DecValTok{42}\NormalTok{)}\OperatorTok{;} \CommentTok{// "number"}
\BuiltInTok{console}\OperatorTok{.}\FunctionTok{log}\NormalTok{(}\KeywordTok{typeof} \StringTok{"hello"}\NormalTok{)}\OperatorTok{;} \CommentTok{// "string"}
\BuiltInTok{console}\OperatorTok{.}\FunctionTok{log}\NormalTok{(}\KeywordTok{typeof} \KeywordTok{true}\NormalTok{)}\OperatorTok{;} \CommentTok{// "boolean"}
\BuiltInTok{console}\OperatorTok{.}\FunctionTok{log}\NormalTok{(}\KeywordTok{typeof} \KeywordTok{null}\NormalTok{)}\OperatorTok{;} \CommentTok{// "object" (known quirk!)}
\BuiltInTok{console}\OperatorTok{.}\FunctionTok{log}\NormalTok{(}\KeywordTok{typeof} \KeywordTok{undefined}\NormalTok{)}\OperatorTok{;} \CommentTok{// "undefined"}
\end{Highlighting}
\end{Shaded}

\subsubsection{Arithmetic Operators}\label{arithmetic-operators}

\begin{Shaded}
\begin{Highlighting}[]
\BuiltInTok{console}\OperatorTok{.}\FunctionTok{log}\NormalTok{(}\DecValTok{5} \OperatorTok{+} \DecValTok{3}\NormalTok{)}\OperatorTok{;} \CommentTok{// 8}
\BuiltInTok{console}\OperatorTok{.}\FunctionTok{log}\NormalTok{(}\DecValTok{10} \OperatorTok{{-}} \DecValTok{4}\NormalTok{)}\OperatorTok{;} \CommentTok{// 6}
\BuiltInTok{console}\OperatorTok{.}\FunctionTok{log}\NormalTok{(}\DecValTok{7} \OperatorTok{*} \DecValTok{6}\NormalTok{)}\OperatorTok{;} \CommentTok{// 42}
\BuiltInTok{console}\OperatorTok{.}\FunctionTok{log}\NormalTok{(}\DecValTok{15} \OperatorTok{/} \DecValTok{4}\NormalTok{)}\OperatorTok{;} \CommentTok{// 3.75  (always decimal, unlike Python 2 or Java)}
\BuiltInTok{console}\OperatorTok{.}\FunctionTok{log}\NormalTok{(}\DecValTok{10} \OperatorTok{\%} \DecValTok{3}\NormalTok{)}\OperatorTok{;} \CommentTok{// 1  (remainder)}
\BuiltInTok{console}\OperatorTok{.}\FunctionTok{log}\NormalTok{(}\DecValTok{2} \OperatorTok{**} \DecValTok{10}\NormalTok{)}\OperatorTok{;} \CommentTok{// 1024  (exponentiation)}

\CommentTok{// Math.floor for integer division equivalent}
\BuiltInTok{console}\OperatorTok{.}\FunctionTok{log}\NormalTok{(}\BuiltInTok{Math}\OperatorTok{.}\FunctionTok{floor}\NormalTok{(}\DecValTok{15} \OperatorTok{/} \DecValTok{4}\NormalTok{))}\OperatorTok{;} \CommentTok{// 3}
\end{Highlighting}
\end{Shaded}

\subsubsection{Type Conversion}\label{type-conversion}

\begin{Shaded}
\begin{Highlighting}[]
\CommentTok{// String to Number}
\KeywordTok{let}\NormalTok{ str }\OperatorTok{=} \StringTok{"123"}\OperatorTok{;}
\KeywordTok{let}\NormalTok{ num }\OperatorTok{=} \BuiltInTok{Number}\NormalTok{(str)}\OperatorTok{;} \CommentTok{// 123}
\KeywordTok{let}\NormalTok{ numParsed }\OperatorTok{=} \PreprocessorTok{parseInt}\NormalTok{(str)}\OperatorTok{;} \CommentTok{// 123}
\KeywordTok{let}\NormalTok{ floatParsed }\OperatorTok{=} \PreprocessorTok{parseFloat}\NormalTok{(}\StringTok{"3.14"}\NormalTok{)}\OperatorTok{;} \CommentTok{// 3.14}

\CommentTok{// Number to String}
\KeywordTok{let}\NormalTok{ n }\OperatorTok{=} \DecValTok{42}\OperatorTok{;}
\KeywordTok{let}\NormalTok{ s }\OperatorTok{=} \BuiltInTok{String}\NormalTok{(n)}\OperatorTok{;} \CommentTok{// "42"}
\KeywordTok{let}\NormalTok{ s2 }\OperatorTok{=}\NormalTok{ n}\OperatorTok{.}\FunctionTok{toString}\NormalTok{()}\OperatorTok{;} \CommentTok{// "42"}
\KeywordTok{let}\NormalTok{ s3 }\OperatorTok{=} \VerbatimStringTok{\textasciigrave{}}\SpecialCharTok{$\{}\NormalTok{n}\SpecialCharTok{\}}\VerbatimStringTok{\textasciigrave{}}\OperatorTok{;} \CommentTok{// "42" (template literal)}

\CommentTok{// Checking}
\BuiltInTok{console}\OperatorTok{.}\FunctionTok{log}\NormalTok{(}\PreprocessorTok{isNaN}\NormalTok{(}\StringTok{"abc"}\NormalTok{))}\OperatorTok{;} \CommentTok{// true {-}{-}{-} "not a number"}
\BuiltInTok{console}\OperatorTok{.}\FunctionTok{log}\NormalTok{(}\PreprocessorTok{isNaN}\NormalTok{(}\DecValTok{42}\NormalTok{))}\OperatorTok{;} \CommentTok{// false}
\BuiltInTok{console}\OperatorTok{.}\FunctionTok{log}\NormalTok{(}\BuiltInTok{Number}\NormalTok{(}\StringTok{"abc"}\NormalTok{))}\OperatorTok{;} \CommentTok{// NaN}
\end{Highlighting}
\end{Shaded}

\subsubsection{Constants}\label{constants}

\begin{Shaded}
\begin{Highlighting}[]
\KeywordTok{const}\NormalTok{ MAX\_SIZE }\OperatorTok{=} \DecValTok{100}\OperatorTok{;}
\KeywordTok{const}\NormalTok{ APP\_NAME }\OperatorTok{=} \StringTok{"My App"}\OperatorTok{;}
\KeywordTok{const} \ConstantTok{PI} \OperatorTok{=} \BuiltInTok{Math}\OperatorTok{.}\ConstantTok{PI}\OperatorTok{;} \CommentTok{// Use Math.PI {-}{-}{-} already defined}
\end{Highlighting}
\end{Shaded}

\subsubsection{Examples}\label{examples}

\textbf{Example 1.3.1: Personal Information}

\begin{Shaded}
\begin{Highlighting}[]
\KeywordTok{const}\NormalTok{ myName }\OperatorTok{=} \StringTok{"Jamie"}\OperatorTok{;}
\KeywordTok{const}\NormalTok{ favoriteNumber }\OperatorTok{=} \DecValTok{7}\OperatorTok{;}
\KeywordTok{const}\NormalTok{ heightInMeters }\OperatorTok{=} \FloatTok{1.75}\OperatorTok{;}
\KeywordTok{const}\NormalTok{ likesPizza }\OperatorTok{=} \KeywordTok{true}\OperatorTok{;}

\BuiltInTok{console}\OperatorTok{.}\FunctionTok{log}\NormalTok{(}\VerbatimStringTok{\textasciigrave{}Name: }\SpecialCharTok{$\{}\NormalTok{myName}\SpecialCharTok{\}}\VerbatimStringTok{\textasciigrave{}}\NormalTok{)}\OperatorTok{;}
\BuiltInTok{console}\OperatorTok{.}\FunctionTok{log}\NormalTok{(}\VerbatimStringTok{\textasciigrave{}Favorite number: }\SpecialCharTok{$\{}\NormalTok{favoriteNumber}\SpecialCharTok{\}}\VerbatimStringTok{\textasciigrave{}}\NormalTok{)}\OperatorTok{;}
\BuiltInTok{console}\OperatorTok{.}\FunctionTok{log}\NormalTok{(}\VerbatimStringTok{\textasciigrave{}Height: }\SpecialCharTok{$\{}\NormalTok{heightInMeters}\SpecialCharTok{\}}\VerbatimStringTok{ meters\textasciigrave{}}\NormalTok{)}\OperatorTok{;}
\BuiltInTok{console}\OperatorTok{.}\FunctionTok{log}\NormalTok{(}\VerbatimStringTok{\textasciigrave{}Likes pizza: }\SpecialCharTok{$\{}\NormalTok{likesPizza}\SpecialCharTok{\}}\VerbatimStringTok{\textasciigrave{}}\NormalTok{)}\OperatorTok{;}
\end{Highlighting}
\end{Shaded}

\textbf{Example 1.3.2: Simple Math}

\begin{Shaded}
\begin{Highlighting}[]
\KeywordTok{const}\NormalTok{ num1 }\OperatorTok{=} \DecValTok{15}\OperatorTok{;}
\KeywordTok{const}\NormalTok{ num2 }\OperatorTok{=} \DecValTok{7}\OperatorTok{;}

\BuiltInTok{console}\OperatorTok{.}\FunctionTok{log}\NormalTok{(}\VerbatimStringTok{\textasciigrave{}Sum: }\SpecialCharTok{$\{}\NormalTok{num1 }\OperatorTok{+}\NormalTok{ num2}\SpecialCharTok{\}}\VerbatimStringTok{\textasciigrave{}}\NormalTok{)}\OperatorTok{;}
\BuiltInTok{console}\OperatorTok{.}\FunctionTok{log}\NormalTok{(}\VerbatimStringTok{\textasciigrave{}Difference: }\SpecialCharTok{$\{}\NormalTok{num1 }\OperatorTok{{-}}\NormalTok{ num2}\SpecialCharTok{\}}\VerbatimStringTok{\textasciigrave{}}\NormalTok{)}\OperatorTok{;}
\BuiltInTok{console}\OperatorTok{.}\FunctionTok{log}\NormalTok{(}\VerbatimStringTok{\textasciigrave{}Product: }\SpecialCharTok{$\{}\NormalTok{num1 }\OperatorTok{*}\NormalTok{ num2}\SpecialCharTok{\}}\VerbatimStringTok{\textasciigrave{}}\NormalTok{)}\OperatorTok{;}
\BuiltInTok{console}\OperatorTok{.}\FunctionTok{log}\NormalTok{(}\VerbatimStringTok{\textasciigrave{}Quotient: }\SpecialCharTok{$\{}\NormalTok{(num1 }\OperatorTok{/}\NormalTok{ num2)}\OperatorTok{.}\FunctionTok{toFixed}\NormalTok{(}\DecValTok{2}\NormalTok{)}\SpecialCharTok{\}}\VerbatimStringTok{\textasciigrave{}}\NormalTok{)}\OperatorTok{;}
\BuiltInTok{console}\OperatorTok{.}\FunctionTok{log}\NormalTok{(}\VerbatimStringTok{\textasciigrave{}Remainder: }\SpecialCharTok{$\{}\NormalTok{num1 }\OperatorTok{\%}\NormalTok{ num2}\SpecialCharTok{\}}\VerbatimStringTok{\textasciigrave{}}\NormalTok{)}\OperatorTok{;}
\end{Highlighting}
\end{Shaded}

\subsubsection{Activities}\label{activities}

\textbf{Activity 1.3.1: Area Calculator}

\begin{Shaded}
\begin{Highlighting}[]
\KeywordTok{const}\NormalTok{ length }\OperatorTok{=} \DecValTok{12}\OperatorTok{;}
\KeywordTok{const}\NormalTok{ width }\OperatorTok{=} \DecValTok{5}\OperatorTok{;}
\KeywordTok{const}\NormalTok{ height }\OperatorTok{=} \DecValTok{8}\OperatorTok{;}

\KeywordTok{const}\NormalTok{ rectangleArea }\OperatorTok{=}\NormalTok{ length }\OperatorTok{*}\NormalTok{ width}\OperatorTok{;}
\KeywordTok{const}\NormalTok{ boxVolume }\OperatorTok{=}\NormalTok{ length }\OperatorTok{*}\NormalTok{ width }\OperatorTok{*}\NormalTok{ height}\OperatorTok{;}
\KeywordTok{const}\NormalTok{ circleArea }\OperatorTok{=} \BuiltInTok{Math}\OperatorTok{.}\ConstantTok{PI} \OperatorTok{*}\NormalTok{ (length }\OperatorTok{/} \DecValTok{2}\NormalTok{) }\OperatorTok{**} \DecValTok{2}\OperatorTok{;}

\BuiltInTok{console}\OperatorTok{.}\FunctionTok{log}\NormalTok{(}\VerbatimStringTok{\textasciigrave{}Rectangle area: }\SpecialCharTok{$\{}\NormalTok{rectangleArea}\SpecialCharTok{\}}\VerbatimStringTok{\textasciigrave{}}\NormalTok{)}\OperatorTok{;}
\BuiltInTok{console}\OperatorTok{.}\FunctionTok{log}\NormalTok{(}\VerbatimStringTok{\textasciigrave{}Box volume: }\SpecialCharTok{$\{}\NormalTok{boxVolume}\SpecialCharTok{\}}\VerbatimStringTok{\textasciigrave{}}\NormalTok{)}\OperatorTok{;}
\BuiltInTok{console}\OperatorTok{.}\FunctionTok{log}\NormalTok{(}\VerbatimStringTok{\textasciigrave{}Circle area: }\SpecialCharTok{$\{}\NormalTok{circleArea}\OperatorTok{.}\FunctionTok{toFixed}\NormalTok{(}\DecValTok{2}\NormalTok{)}\SpecialCharTok{\}}\VerbatimStringTok{\textasciigrave{}}\NormalTok{)}\OperatorTok{;}
\end{Highlighting}
\end{Shaded}

\subsubsection{Skills Checklist}\label{skills-checklist-3}

\begin{itemize}
\tightlist
\item[$\square$]
  I can identify JavaScript's primitive data types.
\item[$\square$]
  I can perform arithmetic operations in JavaScript.
\item[$\square$]
  I understand that JavaScript's \texttt{/} always gives a decimal
  result.
\item[$\square$]
  I can convert between types using \texttt{Number()},
  \texttt{String()}, \texttt{parseInt()}.
\item[$\square$]
  I understand \texttt{const} vs.~\texttt{let}.
\end{itemize}

\newpage

\section{Unit 2: Numbers, Input, and the Math
Object}\label{unit-2-numbers-input-and-the-math-object}

\begin{center}\rule{0.5\linewidth}{0.5pt}\end{center}

\subsection{Lesson 2.1: The Math
Object}\label{lesson-2.1-the-math-object}

JavaScript's built-in \texttt{Math} object provides mathematical
functions and constants:

\begin{Shaded}
\begin{Highlighting}[]
\CommentTok{// Constants}
\BuiltInTok{console}\OperatorTok{.}\FunctionTok{log}\NormalTok{(}\BuiltInTok{Math}\OperatorTok{.}\ConstantTok{PI}\NormalTok{)}\OperatorTok{;} \CommentTok{// 3.14159265358979...}
\BuiltInTok{console}\OperatorTok{.}\FunctionTok{log}\NormalTok{(}\BuiltInTok{Math}\OperatorTok{.}\ConstantTok{E}\NormalTok{)}\OperatorTok{;} \CommentTok{// 2.71828...}

\CommentTok{// Functions}
\BuiltInTok{console}\OperatorTok{.}\FunctionTok{log}\NormalTok{(}\BuiltInTok{Math}\OperatorTok{.}\FunctionTok{sqrt}\NormalTok{(}\DecValTok{25}\NormalTok{))}\OperatorTok{;} \CommentTok{// 5}
\BuiltInTok{console}\OperatorTok{.}\FunctionTok{log}\NormalTok{(}\BuiltInTok{Math}\OperatorTok{.}\FunctionTok{pow}\NormalTok{(}\DecValTok{2}\OperatorTok{,} \DecValTok{10}\NormalTok{))}\OperatorTok{;} \CommentTok{// 1024}
\BuiltInTok{console}\OperatorTok{.}\FunctionTok{log}\NormalTok{(}\BuiltInTok{Math}\OperatorTok{.}\FunctionTok{abs}\NormalTok{(}\OperatorTok{{-}}\FloatTok{7.5}\NormalTok{))}\OperatorTok{;} \CommentTok{// 7.5}
\BuiltInTok{console}\OperatorTok{.}\FunctionTok{log}\NormalTok{(}\BuiltInTok{Math}\OperatorTok{.}\FunctionTok{ceil}\NormalTok{(}\FloatTok{4.2}\NormalTok{))}\OperatorTok{;} \CommentTok{// 5}
\BuiltInTok{console}\OperatorTok{.}\FunctionTok{log}\NormalTok{(}\BuiltInTok{Math}\OperatorTok{.}\FunctionTok{floor}\NormalTok{(}\FloatTok{4.8}\NormalTok{))}\OperatorTok{;} \CommentTok{// 4}
\BuiltInTok{console}\OperatorTok{.}\FunctionTok{log}\NormalTok{(}\BuiltInTok{Math}\OperatorTok{.}\FunctionTok{round}\NormalTok{(}\FloatTok{4.5}\NormalTok{))}\OperatorTok{;} \CommentTok{// 5}
\BuiltInTok{console}\OperatorTok{.}\FunctionTok{log}\NormalTok{(}\BuiltInTok{Math}\OperatorTok{.}\FunctionTok{round}\NormalTok{(}\FloatTok{4.4}\NormalTok{))}\OperatorTok{;} \CommentTok{// 4}
\BuiltInTok{console}\OperatorTok{.}\FunctionTok{log}\NormalTok{(}\BuiltInTok{Math}\OperatorTok{.}\FunctionTok{max}\NormalTok{(}\DecValTok{3}\OperatorTok{,} \DecValTok{7}\OperatorTok{,} \DecValTok{1}\OperatorTok{,} \DecValTok{9}\OperatorTok{,} \DecValTok{5}\NormalTok{))}\OperatorTok{;} \CommentTok{// 9}
\BuiltInTok{console}\OperatorTok{.}\FunctionTok{log}\NormalTok{(}\BuiltInTok{Math}\OperatorTok{.}\FunctionTok{min}\NormalTok{(}\DecValTok{3}\OperatorTok{,} \DecValTok{7}\OperatorTok{,} \DecValTok{1}\OperatorTok{,} \DecValTok{9}\OperatorTok{,} \DecValTok{5}\NormalTok{))}\OperatorTok{;} \CommentTok{// 1}

\CommentTok{// Random number between 0 (inclusive) and 1 (exclusive)}
\BuiltInTok{console}\OperatorTok{.}\FunctionTok{log}\NormalTok{(}\BuiltInTok{Math}\OperatorTok{.}\FunctionTok{random}\NormalTok{())}\OperatorTok{;}

\CommentTok{// Random integer between min and max (inclusive)}
\KeywordTok{function} \FunctionTok{randomInt}\NormalTok{(min}\OperatorTok{,}\NormalTok{ max) \{}
  \ControlFlowTok{return} \BuiltInTok{Math}\OperatorTok{.}\FunctionTok{floor}\NormalTok{(}\BuiltInTok{Math}\OperatorTok{.}\FunctionTok{random}\NormalTok{() }\OperatorTok{*}\NormalTok{ (max }\OperatorTok{{-}}\NormalTok{ min }\OperatorTok{+} \DecValTok{1}\NormalTok{)) }\OperatorTok{+}\NormalTok{ min}\OperatorTok{;}
\NormalTok{\}}
\BuiltInTok{console}\OperatorTok{.}\FunctionTok{log}\NormalTok{(}\FunctionTok{randomInt}\NormalTok{(}\DecValTok{1}\OperatorTok{,} \DecValTok{6}\NormalTok{))}\OperatorTok{;} \CommentTok{// Simulates a dice roll}
\end{Highlighting}
\end{Shaded}

\begin{center}\rule{0.5\linewidth}{0.5pt}\end{center}

\subsection{Lesson 2.2: Getting User Input in
Node.js}\label{lesson-2.2-getting-user-input-in-node.js}

Getting user input in Node.js requires the built-in \texttt{readline}
module.

\subsubsection{\texorpdfstring{Using \texttt{readline} (Standard
Library)}{Using readline (Standard Library)}}\label{using-readline-standard-library}

\begin{Shaded}
\begin{Highlighting}[]
\KeywordTok{const}\NormalTok{ readline }\OperatorTok{=} \PreprocessorTok{require}\NormalTok{(}\StringTok{"readline"}\NormalTok{)}\OperatorTok{;}

\KeywordTok{const}\NormalTok{ rl }\OperatorTok{=}\NormalTok{ readline}\OperatorTok{.}\FunctionTok{createInterface}\NormalTok{(\{}
  \DataTypeTok{input}\OperatorTok{:} \BuiltInTok{process}\OperatorTok{.}\AttributeTok{stdin}\OperatorTok{,}
  \DataTypeTok{output}\OperatorTok{:} \BuiltInTok{process}\OperatorTok{.}\AttributeTok{stdout}\OperatorTok{,}
\NormalTok{\})}\OperatorTok{;}

\NormalTok{rl}\OperatorTok{.}\FunctionTok{question}\NormalTok{(}\StringTok{"What is your name? "}\OperatorTok{,}\NormalTok{ (name) }\KeywordTok{=\textgreater{}}\NormalTok{ \{}
  \BuiltInTok{console}\OperatorTok{.}\FunctionTok{log}\NormalTok{(}\VerbatimStringTok{\textasciigrave{}Hello, }\SpecialCharTok{$\{}\NormalTok{name}\SpecialCharTok{\}}\VerbatimStringTok{!\textasciigrave{}}\NormalTok{)}\OperatorTok{;}
\NormalTok{  rl}\OperatorTok{.}\FunctionTok{close}\NormalTok{()}\OperatorTok{;}
\NormalTok{\})}\OperatorTok{;}
\end{Highlighting}
\end{Shaded}

\subsubsection{Promisified Input (Modern
Approach)}\label{promisified-input-modern-approach}

\begin{Shaded}
\begin{Highlighting}[]
\KeywordTok{const}\NormalTok{ readline }\OperatorTok{=} \PreprocessorTok{require}\NormalTok{(}\StringTok{"readline"}\NormalTok{)}\OperatorTok{;}

\KeywordTok{const}\NormalTok{ rl }\OperatorTok{=}\NormalTok{ readline}\OperatorTok{.}\FunctionTok{createInterface}\NormalTok{(\{}
  \DataTypeTok{input}\OperatorTok{:} \BuiltInTok{process}\OperatorTok{.}\AttributeTok{stdin}\OperatorTok{,}
  \DataTypeTok{output}\OperatorTok{:} \BuiltInTok{process}\OperatorTok{.}\AttributeTok{stdout}\OperatorTok{,}
\NormalTok{\})}\OperatorTok{;}

\KeywordTok{function} \FunctionTok{askQuestion}\NormalTok{(prompt) \{}
  \ControlFlowTok{return} \KeywordTok{new} \BuiltInTok{Promise}\NormalTok{((resolve) }\KeywordTok{=\textgreater{}}\NormalTok{ \{}
\NormalTok{    rl}\OperatorTok{.}\FunctionTok{question}\NormalTok{(prompt}\OperatorTok{,}\NormalTok{ (answer) }\KeywordTok{=\textgreater{}}\NormalTok{ \{}
      \FunctionTok{resolve}\NormalTok{(answer)}\OperatorTok{;}
\NormalTok{    \})}\OperatorTok{;}
\NormalTok{  \})}\OperatorTok{;}
\NormalTok{\}}

\KeywordTok{async} \KeywordTok{function} \FunctionTok{main}\NormalTok{() \{}
  \KeywordTok{const}\NormalTok{ name }\OperatorTok{=} \ControlFlowTok{await} \FunctionTok{askQuestion}\NormalTok{(}\StringTok{"What is your name? "}\NormalTok{)}\OperatorTok{;}
  \KeywordTok{const}\NormalTok{ ageStr }\OperatorTok{=} \ControlFlowTok{await} \FunctionTok{askQuestion}\NormalTok{(}\StringTok{"How old are you? "}\NormalTok{)}\OperatorTok{;}
  \KeywordTok{const}\NormalTok{ age }\OperatorTok{=} \PreprocessorTok{parseInt}\NormalTok{(ageStr)}\OperatorTok{;}

  \BuiltInTok{console}\OperatorTok{.}\FunctionTok{log}\NormalTok{(}\VerbatimStringTok{\textasciigrave{}Hello, }\SpecialCharTok{$\{}\NormalTok{name}\SpecialCharTok{\}}\VerbatimStringTok{!\textasciigrave{}}\NormalTok{)}\OperatorTok{;}
  \BuiltInTok{console}\OperatorTok{.}\FunctionTok{log}\NormalTok{(}\VerbatimStringTok{\textasciigrave{}Next year you will be }\SpecialCharTok{$\{}\NormalTok{age }\OperatorTok{+} \DecValTok{1}\SpecialCharTok{\}}\VerbatimStringTok{ years old.\textasciigrave{}}\NormalTok{)}\OperatorTok{;}

\NormalTok{  rl}\OperatorTok{.}\FunctionTok{close}\NormalTok{()}\OperatorTok{;}
\NormalTok{\}}

\FunctionTok{main}\NormalTok{()}\OperatorTok{;}
\end{Highlighting}
\end{Shaded}

\subsubsection{Helper Module for Input
(Reusable)}\label{helper-module-for-input-reusable}

Create \texttt{input.js} once and reuse across projects:

\begin{Shaded}
\begin{Highlighting}[]
\CommentTok{// input.js {-}{-}{-} save this file in your project}
\KeywordTok{const}\NormalTok{ readline }\OperatorTok{=} \PreprocessorTok{require}\NormalTok{(}\StringTok{"readline"}\NormalTok{)}\OperatorTok{;}

\KeywordTok{const}\NormalTok{ rl }\OperatorTok{=}\NormalTok{ readline}\OperatorTok{.}\FunctionTok{createInterface}\NormalTok{(\{}
  \DataTypeTok{input}\OperatorTok{:} \BuiltInTok{process}\OperatorTok{.}\AttributeTok{stdin}\OperatorTok{,}
  \DataTypeTok{output}\OperatorTok{:} \BuiltInTok{process}\OperatorTok{.}\AttributeTok{stdout}\OperatorTok{,}
\NormalTok{\})}\OperatorTok{;}

\KeywordTok{function} \FunctionTok{askQuestion}\NormalTok{(prompt) \{}
  \ControlFlowTok{return} \KeywordTok{new} \BuiltInTok{Promise}\NormalTok{((resolve) }\KeywordTok{=\textgreater{}}\NormalTok{ \{}
\NormalTok{    rl}\OperatorTok{.}\FunctionTok{question}\NormalTok{(prompt}\OperatorTok{,}\NormalTok{ resolve)}\OperatorTok{;}
\NormalTok{  \})}\OperatorTok{;}
\NormalTok{\}}

\KeywordTok{function} \FunctionTok{close}\NormalTok{() \{}
\NormalTok{  rl}\OperatorTok{.}\FunctionTok{close}\NormalTok{()}\OperatorTok{;}
\NormalTok{\}}

\NormalTok{module}\OperatorTok{.}\AttributeTok{exports} \OperatorTok{=}\NormalTok{ \{ askQuestion}\OperatorTok{,}\NormalTok{ close \}}\OperatorTok{;}
\end{Highlighting}
\end{Shaded}

Usage in your main file:

\begin{Shaded}
\begin{Highlighting}[]
\CommentTok{// main.js}
\KeywordTok{const}\NormalTok{ \{ askQuestion}\OperatorTok{,}\NormalTok{ close \} }\OperatorTok{=} \PreprocessorTok{require}\NormalTok{(}\StringTok{"./input"}\NormalTok{)}\OperatorTok{;}

\KeywordTok{async} \KeywordTok{function} \FunctionTok{main}\NormalTok{() \{}
  \KeywordTok{const}\NormalTok{ name }\OperatorTok{=} \ControlFlowTok{await} \FunctionTok{askQuestion}\NormalTok{(}\StringTok{"Enter your name: "}\NormalTok{)}\OperatorTok{;}
  \KeywordTok{const}\NormalTok{ numStr }\OperatorTok{=} \ControlFlowTok{await} \FunctionTok{askQuestion}\NormalTok{(}\StringTok{"Enter a number: "}\NormalTok{)}\OperatorTok{;}
  \KeywordTok{const}\NormalTok{ num }\OperatorTok{=} \PreprocessorTok{parseFloat}\NormalTok{(numStr)}\OperatorTok{;}

  \BuiltInTok{console}\OperatorTok{.}\FunctionTok{log}\NormalTok{(}\VerbatimStringTok{\textasciigrave{}Hello, }\SpecialCharTok{$\{}\NormalTok{name}\SpecialCharTok{\}}\VerbatimStringTok{!\textasciigrave{}}\NormalTok{)}\OperatorTok{;}
  \BuiltInTok{console}\OperatorTok{.}\FunctionTok{log}\NormalTok{(}\VerbatimStringTok{\textasciigrave{}Your number squared is: }\SpecialCharTok{$\{}\NormalTok{num }\OperatorTok{**} \DecValTok{2}\SpecialCharTok{\}}\VerbatimStringTok{\textasciigrave{}}\NormalTok{)}\OperatorTok{;}

  \FunctionTok{close}\NormalTok{()}\OperatorTok{;}
\NormalTok{\}}

\FunctionTok{main}\NormalTok{()}\OperatorTok{;}
\end{Highlighting}
\end{Shaded}

\subsubsection{Examples}\label{examples-1}

\textbf{Example 2.2.1: Temperature Converter}

\begin{Shaded}
\begin{Highlighting}[]
\KeywordTok{const}\NormalTok{ \{ askQuestion}\OperatorTok{,}\NormalTok{ close \} }\OperatorTok{=} \PreprocessorTok{require}\NormalTok{(}\StringTok{"./input"}\NormalTok{)}\OperatorTok{;}

\KeywordTok{async} \KeywordTok{function} \FunctionTok{main}\NormalTok{() \{}
  \KeywordTok{const}\NormalTok{ fahrenheitStr }\OperatorTok{=} \ControlFlowTok{await} \FunctionTok{askQuestion}\NormalTok{(}\StringTok{"Enter temperature in Fahrenheit: "}\NormalTok{)}\OperatorTok{;}
  \KeywordTok{const}\NormalTok{ fahrenheit }\OperatorTok{=} \PreprocessorTok{parseFloat}\NormalTok{(fahrenheitStr)}\OperatorTok{;}

  \KeywordTok{const}\NormalTok{ celsius }\OperatorTok{=}\NormalTok{ ((fahrenheit }\OperatorTok{{-}} \DecValTok{32}\NormalTok{) }\OperatorTok{*} \DecValTok{5}\NormalTok{) }\OperatorTok{/} \DecValTok{9}\OperatorTok{;}

  \BuiltInTok{console}\OperatorTok{.}\FunctionTok{log}\NormalTok{(}\VerbatimStringTok{\textasciigrave{}}\SpecialCharTok{$\{}\NormalTok{fahrenheit}\SpecialCharTok{\}}\VerbatimStringTok{°F is }\SpecialCharTok{$\{}\NormalTok{celsius}\OperatorTok{.}\FunctionTok{toFixed}\NormalTok{(}\DecValTok{1}\NormalTok{)}\SpecialCharTok{\}}\VerbatimStringTok{°C\textasciigrave{}}\NormalTok{)}\OperatorTok{;}
  \FunctionTok{close}\NormalTok{()}\OperatorTok{;}
\NormalTok{\}}

\FunctionTok{main}\NormalTok{()}\OperatorTok{;}
\end{Highlighting}
\end{Shaded}

\textbf{Example 2.2.2: Circle Calculator}

\begin{Shaded}
\begin{Highlighting}[]
\KeywordTok{const}\NormalTok{ \{ askQuestion}\OperatorTok{,}\NormalTok{ close \} }\OperatorTok{=} \PreprocessorTok{require}\NormalTok{(}\StringTok{"./input"}\NormalTok{)}\OperatorTok{;}

\KeywordTok{async} \KeywordTok{function} \FunctionTok{main}\NormalTok{() \{}
  \KeywordTok{const}\NormalTok{ radiusStr }\OperatorTok{=} \ControlFlowTok{await} \FunctionTok{askQuestion}\NormalTok{(}\StringTok{"Enter the radius: "}\NormalTok{)}\OperatorTok{;}
  \KeywordTok{const}\NormalTok{ radius }\OperatorTok{=} \PreprocessorTok{parseFloat}\NormalTok{(radiusStr)}\OperatorTok{;}

  \KeywordTok{const}\NormalTok{ area }\OperatorTok{=} \BuiltInTok{Math}\OperatorTok{.}\ConstantTok{PI} \OperatorTok{*}\NormalTok{ radius }\OperatorTok{**} \DecValTok{2}\OperatorTok{;}
  \KeywordTok{const}\NormalTok{ circumference }\OperatorTok{=} \DecValTok{2} \OperatorTok{*} \BuiltInTok{Math}\OperatorTok{.}\ConstantTok{PI} \OperatorTok{*}\NormalTok{ radius}\OperatorTok{;}

  \BuiltInTok{console}\OperatorTok{.}\FunctionTok{log}\NormalTok{(}\VerbatimStringTok{\textasciigrave{}Area: }\SpecialCharTok{$\{}\NormalTok{area}\OperatorTok{.}\FunctionTok{toFixed}\NormalTok{(}\DecValTok{2}\NormalTok{)}\SpecialCharTok{\}}\VerbatimStringTok{\textasciigrave{}}\NormalTok{)}\OperatorTok{;}
  \BuiltInTok{console}\OperatorTok{.}\FunctionTok{log}\NormalTok{(}\VerbatimStringTok{\textasciigrave{}Circumference: }\SpecialCharTok{$\{}\NormalTok{circumference}\OperatorTok{.}\FunctionTok{toFixed}\NormalTok{(}\DecValTok{2}\NormalTok{)}\SpecialCharTok{\}}\VerbatimStringTok{\textasciigrave{}}\NormalTok{)}\OperatorTok{;}
  \FunctionTok{close}\NormalTok{()}\OperatorTok{;}
\NormalTok{\}}

\FunctionTok{main}\NormalTok{()}\OperatorTok{;}
\end{Highlighting}
\end{Shaded}

\subsubsection{Skills Checklist}\label{skills-checklist-4}

\begin{itemize}
\tightlist
\item[$\square$]
  I can use \texttt{Math} object methods and constants.
\item[$\square$]
  I can get user input using the \texttt{readline} module.
\item[$\square$]
  I understand \texttt{async}/\texttt{await} for asynchronous input.
\item[$\square$]
  I can convert string input to numbers using \texttt{parseInt()} and
  \texttt{parseFloat()}.
\end{itemize}

\begin{center}\rule{0.5\linewidth}{0.5pt}\end{center}

\section{Unit 3: Control Flow --- Decisions and
Loops}\label{unit-3-control-flow-decisions-and-loops}

\begin{center}\rule{0.5\linewidth}{0.5pt}\end{center}

\subsection{\texorpdfstring{Lesson 3.1: Boolean Expressions and
\texttt{if}
Statements}{Lesson 3.1: Boolean Expressions and if Statements}}\label{lesson-3.1-boolean-expressions-and-if-statements}

\subsubsection{Comparison Operators}\label{comparison-operators}

{\def\LTcaptype{none} % do not increment counter
\begin{longtable}[]{@{}lll@{}}
\toprule\noalign{}
Operator & Meaning & Note \\
\midrule\noalign{}
\endhead
\bottomrule\noalign{}
\endlastfoot
\texttt{===} & Strict equal & Checks value AND type (preferred) \\
\texttt{!==} & Strict not equal & Checks value AND type (preferred) \\
\texttt{==} & Loose equal & Type coercion (avoid) \\
\texttt{!=} & Loose not equal & Type coercion (avoid) \\
\texttt{\textless{}} & Less than & \\
\texttt{\textgreater{}} & Greater than & \\
\texttt{\textless{}=} & Less than or equal & \\
\texttt{\textgreater{}=} & Greater than or equal & \\
\end{longtable}
}

\textbf{Always use \texttt{===} and \texttt{!==}} --- \texttt{==} does
unexpected type conversions.

\subsubsection{Logical Operators}\label{logical-operators}

{\def\LTcaptype{none} % do not increment counter
\begin{longtable}[]{@{}lll@{}}
\toprule\noalign{}
Operator & Meaning & Example \\
\midrule\noalign{}
\endhead
\bottomrule\noalign{}
\endlastfoot
\texttt{\&\&} & AND & \texttt{age\ \textgreater{}=\ 18\ \&\&\ hasId} \\
\texttt{\textbackslash{}\textbar{}\textbackslash{}\textbar{}} & OR &
\texttt{day\ ===\ \textquotesingle{}Sat\textquotesingle{}\ \textbackslash{}\textbar{}\textbackslash{}\textbar{}\ day\ ===\ \textquotesingle{}Sun\textquotesingle{}} \\
\texttt{!} & NOT & \texttt{!isRaining} \\
\end{longtable}
}

\subsubsection{\texorpdfstring{\texttt{if-else\ if-else}}{if-else if-else}}\label{if-else-if-else}

\begin{Shaded}
\begin{Highlighting}[]
\KeywordTok{const}\NormalTok{ score }\OperatorTok{=} \DecValTok{85}\OperatorTok{;}

\ControlFlowTok{if}\NormalTok{ (score }\OperatorTok{\textgreater{}=} \DecValTok{90}\NormalTok{) \{}
  \BuiltInTok{console}\OperatorTok{.}\FunctionTok{log}\NormalTok{(}\StringTok{"Grade: A"}\NormalTok{)}\OperatorTok{;}
\NormalTok{\} }\ControlFlowTok{else} \ControlFlowTok{if}\NormalTok{ (score }\OperatorTok{\textgreater{}=} \DecValTok{80}\NormalTok{) \{}
  \BuiltInTok{console}\OperatorTok{.}\FunctionTok{log}\NormalTok{(}\StringTok{"Grade: B"}\NormalTok{)}\OperatorTok{;}
\NormalTok{\} }\ControlFlowTok{else} \ControlFlowTok{if}\NormalTok{ (score }\OperatorTok{\textgreater{}=} \DecValTok{70}\NormalTok{) \{}
  \BuiltInTok{console}\OperatorTok{.}\FunctionTok{log}\NormalTok{(}\StringTok{"Grade: C"}\NormalTok{)}\OperatorTok{;}
\NormalTok{\} }\ControlFlowTok{else}\NormalTok{ \{}
  \BuiltInTok{console}\OperatorTok{.}\FunctionTok{log}\NormalTok{(}\StringTok{"Grade: F"}\NormalTok{)}\OperatorTok{;}
\NormalTok{\}}
\end{Highlighting}
\end{Shaded}

\subsubsection{Ternary Operator}\label{ternary-operator}

\begin{Shaded}
\begin{Highlighting}[]
\KeywordTok{const}\NormalTok{ age }\OperatorTok{=} \DecValTok{20}\OperatorTok{;}
\KeywordTok{const}\NormalTok{ status }\OperatorTok{=}\NormalTok{ age }\OperatorTok{\textgreater{}=} \DecValTok{18} \OperatorTok{?} \StringTok{"adult"} \OperatorTok{:} \StringTok{"minor"}\OperatorTok{;}
\BuiltInTok{console}\OperatorTok{.}\FunctionTok{log}\NormalTok{(status)}\OperatorTok{;} \CommentTok{// "adult"}
\end{Highlighting}
\end{Shaded}

\subsubsection{\texorpdfstring{\texttt{switch}
Statement}{switch Statement}}\label{switch-statement}

\begin{Shaded}
\begin{Highlighting}[]
\KeywordTok{const}\NormalTok{ day }\OperatorTok{=} \StringTok{"Wednesday"}\OperatorTok{;}
\ControlFlowTok{switch}\NormalTok{ (day) \{}
  \ControlFlowTok{case} \StringTok{"Monday"}\OperatorTok{:}
  \ControlFlowTok{case} \StringTok{"Tuesday"}\OperatorTok{:}
  \ControlFlowTok{case} \StringTok{"Wednesday"}\OperatorTok{:}
  \ControlFlowTok{case} \StringTok{"Thursday"}\OperatorTok{:}
  \ControlFlowTok{case} \StringTok{"Friday"}\OperatorTok{:}
    \BuiltInTok{console}\OperatorTok{.}\FunctionTok{log}\NormalTok{(}\StringTok{"Weekday"}\NormalTok{)}\OperatorTok{;}
    \ControlFlowTok{break}\OperatorTok{;}
  \ControlFlowTok{case} \StringTok{"Saturday"}\OperatorTok{:}
  \ControlFlowTok{case} \StringTok{"Sunday"}\OperatorTok{:}
    \BuiltInTok{console}\OperatorTok{.}\FunctionTok{log}\NormalTok{(}\StringTok{"Weekend"}\NormalTok{)}\OperatorTok{;}
    \ControlFlowTok{break}\OperatorTok{;}
  \ControlFlowTok{default}\OperatorTok{:}
    \BuiltInTok{console}\OperatorTok{.}\FunctionTok{log}\NormalTok{(}\StringTok{"Unknown day"}\NormalTok{)}\OperatorTok{;}
\NormalTok{\}}
\end{Highlighting}
\end{Shaded}

\subsubsection{Truthy and Falsy Values}\label{truthy-and-falsy-values}

In JavaScript, the following values are \textbf{falsy} (evaluate to
\texttt{false} in conditions):

\begin{itemize}
\tightlist
\item
  \texttt{false}, \texttt{0}, \texttt{""} (empty string), \texttt{null},
  \texttt{undefined}, \texttt{NaN}
\end{itemize}

Everything else is \textbf{truthy}.

\begin{Shaded}
\begin{Highlighting}[]
\KeywordTok{const}\NormalTok{ name }\OperatorTok{=} \StringTok{""}\OperatorTok{;}
\ControlFlowTok{if}\NormalTok{ (name) \{}
  \BuiltInTok{console}\OperatorTok{.}\FunctionTok{log}\NormalTok{(}\StringTok{"Name is set"}\NormalTok{)}\OperatorTok{;}
\NormalTok{\} }\ControlFlowTok{else}\NormalTok{ \{}
  \BuiltInTok{console}\OperatorTok{.}\FunctionTok{log}\NormalTok{(}\StringTok{"Name is empty"}\NormalTok{)}\OperatorTok{;} \CommentTok{// This runs}
\NormalTok{\}}
\end{Highlighting}
\end{Shaded}

\subsubsection{Examples}\label{examples-2}

\textbf{Example 3.1.1: Grade Calculator}

\begin{Shaded}
\begin{Highlighting}[]
\KeywordTok{const}\NormalTok{ \{ askQuestion}\OperatorTok{,}\NormalTok{ close \} }\OperatorTok{=} \PreprocessorTok{require}\NormalTok{(}\StringTok{"./input"}\NormalTok{)}\OperatorTok{;}

\KeywordTok{async} \KeywordTok{function} \FunctionTok{main}\NormalTok{() \{}
  \KeywordTok{const}\NormalTok{ scoreStr }\OperatorTok{=} \ControlFlowTok{await} \FunctionTok{askQuestion}\NormalTok{(}\StringTok{"Enter your score (0{-}100): "}\NormalTok{)}\OperatorTok{;}
  \KeywordTok{const}\NormalTok{ score }\OperatorTok{=} \PreprocessorTok{parseInt}\NormalTok{(scoreStr)}\OperatorTok{;}

  \KeywordTok{let}\NormalTok{ grade}\OperatorTok{;}
  \ControlFlowTok{if}\NormalTok{ (score }\OperatorTok{\textgreater{}=} \DecValTok{90}\NormalTok{) grade }\OperatorTok{=} \StringTok{"A"}\OperatorTok{;}
  \ControlFlowTok{else} \ControlFlowTok{if}\NormalTok{ (score }\OperatorTok{\textgreater{}=} \DecValTok{80}\NormalTok{) grade }\OperatorTok{=} \StringTok{"B"}\OperatorTok{;}
  \ControlFlowTok{else} \ControlFlowTok{if}\NormalTok{ (score }\OperatorTok{\textgreater{}=} \DecValTok{70}\NormalTok{) grade }\OperatorTok{=} \StringTok{"C"}\OperatorTok{;}
  \ControlFlowTok{else} \ControlFlowTok{if}\NormalTok{ (score }\OperatorTok{\textgreater{}=} \DecValTok{60}\NormalTok{) grade }\OperatorTok{=} \StringTok{"D"}\OperatorTok{;}
  \ControlFlowTok{else}\NormalTok{ grade }\OperatorTok{=} \StringTok{"F"}\OperatorTok{;}

  \BuiltInTok{console}\OperatorTok{.}\FunctionTok{log}\NormalTok{(}\VerbatimStringTok{\textasciigrave{}Your grade: }\SpecialCharTok{$\{}\NormalTok{grade}\SpecialCharTok{\}}\VerbatimStringTok{\textasciigrave{}}\NormalTok{)}\OperatorTok{;}
  \FunctionTok{close}\NormalTok{()}\OperatorTok{;}
\NormalTok{\}}

\FunctionTok{main}\NormalTok{()}\OperatorTok{;}
\end{Highlighting}
\end{Shaded}

\subsubsection{Skills Checklist}\label{skills-checklist-5}

\begin{itemize}
\tightlist
\item[$\square$]
  I use \texttt{===} and \texttt{!==} for comparisons in JavaScript.
\item[$\square$]
  I can write \texttt{if}, \texttt{else\ if}, and \texttt{else}
  statements.
\item[$\square$]
  I understand truthy and falsy values.
\item[$\square$]
  I can write a \texttt{switch} statement.
\end{itemize}

\begin{center}\rule{0.5\linewidth}{0.5pt}\end{center}

\subsection{\texorpdfstring{Lesson 3.2: \texttt{for}
Loops}{Lesson 3.2: for Loops}}\label{lesson-3.2-for-loops}

\begin{Shaded}
\begin{Highlighting}[]
\CommentTok{// Basic for loop}
\ControlFlowTok{for}\NormalTok{ (}\KeywordTok{let}\NormalTok{ i }\OperatorTok{=} \DecValTok{0}\OperatorTok{;}\NormalTok{ i }\OperatorTok{\textless{}} \DecValTok{5}\OperatorTok{;}\NormalTok{ i}\OperatorTok{++}\NormalTok{) \{}
  \BuiltInTok{console}\OperatorTok{.}\FunctionTok{log}\NormalTok{(}\VerbatimStringTok{\textasciigrave{}i = }\SpecialCharTok{$\{}\NormalTok{i}\SpecialCharTok{\}}\VerbatimStringTok{\textasciigrave{}}\NormalTok{)}\OperatorTok{;}
\NormalTok{\}}

\CommentTok{// Counting down}
\ControlFlowTok{for}\NormalTok{ (}\KeywordTok{let}\NormalTok{ i }\OperatorTok{=} \DecValTok{10}\OperatorTok{;}\NormalTok{ i }\OperatorTok{\textgreater{}} \DecValTok{0}\OperatorTok{;}\NormalTok{ i}\OperatorTok{{-}{-}}\NormalTok{) \{}
  \BuiltInTok{console}\OperatorTok{.}\FunctionTok{log}\NormalTok{(i)}\OperatorTok{;}
\NormalTok{\}}
\BuiltInTok{console}\OperatorTok{.}\FunctionTok{log}\NormalTok{(}\StringTok{"Blast off!"}\NormalTok{)}\OperatorTok{;}

\CommentTok{// for...of {-}{-}{-} iterate over array elements (like Python\textquotesingle{}s for{-}in)}
\KeywordTok{const}\NormalTok{ fruits }\OperatorTok{=}\NormalTok{ [}\StringTok{"apple"}\OperatorTok{,} \StringTok{"banana"}\OperatorTok{,} \StringTok{"cherry"}\NormalTok{]}\OperatorTok{;}
\ControlFlowTok{for}\NormalTok{ (}\KeywordTok{const}\NormalTok{ fruit }\KeywordTok{of}\NormalTok{ fruits) \{}
  \BuiltInTok{console}\OperatorTok{.}\FunctionTok{log}\NormalTok{(}\VerbatimStringTok{\textasciigrave{}I like }\SpecialCharTok{$\{}\NormalTok{fruit}\SpecialCharTok{\}}\VerbatimStringTok{.\textasciigrave{}}\NormalTok{)}\OperatorTok{;}
\NormalTok{\}}

\CommentTok{// for...in {-}{-}{-} iterate over object keys}
\KeywordTok{const}\NormalTok{ person }\OperatorTok{=}\NormalTok{ \{ }\DataTypeTok{name}\OperatorTok{:} \StringTok{"Alice"}\OperatorTok{,} \DataTypeTok{age}\OperatorTok{:} \DecValTok{15}\OperatorTok{,} \DataTypeTok{city}\OperatorTok{:} \StringTok{"Portland"}\NormalTok{ \}}\OperatorTok{;}
\ControlFlowTok{for}\NormalTok{ (}\KeywordTok{const}\NormalTok{ key }\KeywordTok{in}\NormalTok{ person) \{}
  \BuiltInTok{console}\OperatorTok{.}\FunctionTok{log}\NormalTok{(}\VerbatimStringTok{\textasciigrave{}}\SpecialCharTok{$\{}\NormalTok{key}\SpecialCharTok{\}}\VerbatimStringTok{: }\SpecialCharTok{$\{}\NormalTok{person[key]}\SpecialCharTok{\}}\VerbatimStringTok{\textasciigrave{}}\NormalTok{)}\OperatorTok{;}
\NormalTok{\}}
\end{Highlighting}
\end{Shaded}

\subsubsection{\texorpdfstring{\texttt{Array.forEach()}}{Array.forEach()}}\label{array.foreach}

\begin{Shaded}
\begin{Highlighting}[]
\KeywordTok{const}\NormalTok{ numbers }\OperatorTok{=}\NormalTok{ [}\DecValTok{1}\OperatorTok{,} \DecValTok{2}\OperatorTok{,} \DecValTok{3}\OperatorTok{,} \DecValTok{4}\OperatorTok{,} \DecValTok{5}\NormalTok{]}\OperatorTok{;}

\NormalTok{numbers}\OperatorTok{.}\FunctionTok{forEach}\NormalTok{((num) }\KeywordTok{=\textgreater{}}\NormalTok{ \{}
  \BuiltInTok{console}\OperatorTok{.}\FunctionTok{log}\NormalTok{(num }\OperatorTok{*} \DecValTok{2}\NormalTok{)}\OperatorTok{;}
\NormalTok{\})}\OperatorTok{;}

\CommentTok{// With index}
\NormalTok{numbers}\OperatorTok{.}\FunctionTok{forEach}\NormalTok{((num}\OperatorTok{,}\NormalTok{ index) }\KeywordTok{=\textgreater{}}\NormalTok{ \{}
  \BuiltInTok{console}\OperatorTok{.}\FunctionTok{log}\NormalTok{(}\VerbatimStringTok{\textasciigrave{}Index }\SpecialCharTok{$\{}\NormalTok{index}\SpecialCharTok{\}}\VerbatimStringTok{: }\SpecialCharTok{$\{}\NormalTok{num}\SpecialCharTok{\}}\VerbatimStringTok{\textasciigrave{}}\NormalTok{)}\OperatorTok{;}
\NormalTok{\})}\OperatorTok{;}
\end{Highlighting}
\end{Shaded}

\subsubsection{Examples}\label{examples-3}

\textbf{Example 3.2.1: Multiplication Table}

\begin{Shaded}
\begin{Highlighting}[]
\KeywordTok{const}\NormalTok{ n }\OperatorTok{=} \DecValTok{5}\OperatorTok{;}
\BuiltInTok{console}\OperatorTok{.}\FunctionTok{log}\NormalTok{(}\VerbatimStringTok{\textasciigrave{}Multiplication table for }\SpecialCharTok{$\{}\NormalTok{n}\SpecialCharTok{\}}\VerbatimStringTok{:\textasciigrave{}}\NormalTok{)}\OperatorTok{;}
\ControlFlowTok{for}\NormalTok{ (}\KeywordTok{let}\NormalTok{ i }\OperatorTok{=} \DecValTok{1}\OperatorTok{;}\NormalTok{ i }\OperatorTok{\textless{}=} \DecValTok{10}\OperatorTok{;}\NormalTok{ i}\OperatorTok{++}\NormalTok{) \{}
  \BuiltInTok{console}\OperatorTok{.}\FunctionTok{log}\NormalTok{(}\VerbatimStringTok{\textasciigrave{}}\SpecialCharTok{$\{}\NormalTok{n}\SpecialCharTok{\}}\VerbatimStringTok{ x }\SpecialCharTok{$\{}\NormalTok{i}\SpecialCharTok{\}}\VerbatimStringTok{ = }\SpecialCharTok{$\{}\NormalTok{n }\OperatorTok{*}\NormalTok{ i}\SpecialCharTok{\}}\VerbatimStringTok{\textasciigrave{}}\NormalTok{)}\OperatorTok{;}
\NormalTok{\}}
\end{Highlighting}
\end{Shaded}

\subsubsection{Skills Checklist}\label{skills-checklist-6}

\begin{itemize}
\tightlist
\item[$\square$]
  I can write a \texttt{for} loop with initialization, condition, and
  update.
\item[$\square$]
  I can use \texttt{for...of} to iterate over arrays.
\item[$\square$]
  I can use \texttt{for...in} to iterate over object keys.
\item[$\square$]
  I can use \texttt{Array.forEach()}.
\end{itemize}

\begin{center}\rule{0.5\linewidth}{0.5pt}\end{center}

\subsection{\texorpdfstring{Lesson 3.3: \texttt{while}
Loops}{Lesson 3.3: while Loops}}\label{lesson-3.3-while-loops}

\begin{Shaded}
\begin{Highlighting}[]
\KeywordTok{let}\NormalTok{ count }\OperatorTok{=} \DecValTok{1}\OperatorTok{;}
\ControlFlowTok{while}\NormalTok{ (count }\OperatorTok{\textless{}=} \DecValTok{5}\NormalTok{) \{}
  \BuiltInTok{console}\OperatorTok{.}\FunctionTok{log}\NormalTok{(}\VerbatimStringTok{\textasciigrave{}Count: }\SpecialCharTok{$\{}\NormalTok{count}\SpecialCharTok{\}}\VerbatimStringTok{\textasciigrave{}}\NormalTok{)}\OperatorTok{;}
\NormalTok{  count}\OperatorTok{++;}
\NormalTok{\}}

\CommentTok{// do{-}while {-}{-}{-} runs at least once}
\KeywordTok{let}\NormalTok{ num}\OperatorTok{;}
\ControlFlowTok{do}\NormalTok{ \{}
\NormalTok{  num }\OperatorTok{=} \BuiltInTok{Math}\OperatorTok{.}\FunctionTok{floor}\NormalTok{(}\BuiltInTok{Math}\OperatorTok{.}\FunctionTok{random}\NormalTok{() }\OperatorTok{*} \DecValTok{10}\NormalTok{)}\OperatorTok{;}
  \BuiltInTok{console}\OperatorTok{.}\FunctionTok{log}\NormalTok{(}\VerbatimStringTok{\textasciigrave{}Rolled: }\SpecialCharTok{$\{}\NormalTok{num}\SpecialCharTok{\}}\VerbatimStringTok{\textasciigrave{}}\NormalTok{)}\OperatorTok{;}
\NormalTok{\} }\ControlFlowTok{while}\NormalTok{ (num }\OperatorTok{\textless{}} \DecValTok{5}\NormalTok{)}\OperatorTok{;}
\end{Highlighting}
\end{Shaded}

\subsubsection{\texorpdfstring{\texttt{break} and
\texttt{continue}}{break and continue}}\label{break-and-continue}

\begin{Shaded}
\begin{Highlighting}[]
\CommentTok{// break {-}{-}{-} exit entirely}
\ControlFlowTok{for}\NormalTok{ (}\KeywordTok{let}\NormalTok{ i }\OperatorTok{=} \DecValTok{0}\OperatorTok{;}\NormalTok{ i }\OperatorTok{\textless{}} \DecValTok{10}\OperatorTok{;}\NormalTok{ i}\OperatorTok{++}\NormalTok{) \{}
  \ControlFlowTok{if}\NormalTok{ (i }\OperatorTok{===} \DecValTok{5}\NormalTok{) }\ControlFlowTok{break}\OperatorTok{;}
  \BuiltInTok{console}\OperatorTok{.}\FunctionTok{log}\NormalTok{(i)}\OperatorTok{;} \CommentTok{// Prints 0{-}4}
\NormalTok{\}}

\CommentTok{// continue {-}{-}{-} skip iteration}
\ControlFlowTok{for}\NormalTok{ (}\KeywordTok{let}\NormalTok{ i }\OperatorTok{=} \DecValTok{0}\OperatorTok{;}\NormalTok{ i }\OperatorTok{\textless{}} \DecValTok{10}\OperatorTok{;}\NormalTok{ i}\OperatorTok{++}\NormalTok{) \{}
  \ControlFlowTok{if}\NormalTok{ (i }\OperatorTok{\%} \DecValTok{2} \OperatorTok{===} \DecValTok{0}\NormalTok{) }\ControlFlowTok{continue}\OperatorTok{;}
  \BuiltInTok{console}\OperatorTok{.}\FunctionTok{log}\NormalTok{(i)}\OperatorTok{;} \CommentTok{// Prints 1, 3, 5, 7, 9}
\NormalTok{\}}
\end{Highlighting}
\end{Shaded}

\subsubsection{Skills Checklist}\label{skills-checklist-7}

\begin{itemize}
\tightlist
\item[$\square$]
  I can write \texttt{while} and \texttt{do-while} loops.
\item[$\square$]
  I can use \texttt{break} and \texttt{continue} appropriately.
\end{itemize}

\begin{center}\rule{0.5\linewidth}{0.5pt}\end{center}

\section{Unit 4: Data Structures --- Arrays and
Objects}\label{unit-4-data-structures-arrays-and-objects}

\begin{center}\rule{0.5\linewidth}{0.5pt}\end{center}

\subsection{Lesson 4.1: Arrays}\label{lesson-4.1-arrays}

\subsubsection{Creating and Accessing
Arrays}\label{creating-and-accessing-arrays}

\begin{Shaded}
\begin{Highlighting}[]
\KeywordTok{const}\NormalTok{ colors }\OperatorTok{=}\NormalTok{ [}\StringTok{"red"}\OperatorTok{,} \StringTok{"blue"}\OperatorTok{,} \StringTok{"green"}\OperatorTok{,} \StringTok{"yellow"}\NormalTok{]}\OperatorTok{;}

\CommentTok{// Access}
\BuiltInTok{console}\OperatorTok{.}\FunctionTok{log}\NormalTok{(colors[}\DecValTok{0}\NormalTok{])}\OperatorTok{;} \CommentTok{// "red"}
\BuiltInTok{console}\OperatorTok{.}\FunctionTok{log}\NormalTok{(colors[colors}\OperatorTok{.}\AttributeTok{length} \OperatorTok{{-}} \DecValTok{1}\NormalTok{])}\OperatorTok{;} \CommentTok{// "yellow"}

\CommentTok{// Modify}
\NormalTok{colors[}\DecValTok{1}\NormalTok{] }\OperatorTok{=} \StringTok{"purple"}\OperatorTok{;}

\CommentTok{// Check}
\BuiltInTok{console}\OperatorTok{.}\FunctionTok{log}\NormalTok{(colors}\OperatorTok{.}\AttributeTok{length}\NormalTok{)}\OperatorTok{;} \CommentTok{// 4}
\BuiltInTok{console}\OperatorTok{.}\FunctionTok{log}\NormalTok{(colors}\OperatorTok{.}\FunctionTok{includes}\NormalTok{(}\StringTok{"red"}\NormalTok{))}\OperatorTok{;} \CommentTok{// true}
\BuiltInTok{console}\OperatorTok{.}\FunctionTok{log}\NormalTok{(colors}\OperatorTok{.}\FunctionTok{indexOf}\NormalTok{(}\StringTok{"green"}\NormalTok{))}\OperatorTok{;} \CommentTok{// 2}
\end{Highlighting}
\end{Shaded}

\subsubsection{Array Methods}\label{array-methods}

\begin{Shaded}
\begin{Highlighting}[]
\KeywordTok{const}\NormalTok{ fruits }\OperatorTok{=}\NormalTok{ [}\StringTok{"apple"}\OperatorTok{,} \StringTok{"banana"}\NormalTok{]}\OperatorTok{;}

\CommentTok{// Adding}
\NormalTok{fruits}\OperatorTok{.}\FunctionTok{push}\NormalTok{(}\StringTok{"cherry"}\NormalTok{)}\OperatorTok{;} \CommentTok{// Add to end {-}\textgreater{} ["apple", "banana", "cherry"]}
\NormalTok{fruits}\OperatorTok{.}\FunctionTok{unshift}\NormalTok{(}\StringTok{"avocado"}\NormalTok{)}\OperatorTok{;} \CommentTok{// Add to start {-}\textgreater{} ["avocado", "apple", "banana", "cherry"]}

\CommentTok{// Removing}
\KeywordTok{const}\NormalTok{ last }\OperatorTok{=}\NormalTok{ fruits}\OperatorTok{.}\FunctionTok{pop}\NormalTok{()}\OperatorTok{;} \CommentTok{// Remove and return last}
\KeywordTok{const}\NormalTok{ first }\OperatorTok{=}\NormalTok{ fruits}\OperatorTok{.}\FunctionTok{shift}\NormalTok{()}\OperatorTok{;} \CommentTok{// Remove and return first}

\CommentTok{// Slicing (non{-}destructive)}
\KeywordTok{const}\NormalTok{ sub }\OperatorTok{=}\NormalTok{ fruits}\OperatorTok{.}\FunctionTok{slice}\NormalTok{(}\DecValTok{1}\OperatorTok{,} \DecValTok{3}\NormalTok{)}\OperatorTok{;} \CommentTok{// Returns new array, original unchanged}

\CommentTok{// Splicing (destructive {-}{-}{-} modifies original)}
\NormalTok{fruits}\OperatorTok{.}\FunctionTok{splice}\NormalTok{(}\DecValTok{1}\OperatorTok{,} \DecValTok{1}\NormalTok{)}\OperatorTok{;} \CommentTok{// Remove 1 element at index 1}
\NormalTok{fruits}\OperatorTok{.}\FunctionTok{splice}\NormalTok{(}\DecValTok{1}\OperatorTok{,} \DecValTok{0}\OperatorTok{,} \StringTok{"mango"}\NormalTok{)}\OperatorTok{;} \CommentTok{// Insert "mango" at index 1}

\CommentTok{// Joining}
\KeywordTok{const}\NormalTok{ str }\OperatorTok{=}\NormalTok{ fruits}\OperatorTok{.}\FunctionTok{join}\NormalTok{(}\StringTok{", "}\NormalTok{)}\OperatorTok{;} \CommentTok{// "apple, mango, cherry"}

\CommentTok{// Sorting}
\KeywordTok{const}\NormalTok{ nums }\OperatorTok{=}\NormalTok{ [}\DecValTok{3}\OperatorTok{,} \DecValTok{1}\OperatorTok{,} \DecValTok{4}\OperatorTok{,} \DecValTok{1}\OperatorTok{,} \DecValTok{5}\OperatorTok{,} \DecValTok{9}\OperatorTok{,} \DecValTok{2}\OperatorTok{,} \DecValTok{6}\NormalTok{]}\OperatorTok{;}
\NormalTok{nums}\OperatorTok{.}\FunctionTok{sort}\NormalTok{((a}\OperatorTok{,}\NormalTok{ b) }\KeywordTok{=\textgreater{}}\NormalTok{ a }\OperatorTok{{-}}\NormalTok{ b)}\OperatorTok{;} \CommentTok{// Ascending: [1, 1, 2, 3, 4, 5, 6, 9]}
\NormalTok{nums}\OperatorTok{.}\FunctionTok{sort}\NormalTok{((a}\OperatorTok{,}\NormalTok{ b) }\KeywordTok{=\textgreater{}}\NormalTok{ b }\OperatorTok{{-}}\NormalTok{ a)}\OperatorTok{;} \CommentTok{// Descending}

\CommentTok{// Reversing}
\KeywordTok{const}\NormalTok{ reversed }\OperatorTok{=}\NormalTok{ [}\OperatorTok{...}\NormalTok{fruits]}\OperatorTok{.}\FunctionTok{reverse}\NormalTok{()}\OperatorTok{;} \CommentTok{// Use spread to avoid mutating original}

\CommentTok{// Concatenating}
\KeywordTok{const}\NormalTok{ combined }\OperatorTok{=}\NormalTok{ fruits}\OperatorTok{.}\FunctionTok{concat}\NormalTok{(nums)}\OperatorTok{;}
\end{Highlighting}
\end{Shaded}

\subsubsection{Functional Array Methods}\label{functional-array-methods}

These are powerful for data transformation:

\begin{Shaded}
\begin{Highlighting}[]
\KeywordTok{const}\NormalTok{ numbers }\OperatorTok{=}\NormalTok{ [}\DecValTok{1}\OperatorTok{,} \DecValTok{2}\OperatorTok{,} \DecValTok{3}\OperatorTok{,} \DecValTok{4}\OperatorTok{,} \DecValTok{5}\OperatorTok{,} \DecValTok{6}\OperatorTok{,} \DecValTok{7}\OperatorTok{,} \DecValTok{8}\OperatorTok{,} \DecValTok{9}\OperatorTok{,} \DecValTok{10}\NormalTok{]}\OperatorTok{;}

\CommentTok{// map {-}{-}{-} transform each element, returns new array}
\KeywordTok{const}\NormalTok{ doubled }\OperatorTok{=}\NormalTok{ numbers}\OperatorTok{.}\FunctionTok{map}\NormalTok{((n) }\KeywordTok{=\textgreater{}}\NormalTok{ n }\OperatorTok{*} \DecValTok{2}\NormalTok{)}\OperatorTok{;}
\BuiltInTok{console}\OperatorTok{.}\FunctionTok{log}\NormalTok{(doubled)}\OperatorTok{;} \CommentTok{// [2, 4, 6, 8, 10, 12, 14, 16, 18, 20]}

\CommentTok{// filter {-}{-}{-} keep elements matching condition}
\KeywordTok{const}\NormalTok{ evens }\OperatorTok{=}\NormalTok{ numbers}\OperatorTok{.}\FunctionTok{filter}\NormalTok{((n) }\KeywordTok{=\textgreater{}}\NormalTok{ n }\OperatorTok{\%} \DecValTok{2} \OperatorTok{===} \DecValTok{0}\NormalTok{)}\OperatorTok{;}
\BuiltInTok{console}\OperatorTok{.}\FunctionTok{log}\NormalTok{(evens)}\OperatorTok{;} \CommentTok{// [2, 4, 6, 8, 10]}

\CommentTok{// reduce {-}{-}{-} combine elements into single value}
\KeywordTok{const}\NormalTok{ sum }\OperatorTok{=}\NormalTok{ numbers}\OperatorTok{.}\FunctionTok{reduce}\NormalTok{((accumulator}\OperatorTok{,}\NormalTok{ current) }\KeywordTok{=\textgreater{}}\NormalTok{ accumulator }\OperatorTok{+}\NormalTok{ current}\OperatorTok{,} \DecValTok{0}\NormalTok{)}\OperatorTok{;}
\BuiltInTok{console}\OperatorTok{.}\FunctionTok{log}\NormalTok{(sum)}\OperatorTok{;} \CommentTok{// 55}

\CommentTok{// find {-}{-}{-} return first matching element}
\KeywordTok{const}\NormalTok{ firstOver5 }\OperatorTok{=}\NormalTok{ numbers}\OperatorTok{.}\FunctionTok{find}\NormalTok{((n) }\KeywordTok{=\textgreater{}}\NormalTok{ n }\OperatorTok{\textgreater{}} \DecValTok{5}\NormalTok{)}\OperatorTok{;}
\BuiltInTok{console}\OperatorTok{.}\FunctionTok{log}\NormalTok{(firstOver5)}\OperatorTok{;} \CommentTok{// 6}

\CommentTok{// some {-}{-}{-} true if any element matches}
\KeywordTok{const}\NormalTok{ hasEven }\OperatorTok{=}\NormalTok{ numbers}\OperatorTok{.}\FunctionTok{some}\NormalTok{((n) }\KeywordTok{=\textgreater{}}\NormalTok{ n }\OperatorTok{\%} \DecValTok{2} \OperatorTok{===} \DecValTok{0}\NormalTok{)}\OperatorTok{;}
\BuiltInTok{console}\OperatorTok{.}\FunctionTok{log}\NormalTok{(hasEven)}\OperatorTok{;} \CommentTok{// true}

\CommentTok{// every {-}{-}{-} true if all elements match}
\KeywordTok{const}\NormalTok{ allPositive }\OperatorTok{=}\NormalTok{ numbers}\OperatorTok{.}\FunctionTok{every}\NormalTok{((n) }\KeywordTok{=\textgreater{}}\NormalTok{ n }\OperatorTok{\textgreater{}} \DecValTok{0}\NormalTok{)}\OperatorTok{;}
\BuiltInTok{console}\OperatorTok{.}\FunctionTok{log}\NormalTok{(allPositive)}\OperatorTok{;} \CommentTok{// true}

\CommentTok{// Chaining}
\KeywordTok{const}\NormalTok{ result }\OperatorTok{=}\NormalTok{ numbers}
  \OperatorTok{.}\FunctionTok{filter}\NormalTok{((n) }\KeywordTok{=\textgreater{}}\NormalTok{ n }\OperatorTok{\%} \DecValTok{2} \OperatorTok{===} \DecValTok{0}\NormalTok{)}
  \OperatorTok{.}\FunctionTok{map}\NormalTok{((n) }\KeywordTok{=\textgreater{}}\NormalTok{ n }\OperatorTok{**} \DecValTok{2}\NormalTok{)}
  \OperatorTok{.}\FunctionTok{reduce}\NormalTok{((sum}\OperatorTok{,}\NormalTok{ n) }\KeywordTok{=\textgreater{}}\NormalTok{ sum }\OperatorTok{+}\NormalTok{ n}\OperatorTok{,} \DecValTok{0}\NormalTok{)}\OperatorTok{;}
\BuiltInTok{console}\OperatorTok{.}\FunctionTok{log}\NormalTok{(result)}\OperatorTok{;} \CommentTok{// 220  (4+16+36+64+100)}
\end{Highlighting}
\end{Shaded}

\subsubsection{Examples}\label{examples-4}

\textbf{Example 4.1.1: Array Analysis}

\begin{Shaded}
\begin{Highlighting}[]
\KeywordTok{const}\NormalTok{ scores }\OperatorTok{=}\NormalTok{ [}\DecValTok{85}\OperatorTok{,} \DecValTok{92}\OperatorTok{,} \DecValTok{78}\OperatorTok{,} \DecValTok{95}\OperatorTok{,} \DecValTok{88}\OperatorTok{,} \DecValTok{72}\OperatorTok{,} \DecValTok{90}\OperatorTok{,} \DecValTok{83}\NormalTok{]}\OperatorTok{;}

\KeywordTok{const}\NormalTok{ average }\OperatorTok{=}\NormalTok{ scores}\OperatorTok{.}\FunctionTok{reduce}\NormalTok{((sum}\OperatorTok{,}\NormalTok{ s) }\KeywordTok{=\textgreater{}}\NormalTok{ sum }\OperatorTok{+}\NormalTok{ s}\OperatorTok{,} \DecValTok{0}\NormalTok{) }\OperatorTok{/}\NormalTok{ scores}\OperatorTok{.}\AttributeTok{length}\OperatorTok{;}
\KeywordTok{const}\NormalTok{ max }\OperatorTok{=} \BuiltInTok{Math}\OperatorTok{.}\FunctionTok{max}\NormalTok{(}\OperatorTok{...}\NormalTok{scores)}\OperatorTok{;}
\KeywordTok{const}\NormalTok{ min }\OperatorTok{=} \BuiltInTok{Math}\OperatorTok{.}\FunctionTok{min}\NormalTok{(}\OperatorTok{...}\NormalTok{scores)}\OperatorTok{;}
\KeywordTok{const}\NormalTok{ passing }\OperatorTok{=}\NormalTok{ scores}\OperatorTok{.}\FunctionTok{filter}\NormalTok{((s) }\KeywordTok{=\textgreater{}}\NormalTok{ s }\OperatorTok{\textgreater{}=} \DecValTok{80}\NormalTok{)}\OperatorTok{;}

\BuiltInTok{console}\OperatorTok{.}\FunctionTok{log}\NormalTok{(}\VerbatimStringTok{\textasciigrave{}Scores: }\SpecialCharTok{$\{}\NormalTok{scores}\OperatorTok{.}\FunctionTok{join}\NormalTok{(}\StringTok{", "}\NormalTok{)}\SpecialCharTok{\}}\VerbatimStringTok{\textasciigrave{}}\NormalTok{)}\OperatorTok{;}
\BuiltInTok{console}\OperatorTok{.}\FunctionTok{log}\NormalTok{(}\VerbatimStringTok{\textasciigrave{}Average: }\SpecialCharTok{$\{}\NormalTok{average}\OperatorTok{.}\FunctionTok{toFixed}\NormalTok{(}\DecValTok{1}\NormalTok{)}\SpecialCharTok{\}}\VerbatimStringTok{\textasciigrave{}}\NormalTok{)}\OperatorTok{;}
\BuiltInTok{console}\OperatorTok{.}\FunctionTok{log}\NormalTok{(}\VerbatimStringTok{\textasciigrave{}Highest: }\SpecialCharTok{$\{}\NormalTok{max}\SpecialCharTok{\}}\VerbatimStringTok{\textasciigrave{}}\NormalTok{)}\OperatorTok{;}
\BuiltInTok{console}\OperatorTok{.}\FunctionTok{log}\NormalTok{(}\VerbatimStringTok{\textasciigrave{}Lowest: }\SpecialCharTok{$\{}\NormalTok{min}\SpecialCharTok{\}}\VerbatimStringTok{\textasciigrave{}}\NormalTok{)}\OperatorTok{;}
\BuiltInTok{console}\OperatorTok{.}\FunctionTok{log}\NormalTok{(}\VerbatimStringTok{\textasciigrave{}Passing (\textgreater{}=80): }\SpecialCharTok{$\{}\NormalTok{passing}\OperatorTok{.}\AttributeTok{length}\SpecialCharTok{\}}\VerbatimStringTok{ students\textasciigrave{}}\NormalTok{)}\OperatorTok{;}
\end{Highlighting}
\end{Shaded}

\subsubsection{Skills Checklist}\label{skills-checklist-8}

\begin{itemize}
\tightlist
\item[$\square$]
  I can create, access, and modify arrays.
\item[$\square$]
  I can use \texttt{push}, \texttt{pop}, \texttt{shift},
  \texttt{unshift}, \texttt{splice}, and \texttt{slice}.
\item[$\square$]
  I can use \texttt{map}, \texttt{filter}, and \texttt{reduce}.
\item[$\square$]
  I can sort arrays using a comparator function.
\end{itemize}

\begin{center}\rule{0.5\linewidth}{0.5pt}\end{center}

\subsection{Lesson 4.2: Objects}\label{lesson-4.2-objects-1}

\subsubsection{Creating and Using
Objects}\label{creating-and-using-objects}

Objects in JavaScript are like Python dictionaries and Java Maps
combined with class instances.

\begin{Shaded}
\begin{Highlighting}[]
\CommentTok{// Object literal}
\KeywordTok{const}\NormalTok{ studentProfile }\OperatorTok{=}\NormalTok{ \{}
  \DataTypeTok{name}\OperatorTok{:} \StringTok{"Emily"}\OperatorTok{,}
  \DataTypeTok{age}\OperatorTok{:} \DecValTok{13}\OperatorTok{,}
  \DataTypeTok{grade}\OperatorTok{:} \StringTok{"7th"}\OperatorTok{,}
  \DataTypeTok{favoriteSubject}\OperatorTok{:} \StringTok{"Science"}\OperatorTok{,}
\NormalTok{\}}\OperatorTok{;}

\CommentTok{// Access properties}
\BuiltInTok{console}\OperatorTok{.}\FunctionTok{log}\NormalTok{(studentProfile}\OperatorTok{.}\AttributeTok{name}\NormalTok{)}\OperatorTok{;} \CommentTok{// "Emily" (dot notation)}
\BuiltInTok{console}\OperatorTok{.}\FunctionTok{log}\NormalTok{(studentProfile[}\StringTok{"age"}\NormalTok{])}\OperatorTok{;} \CommentTok{// 13 (bracket notation)}

\CommentTok{// Add/modify}
\NormalTok{studentProfile}\OperatorTok{.}\AttributeTok{city} \OperatorTok{=} \StringTok{"Springfield"}\OperatorTok{;}
\NormalTok{studentProfile}\OperatorTok{.}\AttributeTok{age} \OperatorTok{=} \DecValTok{14}\OperatorTok{;}

\CommentTok{// Delete}
\KeywordTok{delete}\NormalTok{ studentProfile}\OperatorTok{.}\AttributeTok{grade}\OperatorTok{;}

\CommentTok{// Check if property exists}
\BuiltInTok{console}\OperatorTok{.}\FunctionTok{log}\NormalTok{(}\StringTok{"name"} \KeywordTok{in}\NormalTok{ studentProfile)}\OperatorTok{;} \CommentTok{// true}
\BuiltInTok{console}\OperatorTok{.}\FunctionTok{log}\NormalTok{(studentProfile}\OperatorTok{.}\FunctionTok{hasOwnProperty}\NormalTok{(}\StringTok{"city"}\NormalTok{))}\OperatorTok{;} \CommentTok{// true}

\CommentTok{// Get keys, values, entries}
\BuiltInTok{console}\OperatorTok{.}\FunctionTok{log}\NormalTok{(}\BuiltInTok{Object}\OperatorTok{.}\FunctionTok{keys}\NormalTok{(studentProfile))}\OperatorTok{;}
\BuiltInTok{console}\OperatorTok{.}\FunctionTok{log}\NormalTok{(}\BuiltInTok{Object}\OperatorTok{.}\FunctionTok{values}\NormalTok{(studentProfile))}\OperatorTok{;}
\BuiltInTok{console}\OperatorTok{.}\FunctionTok{log}\NormalTok{(}\BuiltInTok{Object}\OperatorTok{.}\FunctionTok{entries}\NormalTok{(studentProfile))}\OperatorTok{;}
\end{Highlighting}
\end{Shaded}

\subsubsection{Iterating Over Objects}\label{iterating-over-objects}

\begin{Shaded}
\begin{Highlighting}[]
\KeywordTok{const}\NormalTok{ inventory }\OperatorTok{=}\NormalTok{ \{ }\DataTypeTok{apples}\OperatorTok{:} \DecValTok{50}\OperatorTok{,} \DataTypeTok{bananas}\OperatorTok{:} \DecValTok{30}\OperatorTok{,} \DataTypeTok{oranges}\OperatorTok{:} \DecValTok{40}\NormalTok{ \}}\OperatorTok{;}

\CommentTok{// for...in}
\ControlFlowTok{for}\NormalTok{ (}\KeywordTok{const}\NormalTok{ key }\KeywordTok{in}\NormalTok{ inventory) \{}
  \BuiltInTok{console}\OperatorTok{.}\FunctionTok{log}\NormalTok{(}\VerbatimStringTok{\textasciigrave{}}\SpecialCharTok{$\{}\NormalTok{key}\SpecialCharTok{\}}\VerbatimStringTok{: }\SpecialCharTok{$\{}\NormalTok{inventory[key]}\SpecialCharTok{\}}\VerbatimStringTok{\textasciigrave{}}\NormalTok{)}\OperatorTok{;}
\NormalTok{\}}

\CommentTok{// Object.entries() with destructuring}
\ControlFlowTok{for}\NormalTok{ (}\KeywordTok{const}\NormalTok{ [key}\OperatorTok{,}\NormalTok{ value] }\KeywordTok{of} \BuiltInTok{Object}\OperatorTok{.}\FunctionTok{entries}\NormalTok{(inventory)) \{}
  \BuiltInTok{console}\OperatorTok{.}\FunctionTok{log}\NormalTok{(}\VerbatimStringTok{\textasciigrave{}}\SpecialCharTok{$\{}\NormalTok{key}\SpecialCharTok{\}}\VerbatimStringTok{: }\SpecialCharTok{$\{}\NormalTok{value}\SpecialCharTok{\}}\VerbatimStringTok{ in stock\textasciigrave{}}\NormalTok{)}\OperatorTok{;}
\NormalTok{\}}
\end{Highlighting}
\end{Shaded}

\subsubsection{Nested Objects and
Arrays}\label{nested-objects-and-arrays}

\begin{Shaded}
\begin{Highlighting}[]
\KeywordTok{const}\NormalTok{ library }\OperatorTok{=}\NormalTok{ \{}
  \DataTypeTok{name}\OperatorTok{:} \StringTok{"City Library"}\OperatorTok{,}
  \DataTypeTok{location}\OperatorTok{:} \StringTok{"Downtown"}\OperatorTok{,}
  \DataTypeTok{books}\OperatorTok{:}\NormalTok{ [}
\NormalTok{    \{ }\DataTypeTok{title}\OperatorTok{:} \StringTok{"Dune"}\OperatorTok{,} \DataTypeTok{author}\OperatorTok{:} \StringTok{"Frank Herbert"}\OperatorTok{,} \DataTypeTok{available}\OperatorTok{:} \KeywordTok{true}\NormalTok{ \}}\OperatorTok{,}
\NormalTok{    \{ }\DataTypeTok{title}\OperatorTok{:} \StringTok{"1984"}\OperatorTok{,} \DataTypeTok{author}\OperatorTok{:} \StringTok{"George Orwell"}\OperatorTok{,} \DataTypeTok{available}\OperatorTok{:} \KeywordTok{false}\NormalTok{ \}}\OperatorTok{,}
\NormalTok{  ]}\OperatorTok{,}
\NormalTok{\}}\OperatorTok{;}

\CommentTok{// Access nested data}
\BuiltInTok{console}\OperatorTok{.}\FunctionTok{log}\NormalTok{(library}\OperatorTok{.}\AttributeTok{books}\NormalTok{[}\DecValTok{0}\NormalTok{]}\OperatorTok{.}\AttributeTok{title}\NormalTok{)}\OperatorTok{;} \CommentTok{// "Dune"}
\BuiltInTok{console}\OperatorTok{.}\FunctionTok{log}\NormalTok{(library}\OperatorTok{.}\AttributeTok{books}\NormalTok{[}\DecValTok{1}\NormalTok{]}\OperatorTok{.}\AttributeTok{author}\NormalTok{)}\OperatorTok{;} \CommentTok{// "George Orwell"}

\CommentTok{// Filter available books}
\KeywordTok{const}\NormalTok{ available }\OperatorTok{=}\NormalTok{ library}\OperatorTok{.}\AttributeTok{books}\OperatorTok{.}\FunctionTok{filter}\NormalTok{((book) }\KeywordTok{=\textgreater{}}\NormalTok{ book}\OperatorTok{.}\AttributeTok{available}\NormalTok{)}\OperatorTok{;}
\BuiltInTok{console}\OperatorTok{.}\FunctionTok{log}\NormalTok{(}\VerbatimStringTok{\textasciigrave{}Available: }\SpecialCharTok{$\{}\NormalTok{available}\OperatorTok{.}\FunctionTok{map}\NormalTok{((b) }\KeywordTok{=\textgreater{}}\NormalTok{ b}\OperatorTok{.}\AttributeTok{title}\NormalTok{)}\OperatorTok{.}\FunctionTok{join}\NormalTok{(}\StringTok{", "}\NormalTok{)}\SpecialCharTok{\}}\VerbatimStringTok{\textasciigrave{}}\NormalTok{)}\OperatorTok{;}
\end{Highlighting}
\end{Shaded}

\subsubsection{Destructuring}\label{destructuring}

\begin{Shaded}
\begin{Highlighting}[]
\CommentTok{// Object destructuring}
\KeywordTok{const}\NormalTok{ \{ name}\OperatorTok{,}\NormalTok{ age \} }\OperatorTok{=}\NormalTok{ studentProfile}\OperatorTok{;}
\BuiltInTok{console}\OperatorTok{.}\FunctionTok{log}\NormalTok{(name}\OperatorTok{,}\NormalTok{ age)}\OperatorTok{;} \CommentTok{// "Emily" 13}

\CommentTok{// With renaming}
\KeywordTok{const}\NormalTok{ \{ }\DataTypeTok{name}\OperatorTok{:}\NormalTok{ studentName}\OperatorTok{,} \DataTypeTok{favoriteSubject}\OperatorTok{:}\NormalTok{ subject \} }\OperatorTok{=}\NormalTok{ studentProfile}\OperatorTok{;}

\CommentTok{// Array destructuring}
\KeywordTok{const}\NormalTok{ [first}\OperatorTok{,}\NormalTok{ second}\OperatorTok{,} \OperatorTok{...}\NormalTok{rest] }\OperatorTok{=}\NormalTok{ [}\DecValTok{1}\OperatorTok{,} \DecValTok{2}\OperatorTok{,} \DecValTok{3}\OperatorTok{,} \DecValTok{4}\OperatorTok{,} \DecValTok{5}\NormalTok{]}\OperatorTok{;}
\BuiltInTok{console}\OperatorTok{.}\FunctionTok{log}\NormalTok{(first)}\OperatorTok{;} \CommentTok{// 1}
\BuiltInTok{console}\OperatorTok{.}\FunctionTok{log}\NormalTok{(rest)}\OperatorTok{;} \CommentTok{// [3, 4, 5]}

\CommentTok{// In function parameters}
\KeywordTok{function} \FunctionTok{printStudent}\NormalTok{(\{ name}\OperatorTok{,}\NormalTok{ age}\OperatorTok{,}\NormalTok{ grade \}) \{}
  \BuiltInTok{console}\OperatorTok{.}\FunctionTok{log}\NormalTok{(}\VerbatimStringTok{\textasciigrave{}}\SpecialCharTok{$\{}\NormalTok{name}\SpecialCharTok{\}}\VerbatimStringTok{ | Age: }\SpecialCharTok{$\{}\NormalTok{age}\SpecialCharTok{\}}\VerbatimStringTok{ | Grade: }\SpecialCharTok{$\{}\NormalTok{grade}\SpecialCharTok{\}}\VerbatimStringTok{\textasciigrave{}}\NormalTok{)}\OperatorTok{;}
\NormalTok{\}}
\FunctionTok{printStudent}\NormalTok{(studentProfile)}\OperatorTok{;}
\end{Highlighting}
\end{Shaded}

\subsubsection{Spread and Rest
Operators}\label{spread-and-rest-operators}

\begin{Shaded}
\begin{Highlighting}[]
\CommentTok{// Spread {-}{-}{-} expand array or object}
\KeywordTok{const}\NormalTok{ arr1 }\OperatorTok{=}\NormalTok{ [}\DecValTok{1}\OperatorTok{,} \DecValTok{2}\OperatorTok{,} \DecValTok{3}\NormalTok{]}\OperatorTok{;}
\KeywordTok{const}\NormalTok{ arr2 }\OperatorTok{=}\NormalTok{ [}\DecValTok{4}\OperatorTok{,} \DecValTok{5}\OperatorTok{,} \DecValTok{6}\NormalTok{]}\OperatorTok{;}
\KeywordTok{const}\NormalTok{ combined }\OperatorTok{=}\NormalTok{ [}\OperatorTok{...}\NormalTok{arr1}\OperatorTok{,} \OperatorTok{...}\NormalTok{arr2]}\OperatorTok{;} \CommentTok{// [1, 2, 3, 4, 5, 6]}

\KeywordTok{const}\NormalTok{ obj1 }\OperatorTok{=}\NormalTok{ \{ }\DataTypeTok{a}\OperatorTok{:} \DecValTok{1}\OperatorTok{,} \DataTypeTok{b}\OperatorTok{:} \DecValTok{2}\NormalTok{ \}}\OperatorTok{;}
\KeywordTok{const}\NormalTok{ obj2 }\OperatorTok{=}\NormalTok{ \{ }\DataTypeTok{c}\OperatorTok{:} \DecValTok{3}\OperatorTok{,} \DataTypeTok{d}\OperatorTok{:} \DecValTok{4}\NormalTok{ \}}\OperatorTok{;}
\KeywordTok{const}\NormalTok{ merged }\OperatorTok{=}\NormalTok{ \{ }\OperatorTok{...}\NormalTok{obj1}\OperatorTok{,} \OperatorTok{...}\NormalTok{obj2 \}}\OperatorTok{;} \CommentTok{// \{ a: 1, b: 2, c: 3, d: 4 \}}

\CommentTok{// Copy without mutating}
\KeywordTok{const}\NormalTok{ originalArr }\OperatorTok{=}\NormalTok{ [}\DecValTok{1}\OperatorTok{,} \DecValTok{2}\OperatorTok{,} \DecValTok{3}\NormalTok{]}\OperatorTok{;}
\KeywordTok{const}\NormalTok{ copyArr }\OperatorTok{=}\NormalTok{ [}\OperatorTok{...}\NormalTok{originalArr]}\OperatorTok{;}

\KeywordTok{const}\NormalTok{ originalObj }\OperatorTok{=}\NormalTok{ \{ }\DataTypeTok{name}\OperatorTok{:} \StringTok{"Alice"}\NormalTok{ \}}\OperatorTok{;}
\KeywordTok{const}\NormalTok{ copyObj }\OperatorTok{=}\NormalTok{ \{ }\OperatorTok{...}\NormalTok{originalObj \}}\OperatorTok{;}
\end{Highlighting}
\end{Shaded}

\subsubsection{Skills Checklist}\label{skills-checklist-9}

\begin{itemize}
\tightlist
\item[$\square$]
  I can create, access, modify, and delete object properties.
\item[$\square$]
  I can iterate over objects using \texttt{for...in} and
  \texttt{Object.entries()}.
\item[$\square$]
  I can work with nested objects and arrays.
\item[$\square$]
  I can use destructuring assignment.
\item[$\square$]
  I can use spread and rest operators.
\end{itemize}

\newpage

\section{Unit 5: Strings and String
Methods}\label{unit-5-strings-and-string-methods}

\begin{center}\rule{0.5\linewidth}{0.5pt}\end{center}

\subsection{Lesson 5.1: Working with
Strings}\label{lesson-5.1-working-with-strings}

\begin{Shaded}
\begin{Highlighting}[]
\KeywordTok{const}\NormalTok{ name }\OperatorTok{=} \StringTok{"Alice Smith"}\OperatorTok{;}

\CommentTok{// Properties}
\BuiltInTok{console}\OperatorTok{.}\FunctionTok{log}\NormalTok{(name}\OperatorTok{.}\AttributeTok{length}\NormalTok{)}\OperatorTok{;} \CommentTok{// 11}

\CommentTok{// Accessing characters}
\BuiltInTok{console}\OperatorTok{.}\FunctionTok{log}\NormalTok{(name[}\DecValTok{0}\NormalTok{])}\OperatorTok{;} \CommentTok{// "A"}
\BuiltInTok{console}\OperatorTok{.}\FunctionTok{log}\NormalTok{(name}\OperatorTok{.}\FunctionTok{charAt}\NormalTok{(}\DecValTok{0}\NormalTok{))}\OperatorTok{;} \CommentTok{// "A"}
\BuiltInTok{console}\OperatorTok{.}\FunctionTok{log}\NormalTok{(name[name}\OperatorTok{.}\AttributeTok{length} \OperatorTok{{-}} \DecValTok{1}\NormalTok{])}\OperatorTok{;} \CommentTok{// "h"}

\CommentTok{// Substrings}
\BuiltInTok{console}\OperatorTok{.}\FunctionTok{log}\NormalTok{(name}\OperatorTok{.}\FunctionTok{slice}\NormalTok{(}\DecValTok{6}\NormalTok{))}\OperatorTok{;} \CommentTok{// "Smith"}
\BuiltInTok{console}\OperatorTok{.}\FunctionTok{log}\NormalTok{(name}\OperatorTok{.}\FunctionTok{slice}\NormalTok{(}\DecValTok{0}\OperatorTok{,} \DecValTok{5}\NormalTok{))}\OperatorTok{;} \CommentTok{// "Alice"}
\BuiltInTok{console}\OperatorTok{.}\FunctionTok{log}\NormalTok{(name}\OperatorTok{.}\FunctionTok{substring}\NormalTok{(}\DecValTok{6}\OperatorTok{,} \DecValTok{11}\NormalTok{))}\OperatorTok{;} \CommentTok{// "Smith"}

\CommentTok{// Case}
\BuiltInTok{console}\OperatorTok{.}\FunctionTok{log}\NormalTok{(name}\OperatorTok{.}\FunctionTok{toUpperCase}\NormalTok{())}\OperatorTok{;} \CommentTok{// "ALICE SMITH"}
\BuiltInTok{console}\OperatorTok{.}\FunctionTok{log}\NormalTok{(name}\OperatorTok{.}\FunctionTok{toLowerCase}\NormalTok{())}\OperatorTok{;} \CommentTok{// "alice smith"}

\CommentTok{// Trimming}
\KeywordTok{const}\NormalTok{ padded }\OperatorTok{=} \StringTok{"  hello  "}\OperatorTok{;}
\BuiltInTok{console}\OperatorTok{.}\FunctionTok{log}\NormalTok{(padded}\OperatorTok{.}\FunctionTok{trim}\NormalTok{())}\OperatorTok{;} \CommentTok{// "hello"}
\BuiltInTok{console}\OperatorTok{.}\FunctionTok{log}\NormalTok{(padded}\OperatorTok{.}\FunctionTok{trimStart}\NormalTok{())}\OperatorTok{;} \CommentTok{// "hello  "}
\BuiltInTok{console}\OperatorTok{.}\FunctionTok{log}\NormalTok{(padded}\OperatorTok{.}\FunctionTok{trimEnd}\NormalTok{())}\OperatorTok{;} \CommentTok{// "  hello"}

\CommentTok{// Searching}
\BuiltInTok{console}\OperatorTok{.}\FunctionTok{log}\NormalTok{(name}\OperatorTok{.}\FunctionTok{indexOf}\NormalTok{(}\StringTok{"Smith"}\NormalTok{))}\OperatorTok{;} \CommentTok{// 6}
\BuiltInTok{console}\OperatorTok{.}\FunctionTok{log}\NormalTok{(name}\OperatorTok{.}\FunctionTok{lastIndexOf}\NormalTok{(}\StringTok{"l"}\NormalTok{))}\OperatorTok{;} \CommentTok{// 3}
\BuiltInTok{console}\OperatorTok{.}\FunctionTok{log}\NormalTok{(name}\OperatorTok{.}\FunctionTok{includes}\NormalTok{(}\StringTok{"Alice"}\NormalTok{))}\OperatorTok{;} \CommentTok{// true}
\BuiltInTok{console}\OperatorTok{.}\FunctionTok{log}\NormalTok{(name}\OperatorTok{.}\FunctionTok{startsWith}\NormalTok{(}\StringTok{"Ali"}\NormalTok{))}\OperatorTok{;} \CommentTok{// true}
\BuiltInTok{console}\OperatorTok{.}\FunctionTok{log}\NormalTok{(name}\OperatorTok{.}\FunctionTok{endsWith}\NormalTok{(}\StringTok{"ith"}\NormalTok{))}\OperatorTok{;} \CommentTok{// true}

\CommentTok{// Replacing}
\KeywordTok{const}\NormalTok{ result }\OperatorTok{=}\NormalTok{ name}\OperatorTok{.}\FunctionTok{replace}\NormalTok{(}\StringTok{"Alice"}\OperatorTok{,} \StringTok{"Bob"}\NormalTok{)}\OperatorTok{;} \CommentTok{// "Bob Smith"}
\KeywordTok{const}\NormalTok{ allReplaced }\OperatorTok{=} \StringTok{"a,b,a,c,a"}\OperatorTok{.}\FunctionTok{replaceAll}\NormalTok{(}\StringTok{"a"}\OperatorTok{,} \StringTok{"x"}\NormalTok{)}\OperatorTok{;} \CommentTok{// "x,b,x,c,x"}

\CommentTok{// Splitting}
\KeywordTok{const}\NormalTok{ sentence }\OperatorTok{=} \StringTok{"the quick brown fox"}\OperatorTok{;}
\KeywordTok{const}\NormalTok{ words }\OperatorTok{=}\NormalTok{ sentence}\OperatorTok{.}\FunctionTok{split}\NormalTok{(}\StringTok{" "}\NormalTok{)}\OperatorTok{;}
\BuiltInTok{console}\OperatorTok{.}\FunctionTok{log}\NormalTok{(words)}\OperatorTok{;} \CommentTok{// ["the", "quick", "brown", "fox"]}

\KeywordTok{const}\NormalTok{ csv }\OperatorTok{=} \StringTok{"Alice,15,Portland"}\OperatorTok{;}
\KeywordTok{const}\NormalTok{ parts }\OperatorTok{=}\NormalTok{ csv}\OperatorTok{.}\FunctionTok{split}\NormalTok{(}\StringTok{","}\NormalTok{)}\OperatorTok{;}
\BuiltInTok{console}\OperatorTok{.}\FunctionTok{log}\NormalTok{(parts)}\OperatorTok{;} \CommentTok{// ["Alice", "15", "Portland"]}

\CommentTok{// Repeating and Padding}
\BuiltInTok{console}\OperatorTok{.}\FunctionTok{log}\NormalTok{(}\StringTok{"ha"}\OperatorTok{.}\FunctionTok{repeat}\NormalTok{(}\DecValTok{3}\NormalTok{))}\OperatorTok{;} \CommentTok{// "hahaha"}
\BuiltInTok{console}\OperatorTok{.}\FunctionTok{log}\NormalTok{(}\StringTok{"5"}\OperatorTok{.}\FunctionTok{padStart}\NormalTok{(}\DecValTok{3}\OperatorTok{,} \StringTok{"0"}\NormalTok{))}\OperatorTok{;} \CommentTok{// "005"}
\BuiltInTok{console}\OperatorTok{.}\FunctionTok{log}\NormalTok{(}\StringTok{"5"}\OperatorTok{.}\FunctionTok{padEnd}\NormalTok{(}\DecValTok{3}\OperatorTok{,} \StringTok{"0"}\NormalTok{))}\OperatorTok{;} \CommentTok{// "500"}

\CommentTok{// Template literals (multi{-}line)}
\KeywordTok{const}\NormalTok{ multiLine }\OperatorTok{=} \VerbatimStringTok{\textasciigrave{}}
\VerbatimStringTok{Name: }\SpecialCharTok{$\{}\NormalTok{name}\SpecialCharTok{\}}
\VerbatimStringTok{Length: }\SpecialCharTok{$\{}\NormalTok{name}\OperatorTok{.}\AttributeTok{length}\SpecialCharTok{\}}
\VerbatimStringTok{Uppercase: }\SpecialCharTok{$\{}\NormalTok{name}\OperatorTok{.}\FunctionTok{toUpperCase}\NormalTok{()}\SpecialCharTok{\}}
\VerbatimStringTok{\textasciigrave{}}\OperatorTok{;}
\BuiltInTok{console}\OperatorTok{.}\FunctionTok{log}\NormalTok{(multiLine)}\OperatorTok{;}
\end{Highlighting}
\end{Shaded}

\subsubsection{Regular Expressions
(Basics)}\label{regular-expressions-basics}

\begin{Shaded}
\begin{Highlighting}[]
\KeywordTok{const}\NormalTok{ text }\OperatorTok{=} \StringTok{"Hello World! Hello JavaScript!"}\OperatorTok{;}

\CommentTok{// Test if pattern exists}
\BuiltInTok{console}\OperatorTok{.}\FunctionTok{log}\NormalTok{(}\SpecialStringTok{/hello/i}\OperatorTok{.}\FunctionTok{test}\NormalTok{(text))}\OperatorTok{;} \CommentTok{// true (i = case{-}insensitive)}

\CommentTok{// Find matches}
\KeywordTok{const}\NormalTok{ matches }\OperatorTok{=}\NormalTok{ text}\OperatorTok{.}\FunctionTok{match}\NormalTok{(}\SpecialStringTok{/Hello/g}\NormalTok{)}\OperatorTok{;} \CommentTok{// g = global (find all)}
\BuiltInTok{console}\OperatorTok{.}\FunctionTok{log}\NormalTok{(matches)}\OperatorTok{;} \CommentTok{// ["Hello", "Hello"]}

\CommentTok{// Replace with regex}
\KeywordTok{const}\NormalTok{ cleaned }\OperatorTok{=}\NormalTok{ text}\OperatorTok{.}\FunctionTok{replace}\NormalTok{(}\SpecialStringTok{/}\SpecialCharTok{[!]}\SpecialStringTok{/g}\OperatorTok{,} \StringTok{"."}\NormalTok{)}\OperatorTok{;} \CommentTok{// "Hello World. Hello JavaScript."}

\CommentTok{// Split with regex}
\KeywordTok{const}\NormalTok{ words2 }\OperatorTok{=} \StringTok{"one  two   three"}\OperatorTok{.}\FunctionTok{split}\NormalTok{(}\SpecialStringTok{/}\SpecialCharTok{\textbackslash{}s+}\SpecialStringTok{/}\NormalTok{)}\OperatorTok{;}
\BuiltInTok{console}\OperatorTok{.}\FunctionTok{log}\NormalTok{(words2)}\OperatorTok{;} \CommentTok{// ["one", "two", "three"]}
\end{Highlighting}
\end{Shaded}

\subsubsection{Skills Checklist}\label{skills-checklist-10}

\begin{itemize}
\tightlist
\item[$\square$]
  I can use common \texttt{String} methods in JavaScript.
\item[$\square$]
  I can use template literals for multi-line strings and interpolation.
\item[$\square$]
  I can split strings into arrays and join arrays into strings.
\item[$\square$]
  I understand basic regular expressions.
\end{itemize}

\begin{center}\rule{0.5\linewidth}{0.5pt}\end{center}

\section{Unit 6: Functions}\label{unit-6-functions-1}

\begin{center}\rule{0.5\linewidth}{0.5pt}\end{center}

\subsection{Lesson 6.1: Defining
Functions}\label{lesson-6.1-defining-functions}

JavaScript has multiple ways to define functions:

\subsubsection{Function Declaration}\label{function-declaration}

\begin{Shaded}
\begin{Highlighting}[]
\KeywordTok{function} \FunctionTok{greetUser}\NormalTok{() \{}
  \BuiltInTok{console}\OperatorTok{.}\FunctionTok{log}\NormalTok{(}\StringTok{"Hello, JavaScript programmer!"}\NormalTok{)}\OperatorTok{;}
\NormalTok{\}}

\FunctionTok{greetUser}\NormalTok{()}\OperatorTok{;}
\FunctionTok{greetUser}\NormalTok{()}\OperatorTok{;}
\end{Highlighting}
\end{Shaded}

\subsubsection{Function Expression}\label{function-expression}

\begin{Shaded}
\begin{Highlighting}[]
\KeywordTok{const}\NormalTok{ greetByName }\OperatorTok{=} \KeywordTok{function}\NormalTok{ (name) \{}
  \BuiltInTok{console}\OperatorTok{.}\FunctionTok{log}\NormalTok{(}\VerbatimStringTok{\textasciigrave{}Hello, }\SpecialCharTok{$\{}\NormalTok{name}\SpecialCharTok{\}}\VerbatimStringTok{!\textasciigrave{}}\NormalTok{)}\OperatorTok{;}
\NormalTok{\}}\OperatorTok{;}

\FunctionTok{greetByName}\NormalTok{(}\StringTok{"Alice"}\NormalTok{)}\OperatorTok{;}
\end{Highlighting}
\end{Shaded}

\subsubsection{Arrow Functions (ES6) --- Modern Preferred
Style}\label{arrow-functions-es6-modern-preferred-style}

\begin{Shaded}
\begin{Highlighting}[]
\CommentTok{// Full arrow function}
\KeywordTok{const}\NormalTok{ add }\OperatorTok{=}\NormalTok{ (a}\OperatorTok{,}\NormalTok{ b) }\KeywordTok{=\textgreater{}}\NormalTok{ \{}
  \ControlFlowTok{return}\NormalTok{ a }\OperatorTok{+}\NormalTok{ b}\OperatorTok{;}
\NormalTok{\}}\OperatorTok{;}

\CommentTok{// Concise body (implicit return for single expressions)}
\KeywordTok{const}\NormalTok{ addShort }\OperatorTok{=}\NormalTok{ (a}\OperatorTok{,}\NormalTok{ b) }\KeywordTok{=\textgreater{}}\NormalTok{ a }\OperatorTok{+}\NormalTok{ b}\OperatorTok{;}

\CommentTok{// Single parameter (no parentheses needed)}
\KeywordTok{const}\NormalTok{ double }\OperatorTok{=}\NormalTok{ (n) }\KeywordTok{=\textgreater{}}\NormalTok{ n }\OperatorTok{*} \DecValTok{2}\OperatorTok{;}

\CommentTok{// No parameters}
\KeywordTok{const}\NormalTok{ sayHello }\OperatorTok{=}\NormalTok{ () }\KeywordTok{=\textgreater{}} \BuiltInTok{console}\OperatorTok{.}\FunctionTok{log}\NormalTok{(}\StringTok{"Hello!"}\NormalTok{)}\OperatorTok{;}

\BuiltInTok{console}\OperatorTok{.}\FunctionTok{log}\NormalTok{(}\FunctionTok{add}\NormalTok{(}\DecValTok{5}\OperatorTok{,} \DecValTok{3}\NormalTok{))}\OperatorTok{;} \CommentTok{// 8}
\BuiltInTok{console}\OperatorTok{.}\FunctionTok{log}\NormalTok{(}\FunctionTok{addShort}\NormalTok{(}\DecValTok{5}\OperatorTok{,} \DecValTok{3}\NormalTok{))}\OperatorTok{;} \CommentTok{// 8}
\BuiltInTok{console}\OperatorTok{.}\FunctionTok{log}\NormalTok{(}\FunctionTok{double}\NormalTok{(}\DecValTok{7}\NormalTok{))}\OperatorTok{;} \CommentTok{// 14}
\FunctionTok{sayHello}\NormalTok{()}\OperatorTok{;}
\end{Highlighting}
\end{Shaded}

\subsubsection{Default Parameters}\label{default-parameters}

\begin{Shaded}
\begin{Highlighting}[]
\KeywordTok{function} \FunctionTok{makeShirt}\NormalTok{(size}\OperatorTok{,}\NormalTok{ message }\OperatorTok{=} \StringTok{"I love JavaScript"}\NormalTok{) \{}
  \BuiltInTok{console}\OperatorTok{.}\FunctionTok{log}\NormalTok{(}\VerbatimStringTok{\textasciigrave{}Making a }\SpecialCharTok{$\{}\NormalTok{size}\SpecialCharTok{\}}\VerbatimStringTok{{-}sized shirt: \textquotesingle{}}\SpecialCharTok{$\{}\NormalTok{message}\SpecialCharTok{\}}\VerbatimStringTok{\textquotesingle{}\textasciigrave{}}\NormalTok{)}\OperatorTok{;}
\NormalTok{\}}

\FunctionTok{makeShirt}\NormalTok{(}\StringTok{"large"}\NormalTok{)}\OperatorTok{;} \CommentTok{// Uses default}
\FunctionTok{makeShirt}\NormalTok{(}\StringTok{"medium"}\OperatorTok{,} \StringTok{"Hello World"}\NormalTok{)}\OperatorTok{;} \CommentTok{// Custom message}
\end{Highlighting}
\end{Shaded}

\subsubsection{Rest Parameters}\label{rest-parameters}

\begin{Shaded}
\begin{Highlighting}[]
\KeywordTok{function} \FunctionTok{sum}\NormalTok{(}\OperatorTok{...}\NormalTok{numbers) \{}
  \ControlFlowTok{return}\NormalTok{ numbers}\OperatorTok{.}\FunctionTok{reduce}\NormalTok{((total}\OperatorTok{,}\NormalTok{ n) }\KeywordTok{=\textgreater{}}\NormalTok{ total }\OperatorTok{+}\NormalTok{ n}\OperatorTok{,} \DecValTok{0}\NormalTok{)}\OperatorTok{;}
\NormalTok{\}}

\BuiltInTok{console}\OperatorTok{.}\FunctionTok{log}\NormalTok{(}\FunctionTok{sum}\NormalTok{(}\DecValTok{1}\OperatorTok{,} \DecValTok{2}\OperatorTok{,} \DecValTok{3}\NormalTok{))}\OperatorTok{;} \CommentTok{// 6}
\BuiltInTok{console}\OperatorTok{.}\FunctionTok{log}\NormalTok{(}\FunctionTok{sum}\NormalTok{(}\DecValTok{10}\OperatorTok{,} \DecValTok{20}\OperatorTok{,} \DecValTok{30}\OperatorTok{,} \DecValTok{40}\NormalTok{))}\OperatorTok{;} \CommentTok{// 100}
\end{Highlighting}
\end{Shaded}

\subsubsection{Examples}\label{examples-5}

\textbf{Example 6.1.1: Calculator Functions}

\begin{Shaded}
\begin{Highlighting}[]
\KeywordTok{const}\NormalTok{ add }\OperatorTok{=}\NormalTok{ (a}\OperatorTok{,}\NormalTok{ b) }\KeywordTok{=\textgreater{}}\NormalTok{ a }\OperatorTok{+}\NormalTok{ b}\OperatorTok{;}
\KeywordTok{const}\NormalTok{ subtract }\OperatorTok{=}\NormalTok{ (a}\OperatorTok{,}\NormalTok{ b) }\KeywordTok{=\textgreater{}}\NormalTok{ a }\OperatorTok{{-}}\NormalTok{ b}\OperatorTok{;}
\KeywordTok{const}\NormalTok{ multiply }\OperatorTok{=}\NormalTok{ (a}\OperatorTok{,}\NormalTok{ b) }\KeywordTok{=\textgreater{}}\NormalTok{ a }\OperatorTok{*}\NormalTok{ b}\OperatorTok{;}
\KeywordTok{const}\NormalTok{ divide }\OperatorTok{=}\NormalTok{ (a}\OperatorTok{,}\NormalTok{ b) }\KeywordTok{=\textgreater{}}\NormalTok{ \{}
  \ControlFlowTok{if}\NormalTok{ (b }\OperatorTok{===} \DecValTok{0}\NormalTok{) }\ControlFlowTok{throw} \KeywordTok{new} \BuiltInTok{Error}\NormalTok{(}\StringTok{"Cannot divide by zero"}\NormalTok{)}\OperatorTok{;}
  \ControlFlowTok{return}\NormalTok{ a }\OperatorTok{/}\NormalTok{ b}\OperatorTok{;}
\NormalTok{\}}\OperatorTok{;}

\KeywordTok{const}\NormalTok{ circleArea }\OperatorTok{=}\NormalTok{ (radius) }\KeywordTok{=\textgreater{}} \BuiltInTok{Math}\OperatorTok{.}\ConstantTok{PI} \OperatorTok{*}\NormalTok{ radius }\OperatorTok{**} \DecValTok{2}\OperatorTok{;}
\KeywordTok{const}\NormalTok{ getGrade }\OperatorTok{=}\NormalTok{ (score) }\KeywordTok{=\textgreater{}}\NormalTok{ \{}
  \ControlFlowTok{if}\NormalTok{ (score }\OperatorTok{\textgreater{}=} \DecValTok{90}\NormalTok{) }\ControlFlowTok{return} \StringTok{"A"}\OperatorTok{;}
  \ControlFlowTok{if}\NormalTok{ (score }\OperatorTok{\textgreater{}=} \DecValTok{80}\NormalTok{) }\ControlFlowTok{return} \StringTok{"B"}\OperatorTok{;}
  \ControlFlowTok{if}\NormalTok{ (score }\OperatorTok{\textgreater{}=} \DecValTok{70}\NormalTok{) }\ControlFlowTok{return} \StringTok{"C"}\OperatorTok{;}
  \ControlFlowTok{if}\NormalTok{ (score }\OperatorTok{\textgreater{}=} \DecValTok{60}\NormalTok{) }\ControlFlowTok{return} \StringTok{"D"}\OperatorTok{;}
  \ControlFlowTok{return} \StringTok{"F"}\OperatorTok{;}
\NormalTok{\}}\OperatorTok{;}

\BuiltInTok{console}\OperatorTok{.}\FunctionTok{log}\NormalTok{(}\FunctionTok{add}\NormalTok{(}\DecValTok{5}\OperatorTok{,} \DecValTok{3}\NormalTok{))}\OperatorTok{;}
\BuiltInTok{console}\OperatorTok{.}\FunctionTok{log}\NormalTok{(}\FunctionTok{circleArea}\NormalTok{(}\DecValTok{5}\NormalTok{)}\OperatorTok{.}\FunctionTok{toFixed}\NormalTok{(}\DecValTok{2}\NormalTok{))}\OperatorTok{;}
\BuiltInTok{console}\OperatorTok{.}\FunctionTok{log}\NormalTok{(}\FunctionTok{getGrade}\NormalTok{(}\DecValTok{85}\NormalTok{))}\OperatorTok{;}
\end{Highlighting}
\end{Shaded}

\textbf{Example 6.1.2: Higher-Order Functions}

\begin{Shaded}
\begin{Highlighting}[]
\CommentTok{// Functions that take or return other functions}

\CommentTok{// Map using our own function}
\KeywordTok{function} \FunctionTok{doubleAll}\NormalTok{(numbers) \{}
  \ControlFlowTok{return}\NormalTok{ numbers}\OperatorTok{.}\FunctionTok{map}\NormalTok{((n) }\KeywordTok{=\textgreater{}}\NormalTok{ n }\OperatorTok{*} \DecValTok{2}\NormalTok{)}\OperatorTok{;}
\NormalTok{\}}

\CommentTok{// Function that returns a function (closure)}
\KeywordTok{function} \FunctionTok{makeAdder}\NormalTok{(x) \{}
  \ControlFlowTok{return}\NormalTok{ (y) }\KeywordTok{=\textgreater{}}\NormalTok{ x }\OperatorTok{+}\NormalTok{ y}\OperatorTok{;}
\NormalTok{\}}

\KeywordTok{const}\NormalTok{ addFive }\OperatorTok{=} \FunctionTok{makeAdder}\NormalTok{(}\DecValTok{5}\NormalTok{)}\OperatorTok{;}
\KeywordTok{const}\NormalTok{ addTen }\OperatorTok{=} \FunctionTok{makeAdder}\NormalTok{(}\DecValTok{10}\NormalTok{)}\OperatorTok{;}

\BuiltInTok{console}\OperatorTok{.}\FunctionTok{log}\NormalTok{(}\FunctionTok{addFive}\NormalTok{(}\DecValTok{3}\NormalTok{))}\OperatorTok{;} \CommentTok{// 8}
\BuiltInTok{console}\OperatorTok{.}\FunctionTok{log}\NormalTok{(}\FunctionTok{addTen}\NormalTok{(}\DecValTok{3}\NormalTok{))}\OperatorTok{;} \CommentTok{// 13}
\BuiltInTok{console}\OperatorTok{.}\FunctionTok{log}\NormalTok{(}\FunctionTok{doubleAll}\NormalTok{([}\DecValTok{1}\OperatorTok{,} \DecValTok{2}\OperatorTok{,} \DecValTok{3}\OperatorTok{,} \DecValTok{4}\OperatorTok{,} \DecValTok{5}\NormalTok{]))}\OperatorTok{;}
\end{Highlighting}
\end{Shaded}

\subsubsection{Closures}\label{closures}

\begin{Shaded}
\begin{Highlighting}[]
\KeywordTok{function} \FunctionTok{makeCounter}\NormalTok{() \{}
  \KeywordTok{let}\NormalTok{ count }\OperatorTok{=} \DecValTok{0}\OperatorTok{;} \CommentTok{// Private variable}

  \ControlFlowTok{return}\NormalTok{ \{}
    \DataTypeTok{increment}\OperatorTok{:}\NormalTok{ () }\KeywordTok{=\textgreater{}} \OperatorTok{++}\NormalTok{count}\OperatorTok{,}
    \DataTypeTok{decrement}\OperatorTok{:}\NormalTok{ () }\KeywordTok{=\textgreater{}} \OperatorTok{{-}{-}}\NormalTok{count}\OperatorTok{,}
    \DataTypeTok{getCount}\OperatorTok{:}\NormalTok{ () }\KeywordTok{=\textgreater{}}\NormalTok{ count}\OperatorTok{,}
\NormalTok{  \}}\OperatorTok{;}
\NormalTok{\}}

\KeywordTok{const}\NormalTok{ counter }\OperatorTok{=} \FunctionTok{makeCounter}\NormalTok{()}\OperatorTok{;}
\NormalTok{counter}\OperatorTok{.}\FunctionTok{increment}\NormalTok{()}\OperatorTok{;}
\NormalTok{counter}\OperatorTok{.}\FunctionTok{increment}\NormalTok{()}\OperatorTok{;}
\NormalTok{counter}\OperatorTok{.}\FunctionTok{increment}\NormalTok{()}\OperatorTok{;}
\NormalTok{counter}\OperatorTok{.}\FunctionTok{decrement}\NormalTok{()}\OperatorTok{;}
\BuiltInTok{console}\OperatorTok{.}\FunctionTok{log}\NormalTok{(counter}\OperatorTok{.}\FunctionTok{getCount}\NormalTok{())}\OperatorTok{;} \CommentTok{// 2}
\end{Highlighting}
\end{Shaded}

\subsubsection{Skills Checklist}\label{skills-checklist-11}

\begin{itemize}
\tightlist
\item[$\square$]
  I can define functions using declarations, expressions, and arrow
  syntax.
\item[$\square$]
  I can use default and rest parameters.
\item[$\square$]
  I understand closures and how they capture variables.
\item[$\square$]
  I can pass functions as arguments and return them.
\end{itemize}

\begin{center}\rule{0.5\linewidth}{0.5pt}\end{center}

\section{Unit 7: Classes and Object-Oriented
Programming}\label{unit-7-classes-and-object-oriented-programming}

\begin{center}\rule{0.5\linewidth}{0.5pt}\end{center}

\subsection{Lesson 7.1: Classes in
JavaScript}\label{lesson-7.1-classes-in-javascript}

JavaScript uses prototype-based inheritance, but ES6 \texttt{class}
syntax provides a familiar OOP style.

\subsubsection{Class Syntax}\label{class-syntax}

\begin{Shaded}
\begin{Highlighting}[]
\KeywordTok{class}\NormalTok{ Dog \{}
  \CommentTok{// Instance variables declared in constructor}
  \FunctionTok{constructor}\NormalTok{(name}\OperatorTok{,}\NormalTok{ age) \{}
    \KeywordTok{this}\OperatorTok{.}\AttributeTok{name} \OperatorTok{=}\NormalTok{ name}\OperatorTok{;}
    \KeywordTok{this}\OperatorTok{.}\AttributeTok{age} \OperatorTok{=}\NormalTok{ age}\OperatorTok{;}
\NormalTok{  \}}

  \CommentTok{// Methods}
  \FunctionTok{sit}\NormalTok{() \{}
    \BuiltInTok{console}\OperatorTok{.}\FunctionTok{log}\NormalTok{(}\VerbatimStringTok{\textasciigrave{}}\SpecialCharTok{$\{}\KeywordTok{this}\OperatorTok{.}\AttributeTok{name}\SpecialCharTok{\}}\VerbatimStringTok{ is now sitting.\textasciigrave{}}\NormalTok{)}\OperatorTok{;}
\NormalTok{  \}}

  \FunctionTok{rollOver}\NormalTok{() \{}
    \BuiltInTok{console}\OperatorTok{.}\FunctionTok{log}\NormalTok{(}\VerbatimStringTok{\textasciigrave{}}\SpecialCharTok{$\{}\KeywordTok{this}\OperatorTok{.}\AttributeTok{name}\SpecialCharTok{\}}\VerbatimStringTok{ rolled over!\textasciigrave{}}\NormalTok{)}\OperatorTok{;}
\NormalTok{  \}}

  \FunctionTok{getOlder}\NormalTok{(years) \{}
    \KeywordTok{this}\OperatorTok{.}\AttributeTok{age} \OperatorTok{+=}\NormalTok{ years}\OperatorTok{;}
    \BuiltInTok{console}\OperatorTok{.}\FunctionTok{log}\NormalTok{(}\VerbatimStringTok{\textasciigrave{}}\SpecialCharTok{$\{}\KeywordTok{this}\OperatorTok{.}\AttributeTok{name}\SpecialCharTok{\}}\VerbatimStringTok{ is now }\SpecialCharTok{$\{}\KeywordTok{this}\OperatorTok{.}\AttributeTok{age}\SpecialCharTok{\}}\VerbatimStringTok{ years old.\textasciigrave{}}\NormalTok{)}\OperatorTok{;}
\NormalTok{  \}}

  \CommentTok{// Getter}
  \KeywordTok{get} \FunctionTok{info}\NormalTok{() \{}
    \ControlFlowTok{return} \VerbatimStringTok{\textasciigrave{}}\SpecialCharTok{$\{}\KeywordTok{this}\OperatorTok{.}\AttributeTok{name}\SpecialCharTok{\}}\VerbatimStringTok{ (}\SpecialCharTok{$\{}\KeywordTok{this}\OperatorTok{.}\AttributeTok{age}\SpecialCharTok{\}}\VerbatimStringTok{ years old)\textasciigrave{}}\OperatorTok{;}
\NormalTok{  \}}

  \CommentTok{// toString equivalent}
  \FunctionTok{toString}\NormalTok{() \{}
    \ControlFlowTok{return} \VerbatimStringTok{\textasciigrave{}Dog \{ name: \textquotesingle{}}\SpecialCharTok{$\{}\KeywordTok{this}\OperatorTok{.}\AttributeTok{name}\SpecialCharTok{\}}\VerbatimStringTok{\textquotesingle{}, age: }\SpecialCharTok{$\{}\KeywordTok{this}\OperatorTok{.}\AttributeTok{age}\SpecialCharTok{\}}\VerbatimStringTok{ \}\textasciigrave{}}\OperatorTok{;}
\NormalTok{  \}}
\NormalTok{\}}

\CommentTok{// Create instances}
\KeywordTok{const}\NormalTok{ myDog }\OperatorTok{=} \KeywordTok{new} \FunctionTok{Dog}\NormalTok{(}\StringTok{"Buddy"}\OperatorTok{,} \DecValTok{3}\NormalTok{)}\OperatorTok{;}
\KeywordTok{const}\NormalTok{ yourDog }\OperatorTok{=} \KeywordTok{new} \FunctionTok{Dog}\NormalTok{(}\StringTok{"Lucy"}\OperatorTok{,} \DecValTok{1}\NormalTok{)}\OperatorTok{;}

\NormalTok{myDog}\OperatorTok{.}\FunctionTok{sit}\NormalTok{()}\OperatorTok{;}
\NormalTok{myDog}\OperatorTok{.}\FunctionTok{getOlder}\NormalTok{(}\DecValTok{2}\NormalTok{)}\OperatorTok{;}
\BuiltInTok{console}\OperatorTok{.}\FunctionTok{log}\NormalTok{(myDog}\OperatorTok{.}\AttributeTok{info}\NormalTok{)}\OperatorTok{;} \CommentTok{// Uses getter}
\BuiltInTok{console}\OperatorTok{.}\FunctionTok{log}\NormalTok{(myDog}\OperatorTok{.}\FunctionTok{toString}\NormalTok{())}\OperatorTok{;}
\end{Highlighting}
\end{Shaded}

\subsubsection{Private Fields (ES2022)}\label{private-fields-es2022}

\begin{Shaded}
\begin{Highlighting}[]
\KeywordTok{class}\NormalTok{ BankAccount \{}
\NormalTok{  \#balance}\OperatorTok{;} \CommentTok{// Private field {-}{-}{-} only accessible within the class}

  \FunctionTok{constructor}\NormalTok{(initialBalance) \{}
    \KeywordTok{this}\OperatorTok{.}\NormalTok{\#}\AttributeTok{balance} \OperatorTok{=}\NormalTok{ initialBalance}\OperatorTok{;}
\NormalTok{  \}}

  \FunctionTok{deposit}\NormalTok{(amount) \{}
    \ControlFlowTok{if}\NormalTok{ (amount }\OperatorTok{\textgreater{}} \DecValTok{0}\NormalTok{) \{}
      \KeywordTok{this}\OperatorTok{.}\NormalTok{\#}\AttributeTok{balance} \OperatorTok{+=}\NormalTok{ amount}\OperatorTok{;}
      \BuiltInTok{console}\OperatorTok{.}\FunctionTok{log}\NormalTok{(}\VerbatimStringTok{\textasciigrave{}Deposited $}\SpecialCharTok{$\{}\NormalTok{amount}\SpecialCharTok{\}}\VerbatimStringTok{. New balance: $}\SpecialCharTok{$\{}\KeywordTok{this}\OperatorTok{.}\NormalTok{\#}\AttributeTok{balance}\SpecialCharTok{\}}\VerbatimStringTok{\textasciigrave{}}\NormalTok{)}\OperatorTok{;}
\NormalTok{    \}}
\NormalTok{  \}}

  \FunctionTok{withdraw}\NormalTok{(amount) \{}
    \ControlFlowTok{if}\NormalTok{ (amount }\OperatorTok{\textgreater{}} \KeywordTok{this}\OperatorTok{.}\NormalTok{\#}\AttributeTok{balance}\NormalTok{) \{}
      \BuiltInTok{console}\OperatorTok{.}\FunctionTok{log}\NormalTok{(}\StringTok{"Insufficient funds!"}\NormalTok{)}\OperatorTok{;}
\NormalTok{    \} }\ControlFlowTok{else}\NormalTok{ \{}
      \KeywordTok{this}\OperatorTok{.}\NormalTok{\#}\AttributeTok{balance} \OperatorTok{{-}=}\NormalTok{ amount}\OperatorTok{;}
      \BuiltInTok{console}\OperatorTok{.}\FunctionTok{log}\NormalTok{(}\VerbatimStringTok{\textasciigrave{}Withdrew $}\SpecialCharTok{$\{}\NormalTok{amount}\SpecialCharTok{\}}\VerbatimStringTok{. New balance: $}\SpecialCharTok{$\{}\KeywordTok{this}\OperatorTok{.}\NormalTok{\#}\AttributeTok{balance}\SpecialCharTok{\}}\VerbatimStringTok{\textasciigrave{}}\NormalTok{)}\OperatorTok{;}
\NormalTok{    \}}
\NormalTok{  \}}

  \KeywordTok{get} \FunctionTok{balance}\NormalTok{() \{}
    \ControlFlowTok{return} \KeywordTok{this}\OperatorTok{.}\NormalTok{\#}\AttributeTok{balance}\OperatorTok{;}
\NormalTok{  \}}
\NormalTok{\}}

\KeywordTok{const}\NormalTok{ account }\OperatorTok{=} \KeywordTok{new} \FunctionTok{BankAccount}\NormalTok{(}\DecValTok{1000}\NormalTok{)}\OperatorTok{;}
\NormalTok{account}\OperatorTok{.}\FunctionTok{deposit}\NormalTok{(}\DecValTok{500}\NormalTok{)}\OperatorTok{;}
\NormalTok{account}\OperatorTok{.}\FunctionTok{withdraw}\NormalTok{(}\DecValTok{200}\NormalTok{)}\OperatorTok{;}
\BuiltInTok{console}\OperatorTok{.}\FunctionTok{log}\NormalTok{(}\VerbatimStringTok{\textasciigrave{}Balance: $}\SpecialCharTok{$\{}\NormalTok{account}\OperatorTok{.}\AttributeTok{balance}\SpecialCharTok{\}}\VerbatimStringTok{\textasciigrave{}}\NormalTok{)}\OperatorTok{;}
\end{Highlighting}
\end{Shaded}

\subsubsection{Static Methods}\label{static-methods}

\begin{Shaded}
\begin{Highlighting}[]
\KeywordTok{class}\NormalTok{ MathHelper \{}
  \KeywordTok{static} \FunctionTok{add}\NormalTok{(a}\OperatorTok{,}\NormalTok{ b) \{}
    \ControlFlowTok{return}\NormalTok{ a }\OperatorTok{+}\NormalTok{ b}\OperatorTok{;}
\NormalTok{  \}}
  \KeywordTok{static} \FunctionTok{multiply}\NormalTok{(a}\OperatorTok{,}\NormalTok{ b) \{}
    \ControlFlowTok{return}\NormalTok{ a }\OperatorTok{*}\NormalTok{ b}\OperatorTok{;}
\NormalTok{  \}}
  \KeywordTok{static} \ConstantTok{PI} \OperatorTok{=} \FloatTok{3.14159265358979}\OperatorTok{;}
\NormalTok{\}}

\CommentTok{// Call without creating an instance}
\BuiltInTok{console}\OperatorTok{.}\FunctionTok{log}\NormalTok{(MathHelper}\OperatorTok{.}\FunctionTok{add}\NormalTok{(}\DecValTok{5}\OperatorTok{,} \DecValTok{3}\NormalTok{))}\OperatorTok{;}
\BuiltInTok{console}\OperatorTok{.}\FunctionTok{log}\NormalTok{(MathHelper}\OperatorTok{.}\AttributeTok{PI}\NormalTok{)}\OperatorTok{;}
\end{Highlighting}
\end{Shaded}

\begin{center}\rule{0.5\linewidth}{0.5pt}\end{center}

\subsection{Lesson 7.2: Inheritance}\label{lesson-7.2-inheritance}

\begin{Shaded}
\begin{Highlighting}[]
\KeywordTok{class}\NormalTok{ Vehicle \{}
  \FunctionTok{constructor}\NormalTok{(make}\OperatorTok{,}\NormalTok{ model}\OperatorTok{,}\NormalTok{ year) \{}
    \KeywordTok{this}\OperatorTok{.}\AttributeTok{make} \OperatorTok{=}\NormalTok{ make}\OperatorTok{;}
    \KeywordTok{this}\OperatorTok{.}\AttributeTok{model} \OperatorTok{=}\NormalTok{ model}\OperatorTok{;}
    \KeywordTok{this}\OperatorTok{.}\AttributeTok{year} \OperatorTok{=}\NormalTok{ year}\OperatorTok{;}
    \KeywordTok{this}\OperatorTok{.}\AttributeTok{odometer} \OperatorTok{=} \DecValTok{0}\OperatorTok{;}
\NormalTok{  \}}

  \FunctionTok{getDescriptiveName}\NormalTok{() \{}
    \ControlFlowTok{return} \VerbatimStringTok{\textasciigrave{}}\SpecialCharTok{$\{}\KeywordTok{this}\OperatorTok{.}\AttributeTok{year}\SpecialCharTok{\}}\VerbatimStringTok{ }\SpecialCharTok{$\{}\KeywordTok{this}\OperatorTok{.}\AttributeTok{make}\SpecialCharTok{\}}\VerbatimStringTok{ }\SpecialCharTok{$\{}\KeywordTok{this}\OperatorTok{.}\AttributeTok{model}\SpecialCharTok{\}}\VerbatimStringTok{\textasciigrave{}}\OperatorTok{;}
\NormalTok{  \}}

  \FunctionTok{readOdometer}\NormalTok{() \{}
    \BuiltInTok{console}\OperatorTok{.}\FunctionTok{log}\NormalTok{(}\VerbatimStringTok{\textasciigrave{}This vehicle has }\SpecialCharTok{$\{}\KeywordTok{this}\OperatorTok{.}\AttributeTok{odometer}\SpecialCharTok{\}}\VerbatimStringTok{ miles on it.\textasciigrave{}}\NormalTok{)}\OperatorTok{;}
\NormalTok{  \}}
\NormalTok{\}}

\KeywordTok{class}\NormalTok{ Car }\KeywordTok{extends}\NormalTok{ Vehicle \{}
  \FunctionTok{constructor}\NormalTok{(make}\OperatorTok{,}\NormalTok{ model}\OperatorTok{,}\NormalTok{ year) \{}
    \KeywordTok{super}\NormalTok{(make}\OperatorTok{,}\NormalTok{ model}\OperatorTok{,}\NormalTok{ year)}\OperatorTok{;} \CommentTok{// Call parent constructor}
    \KeywordTok{this}\OperatorTok{.}\AttributeTok{numDoors} \OperatorTok{=} \DecValTok{4}\OperatorTok{;}
\NormalTok{  \}}

  \FunctionTok{readOdometer}\NormalTok{() \{}
    \CommentTok{// Override parent method}
    \BuiltInTok{console}\OperatorTok{.}\FunctionTok{log}\NormalTok{(}\VerbatimStringTok{\textasciigrave{}This car has }\SpecialCharTok{$\{}\KeywordTok{this}\OperatorTok{.}\AttributeTok{odometer}\SpecialCharTok{\}}\VerbatimStringTok{ miles on it.\textasciigrave{}}\NormalTok{)}\OperatorTok{;}
\NormalTok{  \}}
\NormalTok{\}}

\KeywordTok{class}\NormalTok{ ElectricCar }\KeywordTok{extends}\NormalTok{ Car \{}
\NormalTok{  \#batterySize}\OperatorTok{;}

  \FunctionTok{constructor}\NormalTok{(make}\OperatorTok{,}\NormalTok{ model}\OperatorTok{,}\NormalTok{ year) \{}
    \KeywordTok{super}\NormalTok{(make}\OperatorTok{,}\NormalTok{ model}\OperatorTok{,}\NormalTok{ year)}\OperatorTok{;}
    \KeywordTok{this}\OperatorTok{.}\NormalTok{\#}\AttributeTok{batterySize} \OperatorTok{=} \DecValTok{75}\OperatorTok{;}
\NormalTok{  \}}

  \FunctionTok{describeBattery}\NormalTok{() \{}
    \BuiltInTok{console}\OperatorTok{.}\FunctionTok{log}\NormalTok{(}\VerbatimStringTok{\textasciigrave{}This car has a }\SpecialCharTok{$\{}\KeywordTok{this}\OperatorTok{.}\NormalTok{\#}\AttributeTok{batterySize}\SpecialCharTok{\}}\VerbatimStringTok{{-}kWh battery.\textasciigrave{}}\NormalTok{)}\OperatorTok{;}
\NormalTok{  \}}
\NormalTok{\}}

\KeywordTok{const}\NormalTok{ myEV }\OperatorTok{=} \KeywordTok{new} \FunctionTok{ElectricCar}\NormalTok{(}\StringTok{"Ford"}\OperatorTok{,} \StringTok{"Mustang Mach{-}E"}\OperatorTok{,} \DecValTok{2023}\NormalTok{)}\OperatorTok{;}
\BuiltInTok{console}\OperatorTok{.}\FunctionTok{log}\NormalTok{(myEV}\OperatorTok{.}\FunctionTok{getDescriptiveName}\NormalTok{())}\OperatorTok{;} \CommentTok{// Inherited from Vehicle}
\NormalTok{myEV}\OperatorTok{.}\FunctionTok{readOdometer}\NormalTok{()}\OperatorTok{;} \CommentTok{// Overridden in Car}
\NormalTok{myEV}\OperatorTok{.}\FunctionTok{describeBattery}\NormalTok{()}\OperatorTok{;} \CommentTok{// Defined in ElectricCar}
\end{Highlighting}
\end{Shaded}

\subsubsection{Skills Checklist}\label{skills-checklist-12}

\begin{itemize}
\tightlist
\item[$\square$]
  I can define a JavaScript class with a constructor and methods.
\item[$\square$]
  I can use \texttt{extends} to create a subclass.
\item[$\square$]
  I can use \texttt{super()} to call the parent constructor.
\item[$\square$]
  I can use \texttt{\#} for private fields.
\item[$\square$]
  I can define static methods and properties.
\end{itemize}

\begin{center}\rule{0.5\linewidth}{0.5pt}\end{center}

\section{Unit 8: Files, Modules, and Error
Handling}\label{unit-8-files-modules-and-error-handling-1}

\begin{center}\rule{0.5\linewidth}{0.5pt}\end{center}

\subsection{Lesson 8.1: Reading and Writing
Files}\label{lesson-8.1-reading-and-writing-files}

\begin{Shaded}
\begin{Highlighting}[]
\KeywordTok{const}\NormalTok{ fs }\OperatorTok{=} \PreprocessorTok{require}\NormalTok{(}\StringTok{"fs"}\NormalTok{)}\OperatorTok{;}
\KeywordTok{const}\NormalTok{ path }\OperatorTok{=} \PreprocessorTok{require}\NormalTok{(}\StringTok{"path"}\NormalTok{)}\OperatorTok{;}

\CommentTok{// Synchronous reading (simple but blocks execution)}
\ControlFlowTok{try}\NormalTok{ \{}
  \KeywordTok{const}\NormalTok{ content }\OperatorTok{=}\NormalTok{ fs}\OperatorTok{.}\FunctionTok{readFileSync}\NormalTok{(}\StringTok{"my\_notes.txt"}\OperatorTok{,} \StringTok{"utf8"}\NormalTok{)}\OperatorTok{;}
  \BuiltInTok{console}\OperatorTok{.}\FunctionTok{log}\NormalTok{(content)}\OperatorTok{;}
\NormalTok{\} }\ControlFlowTok{catch}\NormalTok{ (err) \{}
  \BuiltInTok{console}\OperatorTok{.}\FunctionTok{log}\NormalTok{(}\VerbatimStringTok{\textasciigrave{}Error: }\SpecialCharTok{$\{}\NormalTok{err}\OperatorTok{.}\AttributeTok{message}\SpecialCharTok{\}}\VerbatimStringTok{\textasciigrave{}}\NormalTok{)}\OperatorTok{;}
\NormalTok{\}}

\CommentTok{// Reading line by line}
\KeywordTok{const}\NormalTok{ lines }\OperatorTok{=}\NormalTok{ content}\OperatorTok{.}\FunctionTok{split}\NormalTok{(}\StringTok{"}\SpecialCharTok{\textbackslash{}n}\StringTok{"}\NormalTok{)}\OperatorTok{;}
\NormalTok{lines}\OperatorTok{.}\FunctionTok{forEach}\NormalTok{((line}\OperatorTok{,}\NormalTok{ index) }\KeywordTok{=\textgreater{}}\NormalTok{ \{}
  \BuiltInTok{console}\OperatorTok{.}\FunctionTok{log}\NormalTok{(}\VerbatimStringTok{\textasciigrave{}Line }\SpecialCharTok{$\{}\NormalTok{index }\OperatorTok{+} \DecValTok{1}\SpecialCharTok{\}}\VerbatimStringTok{: }\SpecialCharTok{$\{}\NormalTok{line}\SpecialCharTok{\}}\VerbatimStringTok{\textasciigrave{}}\NormalTok{)}\OperatorTok{;}
\NormalTok{\})}\OperatorTok{;}

\CommentTok{// Writing (overwrites existing content)}
\ControlFlowTok{try}\NormalTok{ \{}
\NormalTok{  fs}\OperatorTok{.}\FunctionTok{writeFileSync}\NormalTok{(}\StringTok{"output.txt"}\OperatorTok{,} \StringTok{"Hello, Node.js file I/O!}\SpecialCharTok{\textbackslash{}n}\StringTok{"}\NormalTok{)}\OperatorTok{;}
  \BuiltInTok{console}\OperatorTok{.}\FunctionTok{log}\NormalTok{(}\StringTok{"File written successfully."}\NormalTok{)}\OperatorTok{;}
\NormalTok{\} }\ControlFlowTok{catch}\NormalTok{ (err) \{}
  \BuiltInTok{console}\OperatorTok{.}\FunctionTok{log}\NormalTok{(}\VerbatimStringTok{\textasciigrave{}Error writing file: }\SpecialCharTok{$\{}\NormalTok{err}\OperatorTok{.}\AttributeTok{message}\SpecialCharTok{\}}\VerbatimStringTok{\textasciigrave{}}\NormalTok{)}\OperatorTok{;}
\NormalTok{\}}

\CommentTok{// Appending}
\ControlFlowTok{try}\NormalTok{ \{}
\NormalTok{  fs}\OperatorTok{.}\FunctionTok{appendFileSync}\NormalTok{(}\StringTok{"output.txt"}\OperatorTok{,} \StringTok{"This line is appended.}\SpecialCharTok{\textbackslash{}n}\StringTok{"}\NormalTok{)}\OperatorTok{;}
  \BuiltInTok{console}\OperatorTok{.}\FunctionTok{log}\NormalTok{(}\StringTok{"Line appended."}\NormalTok{)}\OperatorTok{;}
\NormalTok{\} }\ControlFlowTok{catch}\NormalTok{ (err) \{}
  \BuiltInTok{console}\OperatorTok{.}\FunctionTok{log}\NormalTok{(}\VerbatimStringTok{\textasciigrave{}Error appending: }\SpecialCharTok{$\{}\NormalTok{err}\OperatorTok{.}\AttributeTok{message}\SpecialCharTok{\}}\VerbatimStringTok{\textasciigrave{}}\NormalTok{)}\OperatorTok{;}
\NormalTok{\}}

\CommentTok{// Check if file exists}
\ControlFlowTok{if}\NormalTok{ (fs}\OperatorTok{.}\FunctionTok{existsSync}\NormalTok{(}\StringTok{"my\_notes.txt"}\NormalTok{)) \{}
  \BuiltInTok{console}\OperatorTok{.}\FunctionTok{log}\NormalTok{(}\StringTok{"File exists."}\NormalTok{)}\OperatorTok{;}
\NormalTok{\}}
\end{Highlighting}
\end{Shaded}

\subsubsection{Asynchronous File I/O (Modern
Approach)}\label{asynchronous-file-io-modern-approach}

\begin{Shaded}
\begin{Highlighting}[]
\KeywordTok{const}\NormalTok{ fs }\OperatorTok{=} \PreprocessorTok{require}\NormalTok{(}\StringTok{"fs"}\NormalTok{)}\OperatorTok{.}\AttributeTok{promises}\OperatorTok{;} \CommentTok{// Promise{-}based fs}

\KeywordTok{async} \KeywordTok{function} \FunctionTok{readFile}\NormalTok{() \{}
  \ControlFlowTok{try}\NormalTok{ \{}
    \KeywordTok{const}\NormalTok{ content }\OperatorTok{=} \ControlFlowTok{await}\NormalTok{ fs}\OperatorTok{.}\FunctionTok{readFile}\NormalTok{(}\StringTok{"my\_notes.txt"}\OperatorTok{,} \StringTok{"utf8"}\NormalTok{)}\OperatorTok{;}
    \BuiltInTok{console}\OperatorTok{.}\FunctionTok{log}\NormalTok{(content)}\OperatorTok{;}
\NormalTok{  \} }\ControlFlowTok{catch}\NormalTok{ (err) \{}
    \BuiltInTok{console}\OperatorTok{.}\FunctionTok{log}\NormalTok{(}\VerbatimStringTok{\textasciigrave{}Error reading file: }\SpecialCharTok{$\{}\NormalTok{err}\OperatorTok{.}\AttributeTok{message}\SpecialCharTok{\}}\VerbatimStringTok{\textasciigrave{}}\NormalTok{)}\OperatorTok{;}
\NormalTok{  \}}
\NormalTok{\}}

\KeywordTok{async} \KeywordTok{function} \FunctionTok{writeFile}\NormalTok{() \{}
  \ControlFlowTok{try}\NormalTok{ \{}
    \ControlFlowTok{await}\NormalTok{ fs}\OperatorTok{.}\FunctionTok{writeFile}\NormalTok{(}\StringTok{"output.txt"}\OperatorTok{,} \StringTok{"Async write!}\SpecialCharTok{\textbackslash{}n}\StringTok{"}\NormalTok{)}\OperatorTok{;}
    \ControlFlowTok{await}\NormalTok{ fs}\OperatorTok{.}\FunctionTok{appendFile}\NormalTok{(}\StringTok{"output.txt"}\OperatorTok{,} \StringTok{"Appended!}\SpecialCharTok{\textbackslash{}n}\StringTok{"}\NormalTok{)}\OperatorTok{;}
    \BuiltInTok{console}\OperatorTok{.}\FunctionTok{log}\NormalTok{(}\StringTok{"Files written."}\NormalTok{)}\OperatorTok{;}
\NormalTok{  \} }\ControlFlowTok{catch}\NormalTok{ (err) \{}
    \BuiltInTok{console}\OperatorTok{.}\FunctionTok{log}\NormalTok{(}\VerbatimStringTok{\textasciigrave{}Error: }\SpecialCharTok{$\{}\NormalTok{err}\OperatorTok{.}\AttributeTok{message}\SpecialCharTok{\}}\VerbatimStringTok{\textasciigrave{}}\NormalTok{)}\OperatorTok{;}
\NormalTok{  \}}
\NormalTok{\}}

\FunctionTok{readFile}\NormalTok{()}\OperatorTok{;}
\FunctionTok{writeFile}\NormalTok{()}\OperatorTok{;}
\end{Highlighting}
\end{Shaded}

\begin{center}\rule{0.5\linewidth}{0.5pt}\end{center}

\subsection{Lesson 8.2: Error Handling}\label{lesson-8.2-error-handling}

\begin{Shaded}
\begin{Highlighting}[]
\CommentTok{// try{-}catch{-}finally}
\ControlFlowTok{try}\NormalTok{ \{}
  \KeywordTok{const}\NormalTok{ result }\OperatorTok{=} \DecValTok{10} \OperatorTok{/} \DecValTok{0}\OperatorTok{;} \CommentTok{// Doesn\textquotesingle{}t throw in JS (gives Infinity)}
  \KeywordTok{const}\NormalTok{ parsed }\OperatorTok{=} \BuiltInTok{JSON}\OperatorTok{.}\FunctionTok{parse}\NormalTok{(}\StringTok{"invalid json"}\NormalTok{)}\OperatorTok{;} \CommentTok{// This throws}
\NormalTok{\} }\ControlFlowTok{catch}\NormalTok{ (err) \{}
  \BuiltInTok{console}\OperatorTok{.}\FunctionTok{log}\NormalTok{(}\VerbatimStringTok{\textasciigrave{}Error type: }\SpecialCharTok{$\{}\NormalTok{err}\OperatorTok{.}\AttributeTok{constructor}\OperatorTok{.}\AttributeTok{name}\SpecialCharTok{\}}\VerbatimStringTok{\textasciigrave{}}\NormalTok{)}\OperatorTok{;}
  \BuiltInTok{console}\OperatorTok{.}\FunctionTok{log}\NormalTok{(}\VerbatimStringTok{\textasciigrave{}Error message: }\SpecialCharTok{$\{}\NormalTok{err}\OperatorTok{.}\AttributeTok{message}\SpecialCharTok{\}}\VerbatimStringTok{\textasciigrave{}}\NormalTok{)}\OperatorTok{;}
\NormalTok{\} }\ControlFlowTok{finally}\NormalTok{ \{}
  \BuiltInTok{console}\OperatorTok{.}\FunctionTok{log}\NormalTok{(}\StringTok{"This always runs."}\NormalTok{)}\OperatorTok{;}
\NormalTok{\}}

\CommentTok{// Throwing custom errors}
\KeywordTok{function} \FunctionTok{divide}\NormalTok{(a}\OperatorTok{,}\NormalTok{ b) \{}
  \ControlFlowTok{if}\NormalTok{ (b }\OperatorTok{===} \DecValTok{0}\NormalTok{) \{}
    \ControlFlowTok{throw} \KeywordTok{new} \BuiltInTok{Error}\NormalTok{(}\StringTok{"Division by zero is not allowed."}\NormalTok{)}\OperatorTok{;}
\NormalTok{  \}}
  \ControlFlowTok{return}\NormalTok{ a }\OperatorTok{/}\NormalTok{ b}\OperatorTok{;}
\NormalTok{\}}

\ControlFlowTok{try}\NormalTok{ \{}
  \BuiltInTok{console}\OperatorTok{.}\FunctionTok{log}\NormalTok{(}\FunctionTok{divide}\NormalTok{(}\DecValTok{10}\OperatorTok{,} \DecValTok{2}\NormalTok{))}\OperatorTok{;} \CommentTok{// 5}
  \BuiltInTok{console}\OperatorTok{.}\FunctionTok{log}\NormalTok{(}\FunctionTok{divide}\NormalTok{(}\DecValTok{10}\OperatorTok{,} \DecValTok{0}\NormalTok{))}\OperatorTok{;} \CommentTok{// Throws}
\NormalTok{\} }\ControlFlowTok{catch}\NormalTok{ (err) \{}
  \BuiltInTok{console}\OperatorTok{.}\FunctionTok{log}\NormalTok{(}\VerbatimStringTok{\textasciigrave{}Caught: }\SpecialCharTok{$\{}\NormalTok{err}\OperatorTok{.}\AttributeTok{message}\SpecialCharTok{\}}\VerbatimStringTok{\textasciigrave{}}\NormalTok{)}\OperatorTok{;}
\NormalTok{\}}
\end{Highlighting}
\end{Shaded}

\subsubsection{Custom Error Classes}\label{custom-error-classes}

\begin{Shaded}
\begin{Highlighting}[]
\KeywordTok{class}\NormalTok{ ValidationError }\KeywordTok{extends} \BuiltInTok{Error}\NormalTok{ \{}
  \FunctionTok{constructor}\NormalTok{(field}\OperatorTok{,}\NormalTok{ message) \{}
    \KeywordTok{super}\NormalTok{(message)}\OperatorTok{;}
    \KeywordTok{this}\OperatorTok{.}\AttributeTok{name} \OperatorTok{=} \StringTok{"ValidationError"}\OperatorTok{;}
    \KeywordTok{this}\OperatorTok{.}\AttributeTok{field} \OperatorTok{=}\NormalTok{ field}\OperatorTok{;}
\NormalTok{  \}}
\NormalTok{\}}

\KeywordTok{class}\NormalTok{ InsufficientFundsError }\KeywordTok{extends} \BuiltInTok{Error}\NormalTok{ \{}
  \FunctionTok{constructor}\NormalTok{(needed) \{}
    \KeywordTok{super}\NormalTok{(}\VerbatimStringTok{\textasciigrave{}Insufficient funds. Need $}\SpecialCharTok{$\{}\NormalTok{needed}\SpecialCharTok{\}}\VerbatimStringTok{ more.\textasciigrave{}}\NormalTok{)}\OperatorTok{;}
    \KeywordTok{this}\OperatorTok{.}\AttributeTok{name} \OperatorTok{=} \StringTok{"InsufficientFundsError"}\OperatorTok{;}
    \KeywordTok{this}\OperatorTok{.}\AttributeTok{needed} \OperatorTok{=}\NormalTok{ needed}\OperatorTok{;}
\NormalTok{  \}}
\NormalTok{\}}

\ControlFlowTok{try}\NormalTok{ \{}
  \ControlFlowTok{throw} \KeywordTok{new} \FunctionTok{ValidationError}\NormalTok{(}\StringTok{"email"}\OperatorTok{,} \StringTok{"Invalid email format."}\NormalTok{)}\OperatorTok{;}
\NormalTok{\} }\ControlFlowTok{catch}\NormalTok{ (err) \{}
  \ControlFlowTok{if}\NormalTok{ (err }\KeywordTok{instanceof}\NormalTok{ ValidationError) \{}
    \BuiltInTok{console}\OperatorTok{.}\FunctionTok{log}\NormalTok{(}\VerbatimStringTok{\textasciigrave{}Validation failed on \textquotesingle{}}\SpecialCharTok{$\{}\NormalTok{err}\OperatorTok{.}\AttributeTok{field}\SpecialCharTok{\}}\VerbatimStringTok{\textquotesingle{}: }\SpecialCharTok{$\{}\NormalTok{err}\OperatorTok{.}\AttributeTok{message}\SpecialCharTok{\}}\VerbatimStringTok{\textasciigrave{}}\NormalTok{)}\OperatorTok{;}
\NormalTok{  \} }\ControlFlowTok{else}\NormalTok{ \{}
    \ControlFlowTok{throw}\NormalTok{ err}\OperatorTok{;} \CommentTok{// Re{-}throw if not our custom error}
\NormalTok{  \}}
\NormalTok{\}}
\end{Highlighting}
\end{Shaded}

\begin{center}\rule{0.5\linewidth}{0.5pt}\end{center}

\subsection{Lesson 8.3: JSON in
Node.js}\label{lesson-8.3-json-in-node.js}

\begin{Shaded}
\begin{Highlighting}[]
\KeywordTok{const}\NormalTok{ fs }\OperatorTok{=} \PreprocessorTok{require}\NormalTok{(}\StringTok{"fs"}\NormalTok{)}\OperatorTok{;}

\CommentTok{// JavaScript object to JSON string}
\KeywordTok{const}\NormalTok{ settings }\OperatorTok{=}\NormalTok{ \{}
  \DataTypeTok{username}\OperatorTok{:} \StringTok{"coder\_kid"}\OperatorTok{,}
  \DataTypeTok{theme}\OperatorTok{:} \StringTok{"dark"}\OperatorTok{,}
  \DataTypeTok{notificationsOn}\OperatorTok{:} \KeywordTok{true}\OperatorTok{,}
  \DataTypeTok{favoriteNumbers}\OperatorTok{:}\NormalTok{ [}\DecValTok{7}\OperatorTok{,} \DecValTok{13}\OperatorTok{,} \DecValTok{42}\NormalTok{]}\OperatorTok{,}
\NormalTok{\}}\OperatorTok{;}

\KeywordTok{const}\NormalTok{ jsonString }\OperatorTok{=} \BuiltInTok{JSON}\OperatorTok{.}\FunctionTok{stringify}\NormalTok{(settings}\OperatorTok{,} \KeywordTok{null}\OperatorTok{,} \DecValTok{2}\NormalTok{)}\OperatorTok{;} \CommentTok{// Pretty{-}printed}
\BuiltInTok{console}\OperatorTok{.}\FunctionTok{log}\NormalTok{(jsonString)}\OperatorTok{;}

\CommentTok{// Save to file}
\NormalTok{fs}\OperatorTok{.}\FunctionTok{writeFileSync}\NormalTok{(}\StringTok{"settings.json"}\OperatorTok{,}\NormalTok{ jsonString)}\OperatorTok{;}

\CommentTok{// Load from file}
\ControlFlowTok{try}\NormalTok{ \{}
  \KeywordTok{const}\NormalTok{ loaded }\OperatorTok{=} \BuiltInTok{JSON}\OperatorTok{.}\FunctionTok{parse}\NormalTok{(fs}\OperatorTok{.}\FunctionTok{readFileSync}\NormalTok{(}\StringTok{"settings.json"}\OperatorTok{,} \StringTok{"utf8"}\NormalTok{))}\OperatorTok{;}
  \BuiltInTok{console}\OperatorTok{.}\FunctionTok{log}\NormalTok{(}\VerbatimStringTok{\textasciigrave{}Username: }\SpecialCharTok{$\{}\NormalTok{loaded}\OperatorTok{.}\AttributeTok{username}\SpecialCharTok{\}}\VerbatimStringTok{\textasciigrave{}}\NormalTok{)}\OperatorTok{;}
  \BuiltInTok{console}\OperatorTok{.}\FunctionTok{log}\NormalTok{(}\VerbatimStringTok{\textasciigrave{}Theme: }\SpecialCharTok{$\{}\NormalTok{loaded}\OperatorTok{.}\AttributeTok{theme}\SpecialCharTok{\}}\VerbatimStringTok{\textasciigrave{}}\NormalTok{)}\OperatorTok{;}
  \BuiltInTok{console}\OperatorTok{.}\FunctionTok{log}\NormalTok{(}\VerbatimStringTok{\textasciigrave{}Numbers: }\SpecialCharTok{$\{}\NormalTok{loaded}\OperatorTok{.}\AttributeTok{favoriteNumbers}\SpecialCharTok{\}}\VerbatimStringTok{\textasciigrave{}}\NormalTok{)}\OperatorTok{;}
\NormalTok{\} }\ControlFlowTok{catch}\NormalTok{ (err) \{}
  \ControlFlowTok{if}\NormalTok{ (err }\KeywordTok{instanceof} \BuiltInTok{SyntaxError}\NormalTok{) \{}
    \BuiltInTok{console}\OperatorTok{.}\FunctionTok{log}\NormalTok{(}\StringTok{"Error: JSON file is corrupted."}\NormalTok{)}\OperatorTok{;}
\NormalTok{  \} }\ControlFlowTok{else} \ControlFlowTok{if}\NormalTok{ (err}\OperatorTok{.}\AttributeTok{code} \OperatorTok{===} \StringTok{"ENOENT"}\NormalTok{) \{}
    \BuiltInTok{console}\OperatorTok{.}\FunctionTok{log}\NormalTok{(}\StringTok{"Error: File not found."}\NormalTok{)}\OperatorTok{;}
\NormalTok{  \} }\ControlFlowTok{else}\NormalTok{ \{}
    \BuiltInTok{console}\OperatorTok{.}\FunctionTok{log}\NormalTok{(}\VerbatimStringTok{\textasciigrave{}Error: }\SpecialCharTok{$\{}\NormalTok{err}\OperatorTok{.}\AttributeTok{message}\SpecialCharTok{\}}\VerbatimStringTok{\textasciigrave{}}\NormalTok{)}\OperatorTok{;}
\NormalTok{  \}}
\NormalTok{\}}
\end{Highlighting}
\end{Shaded}

\subsubsection{Skills Checklist}\label{skills-checklist-13}

\begin{itemize}
\tightlist
\item[$\square$]
  I can read and write files using the \texttt{fs} module.
\item[$\square$]
  I can use \texttt{try-catch-finally} to handle errors.
\item[$\square$]
  I can throw custom errors using \texttt{Error} or custom classes.
\item[$\square$]
  I can serialize/deserialize data using \texttt{JSON.stringify} and
  \texttt{JSON.parse}.
\end{itemize}

\newpage

\section{Unit 9: Modules, Collections, and
Testing}\label{unit-9-modules-collections-and-testing}

\begin{center}\rule{0.5\linewidth}{0.5pt}\end{center}

\subsection{Lesson 9.1: JavaScript
Modules}\label{lesson-9.1-javascript-modules}

\subsubsection{CommonJS (Node.js
default)}\label{commonjs-node.js-default}

\begin{Shaded}
\begin{Highlighting}[]
\CommentTok{// math.js {-}{-}{-} exporting}
\KeywordTok{const}\NormalTok{ add }\OperatorTok{=}\NormalTok{ (a}\OperatorTok{,}\NormalTok{ b) }\KeywordTok{=\textgreater{}}\NormalTok{ a }\OperatorTok{+}\NormalTok{ b}\OperatorTok{;}
\KeywordTok{const}\NormalTok{ multiply }\OperatorTok{=}\NormalTok{ (a}\OperatorTok{,}\NormalTok{ b) }\KeywordTok{=\textgreater{}}\NormalTok{ a }\OperatorTok{*}\NormalTok{ b}\OperatorTok{;}
\KeywordTok{const} \ConstantTok{PI} \OperatorTok{=} \BuiltInTok{Math}\OperatorTok{.}\ConstantTok{PI}\OperatorTok{;}

\NormalTok{module}\OperatorTok{.}\AttributeTok{exports} \OperatorTok{=}\NormalTok{ \{ add}\OperatorTok{,}\NormalTok{ multiply}\OperatorTok{,} \ConstantTok{PI}\NormalTok{ \}}\OperatorTok{;}

\CommentTok{// Or export individually:}
\CommentTok{// module.exports.add = add;}
\end{Highlighting}
\end{Shaded}

\begin{Shaded}
\begin{Highlighting}[]
\CommentTok{// main.js {-}{-}{-} importing}
\KeywordTok{const}\NormalTok{ \{ add}\OperatorTok{,}\NormalTok{ multiply}\OperatorTok{,} \ConstantTok{PI}\NormalTok{ \} }\OperatorTok{=} \PreprocessorTok{require}\NormalTok{(}\StringTok{"./math"}\NormalTok{)}\OperatorTok{;}

\BuiltInTok{console}\OperatorTok{.}\FunctionTok{log}\NormalTok{(}\FunctionTok{add}\NormalTok{(}\DecValTok{5}\OperatorTok{,} \DecValTok{3}\NormalTok{))}\OperatorTok{;} \CommentTok{// 8}
\BuiltInTok{console}\OperatorTok{.}\FunctionTok{log}\NormalTok{(}\FunctionTok{multiply}\NormalTok{(}\DecValTok{4}\OperatorTok{,} \DecValTok{6}\NormalTok{))}\OperatorTok{;} \CommentTok{// 24}
\BuiltInTok{console}\OperatorTok{.}\FunctionTok{log}\NormalTok{(}\ConstantTok{PI}\OperatorTok{.}\FunctionTok{toFixed}\NormalTok{(}\DecValTok{5}\NormalTok{))}\OperatorTok{;} \CommentTok{// 3.14159}
\end{Highlighting}
\end{Shaded}

\subsubsection{\texorpdfstring{ES Modules (Modern --- requires
\texttt{"type":\ "module"} in
package.json)}{ES Modules (Modern --- requires "type": "module" in package.json)}}\label{es-modules-modern-requires-type-module-in-package.json}

\begin{Shaded}
\begin{Highlighting}[]
\CommentTok{// math.mjs}
\ImportTok{export} \KeywordTok{const}\NormalTok{ add }\OperatorTok{=}\NormalTok{ (a}\OperatorTok{,}\NormalTok{ b) }\KeywordTok{=\textgreater{}}\NormalTok{ a }\OperatorTok{+}\NormalTok{ b}\OperatorTok{;}
\ImportTok{export} \KeywordTok{const}\NormalTok{ multiply }\OperatorTok{=}\NormalTok{ (a}\OperatorTok{,}\NormalTok{ b) }\KeywordTok{=\textgreater{}}\NormalTok{ a }\OperatorTok{*}\NormalTok{ b}\OperatorTok{;}
\ImportTok{export} \KeywordTok{const} \ConstantTok{PI} \OperatorTok{=} \BuiltInTok{Math}\OperatorTok{.}\ConstantTok{PI}\OperatorTok{;}
\ImportTok{export} \ImportTok{default} \KeywordTok{function} \FunctionTok{greet}\NormalTok{(name) \{}
  \ControlFlowTok{return} \VerbatimStringTok{\textasciigrave{}Hello, }\SpecialCharTok{$\{}\NormalTok{name}\SpecialCharTok{\}}\VerbatimStringTok{!\textasciigrave{}}\OperatorTok{;}
\NormalTok{\}}
\end{Highlighting}
\end{Shaded}

\begin{Shaded}
\begin{Highlighting}[]
\CommentTok{// main.mjs}
\ImportTok{import}\NormalTok{ greet}\OperatorTok{,}\NormalTok{ \{ add}\OperatorTok{,}\NormalTok{ multiply}\OperatorTok{,} \ConstantTok{PI}\NormalTok{ \} }\ImportTok{from} \StringTok{"./math.mjs"}\OperatorTok{;}

\BuiltInTok{console}\OperatorTok{.}\FunctionTok{log}\NormalTok{(}\FunctionTok{greet}\NormalTok{(}\StringTok{"Alice"}\NormalTok{))}\OperatorTok{;}
\BuiltInTok{console}\OperatorTok{.}\FunctionTok{log}\NormalTok{(}\FunctionTok{add}\NormalTok{(}\DecValTok{5}\OperatorTok{,} \DecValTok{3}\NormalTok{))}\OperatorTok{;}
\end{Highlighting}
\end{Shaded}

\begin{center}\rule{0.5\linewidth}{0.5pt}\end{center}

\subsection{Lesson 9.2: Maps and Sets}\label{lesson-9.2-maps-and-sets}

\subsubsection{\texorpdfstring{\texttt{Map} --- Key-Value (Better than
plain objects for
maps)}{Map --- Key-Value (Better than plain objects for maps)}}\label{map-key-value-better-than-plain-objects-for-maps}

\begin{Shaded}
\begin{Highlighting}[]
\KeywordTok{const}\NormalTok{ inventory }\OperatorTok{=} \KeywordTok{new} \BuiltInTok{Map}\NormalTok{()}\OperatorTok{;}

\CommentTok{// Adding}
\NormalTok{inventory}\OperatorTok{.}\FunctionTok{set}\NormalTok{(}\StringTok{"apples"}\OperatorTok{,} \DecValTok{50}\NormalTok{)}\OperatorTok{;}
\NormalTok{inventory}\OperatorTok{.}\FunctionTok{set}\NormalTok{(}\StringTok{"bananas"}\OperatorTok{,} \DecValTok{30}\NormalTok{)}\OperatorTok{;}
\NormalTok{inventory}\OperatorTok{.}\FunctionTok{set}\NormalTok{(}\StringTok{"oranges"}\OperatorTok{,} \DecValTok{40}\NormalTok{)}\OperatorTok{;}

\CommentTok{// Access}
\BuiltInTok{console}\OperatorTok{.}\FunctionTok{log}\NormalTok{(inventory}\OperatorTok{.}\FunctionTok{get}\NormalTok{(}\StringTok{"apples"}\NormalTok{))}\OperatorTok{;} \CommentTok{// 50}
\BuiltInTok{console}\OperatorTok{.}\FunctionTok{log}\NormalTok{(inventory}\OperatorTok{.}\FunctionTok{has}\NormalTok{(}\StringTok{"bananas"}\NormalTok{))}\OperatorTok{;} \CommentTok{// true}
\BuiltInTok{console}\OperatorTok{.}\FunctionTok{log}\NormalTok{(inventory}\OperatorTok{.}\AttributeTok{size}\NormalTok{)}\OperatorTok{;} \CommentTok{// 3}

\CommentTok{// Delete}
\NormalTok{inventory}\OperatorTok{.}\FunctionTok{delete}\NormalTok{(}\StringTok{"bananas"}\NormalTok{)}\OperatorTok{;}

\CommentTok{// Iterate}
\ControlFlowTok{for}\NormalTok{ (}\KeywordTok{const}\NormalTok{ [key}\OperatorTok{,}\NormalTok{ value] }\KeywordTok{of}\NormalTok{ inventory) \{}
  \BuiltInTok{console}\OperatorTok{.}\FunctionTok{log}\NormalTok{(}\VerbatimStringTok{\textasciigrave{}}\SpecialCharTok{$\{}\NormalTok{key}\SpecialCharTok{\}}\VerbatimStringTok{: }\SpecialCharTok{$\{}\NormalTok{value}\SpecialCharTok{\}}\VerbatimStringTok{\textasciigrave{}}\NormalTok{)}\OperatorTok{;}
\NormalTok{\}}

\NormalTok{inventory}\OperatorTok{.}\FunctionTok{forEach}\NormalTok{((value}\OperatorTok{,}\NormalTok{ key) }\KeywordTok{=\textgreater{}}\NormalTok{ \{}
  \BuiltInTok{console}\OperatorTok{.}\FunctionTok{log}\NormalTok{(}\VerbatimStringTok{\textasciigrave{}}\SpecialCharTok{$\{}\NormalTok{key}\SpecialCharTok{\}}\VerbatimStringTok{ = }\SpecialCharTok{$\{}\NormalTok{value}\SpecialCharTok{\}}\VerbatimStringTok{\textasciigrave{}}\NormalTok{)}\OperatorTok{;}
\NormalTok{\})}\OperatorTok{;}

\BuiltInTok{console}\OperatorTok{.}\FunctionTok{log}\NormalTok{([}\OperatorTok{...}\NormalTok{inventory}\OperatorTok{.}\FunctionTok{keys}\NormalTok{()])}\OperatorTok{;}
\BuiltInTok{console}\OperatorTok{.}\FunctionTok{log}\NormalTok{([}\OperatorTok{...}\NormalTok{inventory}\OperatorTok{.}\FunctionTok{values}\NormalTok{()])}\OperatorTok{;}
\BuiltInTok{console}\OperatorTok{.}\FunctionTok{log}\NormalTok{([}\OperatorTok{...}\NormalTok{inventory}\OperatorTok{.}\FunctionTok{entries}\NormalTok{()])}\OperatorTok{;}
\end{Highlighting}
\end{Shaded}

\subsubsection{\texorpdfstring{\texttt{Set} --- Unique
Values}{Set --- Unique Values}}\label{set-unique-values}

\begin{Shaded}
\begin{Highlighting}[]
\KeywordTok{const}\NormalTok{ uniqueWords }\OperatorTok{=} \KeywordTok{new} \BuiltInTok{Set}\NormalTok{()}\OperatorTok{;}
\NormalTok{uniqueWords}\OperatorTok{.}\FunctionTok{add}\NormalTok{(}\StringTok{"apple"}\NormalTok{)}\OperatorTok{;}
\NormalTok{uniqueWords}\OperatorTok{.}\FunctionTok{add}\NormalTok{(}\StringTok{"banana"}\NormalTok{)}\OperatorTok{;}
\NormalTok{uniqueWords}\OperatorTok{.}\FunctionTok{add}\NormalTok{(}\StringTok{"apple"}\NormalTok{)}\OperatorTok{;} \CommentTok{// Duplicate {-}{-}{-} ignored}

\BuiltInTok{console}\OperatorTok{.}\FunctionTok{log}\NormalTok{(uniqueWords}\OperatorTok{.}\AttributeTok{size}\NormalTok{)}\OperatorTok{;} \CommentTok{// 2}
\BuiltInTok{console}\OperatorTok{.}\FunctionTok{log}\NormalTok{(uniqueWords}\OperatorTok{.}\FunctionTok{has}\NormalTok{(}\StringTok{"banana"}\NormalTok{))}\OperatorTok{;} \CommentTok{// true}
\NormalTok{uniqueWords}\OperatorTok{.}\FunctionTok{delete}\NormalTok{(}\StringTok{"banana"}\NormalTok{)}\OperatorTok{;}

\CommentTok{// Convert array to set (remove duplicates)}
\KeywordTok{const}\NormalTok{ nums }\OperatorTok{=}\NormalTok{ [}\DecValTok{1}\OperatorTok{,} \DecValTok{2}\OperatorTok{,} \DecValTok{2}\OperatorTok{,} \DecValTok{3}\OperatorTok{,} \DecValTok{3}\OperatorTok{,} \DecValTok{3}\OperatorTok{,} \DecValTok{4}\NormalTok{]}\OperatorTok{;}
\KeywordTok{const}\NormalTok{ unique }\OperatorTok{=}\NormalTok{ [}\OperatorTok{...}\KeywordTok{new} \BuiltInTok{Set}\NormalTok{(nums)]}\OperatorTok{;}
\BuiltInTok{console}\OperatorTok{.}\FunctionTok{log}\NormalTok{(unique)}\OperatorTok{;} \CommentTok{// [1, 2, 3, 4]}

\CommentTok{// Iterate}
\ControlFlowTok{for}\NormalTok{ (}\KeywordTok{const}\NormalTok{ item }\KeywordTok{of}\NormalTok{ uniqueWords) \{}
  \BuiltInTok{console}\OperatorTok{.}\FunctionTok{log}\NormalTok{(item)}\OperatorTok{;}
\NormalTok{\}}
\end{Highlighting}
\end{Shaded}

\begin{center}\rule{0.5\linewidth}{0.5pt}\end{center}

\subsection{Lesson 9.3: Word Frequency
Analysis}\label{lesson-9.3-word-frequency-analysis}

\begin{Shaded}
\begin{Highlighting}[]
\KeywordTok{const}\NormalTok{ fs }\OperatorTok{=} \PreprocessorTok{require}\NormalTok{(}\StringTok{"fs"}\NormalTok{)}\OperatorTok{;}

\KeywordTok{function} \FunctionTok{countWords}\NormalTok{(text) \{}
  \KeywordTok{const}\NormalTok{ words }\OperatorTok{=}\NormalTok{ text}
    \OperatorTok{.}\FunctionTok{toLowerCase}\NormalTok{()}
    \OperatorTok{.}\FunctionTok{replace}\NormalTok{(}\SpecialStringTok{/}\SpecialCharTok{[\^{}a{-}z\textbackslash{}s]}\SpecialStringTok{/g}\OperatorTok{,} \StringTok{""}\NormalTok{)}
    \OperatorTok{.}\FunctionTok{trim}\NormalTok{()}
    \OperatorTok{.}\FunctionTok{split}\NormalTok{(}\SpecialStringTok{/}\SpecialCharTok{\textbackslash{}s+}\SpecialStringTok{/}\NormalTok{)}
    \OperatorTok{.}\FunctionTok{filter}\NormalTok{((word) }\KeywordTok{=\textgreater{}}\NormalTok{ word}\OperatorTok{.}\AttributeTok{length} \OperatorTok{\textgreater{}} \DecValTok{0}\NormalTok{)}\OperatorTok{;}

  \KeywordTok{const}\NormalTok{ frequency }\OperatorTok{=} \KeywordTok{new} \BuiltInTok{Map}\NormalTok{()}\OperatorTok{;}
  \ControlFlowTok{for}\NormalTok{ (}\KeywordTok{const}\NormalTok{ word }\KeywordTok{of}\NormalTok{ words) \{}
\NormalTok{    frequency}\OperatorTok{.}\FunctionTok{set}\NormalTok{(word}\OperatorTok{,}\NormalTok{ (frequency}\OperatorTok{.}\FunctionTok{get}\NormalTok{(word) }\OperatorTok{||} \DecValTok{0}\NormalTok{) }\OperatorTok{+} \DecValTok{1}\NormalTok{)}\OperatorTok{;}
\NormalTok{  \}}

  \ControlFlowTok{return}\NormalTok{ frequency}\OperatorTok{;}
\NormalTok{\}}

\KeywordTok{const}\NormalTok{ text }\OperatorTok{=} \StringTok{"the quick brown fox jumps over the lazy dog the quick brown fox"}\OperatorTok{;}
\KeywordTok{const}\NormalTok{ freq }\OperatorTok{=} \FunctionTok{countWords}\NormalTok{(text)}\OperatorTok{;}

\CommentTok{// Sort by frequency (most common first)}
\KeywordTok{const}\NormalTok{ sorted }\OperatorTok{=}\NormalTok{ [}\OperatorTok{...}\NormalTok{freq}\OperatorTok{.}\FunctionTok{entries}\NormalTok{()]}\OperatorTok{.}\FunctionTok{sort}\NormalTok{((a}\OperatorTok{,}\NormalTok{ b) }\KeywordTok{=\textgreater{}}\NormalTok{ b[}\DecValTok{1}\NormalTok{] }\OperatorTok{{-}}\NormalTok{ a[}\DecValTok{1}\NormalTok{])}\OperatorTok{;}

\BuiltInTok{console}\OperatorTok{.}\FunctionTok{log}\NormalTok{(}\StringTok{"Word Frequencies:"}\NormalTok{)}\OperatorTok{;}
\ControlFlowTok{for}\NormalTok{ (}\KeywordTok{const}\NormalTok{ [word}\OperatorTok{,}\NormalTok{ count] }\KeywordTok{of}\NormalTok{ sorted) \{}
  \BuiltInTok{console}\OperatorTok{.}\FunctionTok{log}\NormalTok{(}\VerbatimStringTok{\textasciigrave{}\textquotesingle{}}\SpecialCharTok{$\{}\NormalTok{word}\SpecialCharTok{\}}\VerbatimStringTok{\textquotesingle{}: }\SpecialCharTok{$\{}\NormalTok{count}\SpecialCharTok{\}}\VerbatimStringTok{\textasciigrave{}}\NormalTok{)}\OperatorTok{;}
\NormalTok{\}}

\CommentTok{// Top 5 words}
\BuiltInTok{console}\OperatorTok{.}\FunctionTok{log}\NormalTok{(}\StringTok{"}\SpecialCharTok{\textbackslash{}n}\StringTok{Top 5 words:"}\NormalTok{)}\OperatorTok{;}
\NormalTok{sorted}\OperatorTok{.}\FunctionTok{slice}\NormalTok{(}\DecValTok{0}\OperatorTok{,} \DecValTok{5}\NormalTok{)}\OperatorTok{.}\FunctionTok{forEach}\NormalTok{(([word}\OperatorTok{,}\NormalTok{ count]) }\KeywordTok{=\textgreater{}}\NormalTok{ \{}
  \BuiltInTok{console}\OperatorTok{.}\FunctionTok{log}\NormalTok{(}\VerbatimStringTok{\textasciigrave{}  }\SpecialCharTok{$\{}\NormalTok{word}\SpecialCharTok{\}}\VerbatimStringTok{: }\SpecialCharTok{$\{}\NormalTok{count}\SpecialCharTok{\}}\VerbatimStringTok{\textasciigrave{}}\NormalTok{)}\OperatorTok{;}
\NormalTok{\})}\OperatorTok{;}
\end{Highlighting}
\end{Shaded}

\begin{center}\rule{0.5\linewidth}{0.5pt}\end{center}

\subsection{Lesson 9.4: Working with CSV
Files}\label{lesson-9.4-working-with-csv-files}

\begin{Shaded}
\begin{Highlighting}[]
\KeywordTok{const}\NormalTok{ fs }\OperatorTok{=} \PreprocessorTok{require}\NormalTok{(}\StringTok{"fs"}\NormalTok{)}\OperatorTok{;}

\KeywordTok{function} \FunctionTok{parseCSV}\NormalTok{(content) \{}
  \KeywordTok{const}\NormalTok{ lines }\OperatorTok{=}\NormalTok{ content}\OperatorTok{.}\FunctionTok{trim}\NormalTok{()}\OperatorTok{.}\FunctionTok{split}\NormalTok{(}\StringTok{"}\SpecialCharTok{\textbackslash{}n}\StringTok{"}\NormalTok{)}\OperatorTok{;}
  \KeywordTok{const}\NormalTok{ headers }\OperatorTok{=}\NormalTok{ lines[}\DecValTok{0}\NormalTok{]}\OperatorTok{.}\FunctionTok{split}\NormalTok{(}\StringTok{","}\NormalTok{)}\OperatorTok{.}\FunctionTok{map}\NormalTok{((h) }\KeywordTok{=\textgreater{}}\NormalTok{ h}\OperatorTok{.}\FunctionTok{trim}\NormalTok{())}\OperatorTok{;}

  \ControlFlowTok{return}\NormalTok{ lines}\OperatorTok{.}\FunctionTok{slice}\NormalTok{(}\DecValTok{1}\NormalTok{)}\OperatorTok{.}\FunctionTok{map}\NormalTok{((line) }\KeywordTok{=\textgreater{}}\NormalTok{ \{}
    \KeywordTok{const}\NormalTok{ values }\OperatorTok{=}\NormalTok{ line}\OperatorTok{.}\FunctionTok{split}\NormalTok{(}\StringTok{","}\NormalTok{)}\OperatorTok{.}\FunctionTok{map}\NormalTok{((v) }\KeywordTok{=\textgreater{}}\NormalTok{ v}\OperatorTok{.}\FunctionTok{trim}\NormalTok{())}\OperatorTok{;}
    \KeywordTok{const}\NormalTok{ record }\OperatorTok{=}\NormalTok{ \{\}}\OperatorTok{;}
\NormalTok{    headers}\OperatorTok{.}\FunctionTok{forEach}\NormalTok{((header}\OperatorTok{,}\NormalTok{ index) }\KeywordTok{=\textgreater{}}\NormalTok{ \{}
\NormalTok{      record[header] }\OperatorTok{=}\NormalTok{ values[index]}\OperatorTok{;}
\NormalTok{    \})}\OperatorTok{;}
    \ControlFlowTok{return}\NormalTok{ record}\OperatorTok{;}
\NormalTok{  \})}\OperatorTok{;}
\NormalTok{\}}

\KeywordTok{function} \FunctionTok{generateCSV}\NormalTok{(records) \{}
  \ControlFlowTok{if}\NormalTok{ (records}\OperatorTok{.}\AttributeTok{length} \OperatorTok{===} \DecValTok{0}\NormalTok{) }\ControlFlowTok{return} \StringTok{""}\OperatorTok{;}
  \KeywordTok{const}\NormalTok{ headers }\OperatorTok{=} \BuiltInTok{Object}\OperatorTok{.}\FunctionTok{keys}\NormalTok{(records[}\DecValTok{0}\NormalTok{])}\OperatorTok{;}
  \KeywordTok{const}\NormalTok{ headerLine }\OperatorTok{=}\NormalTok{ headers}\OperatorTok{.}\FunctionTok{join}\NormalTok{(}\StringTok{","}\NormalTok{)}\OperatorTok{;}
  \KeywordTok{const}\NormalTok{ dataLines }\OperatorTok{=}\NormalTok{ records}\OperatorTok{.}\FunctionTok{map}\NormalTok{((record) }\KeywordTok{=\textgreater{}}
\NormalTok{    headers}\OperatorTok{.}\FunctionTok{map}\NormalTok{((h) }\KeywordTok{=\textgreater{}}\NormalTok{ record[h])}\OperatorTok{.}\FunctionTok{join}\NormalTok{(}\StringTok{","}\NormalTok{)}\OperatorTok{,}
\NormalTok{  )}\OperatorTok{;}
  \ControlFlowTok{return}\NormalTok{ [headerLine}\OperatorTok{,} \OperatorTok{...}\NormalTok{dataLines]}\OperatorTok{.}\FunctionTok{join}\NormalTok{(}\StringTok{"}\SpecialCharTok{\textbackslash{}n}\StringTok{"}\NormalTok{)}\OperatorTok{;}
\NormalTok{\}}

\CommentTok{// Read and parse CSV}
\ControlFlowTok{try}\NormalTok{ \{}
  \KeywordTok{const}\NormalTok{ content }\OperatorTok{=}\NormalTok{ fs}\OperatorTok{.}\FunctionTok{readFileSync}\NormalTok{(}\StringTok{"students.csv"}\OperatorTok{,} \StringTok{"utf8"}\NormalTok{)}\OperatorTok{;}
  \KeywordTok{const}\NormalTok{ students }\OperatorTok{=} \FunctionTok{parseCSV}\NormalTok{(content)}\OperatorTok{;}

  \BuiltInTok{console}\OperatorTok{.}\FunctionTok{log}\NormalTok{(}\VerbatimStringTok{\textasciigrave{}Loaded }\SpecialCharTok{$\{}\NormalTok{students}\OperatorTok{.}\AttributeTok{length}\SpecialCharTok{\}}\VerbatimStringTok{ students:\textasciigrave{}}\NormalTok{)}\OperatorTok{;}
\NormalTok{  students}\OperatorTok{.}\FunctionTok{forEach}\NormalTok{((student) }\KeywordTok{=\textgreater{}}\NormalTok{ \{}
    \BuiltInTok{console}\OperatorTok{.}\FunctionTok{log}\NormalTok{(}
      \VerbatimStringTok{\textasciigrave{}  Name: }\SpecialCharTok{$\{}\NormalTok{student}\OperatorTok{.}\AttributeTok{name}\SpecialCharTok{\}}\VerbatimStringTok{, Age: }\SpecialCharTok{$\{}\NormalTok{student}\OperatorTok{.}\AttributeTok{age}\SpecialCharTok{\}}\VerbatimStringTok{, Grade: }\SpecialCharTok{$\{}\NormalTok{student}\OperatorTok{.}\AttributeTok{grade}\SpecialCharTok{\}}\VerbatimStringTok{\textasciigrave{}}\OperatorTok{,}
\NormalTok{    )}\OperatorTok{;}
\NormalTok{  \})}\OperatorTok{;}
\NormalTok{\} }\ControlFlowTok{catch}\NormalTok{ (err) \{}
  \ControlFlowTok{if}\NormalTok{ (err}\OperatorTok{.}\AttributeTok{code} \OperatorTok{===} \StringTok{"ENOENT"}\NormalTok{) \{}
    \CommentTok{// Create a sample CSV for demonstration}
    \KeywordTok{const}\NormalTok{ sampleData }\OperatorTok{=} \VerbatimStringTok{\textasciigrave{}name,age,grade}
\VerbatimStringTok{Alice,15,A}
\VerbatimStringTok{Bob,14,B}
\VerbatimStringTok{Charlie,16,A\textasciigrave{}}\OperatorTok{;}
\NormalTok{    fs}\OperatorTok{.}\FunctionTok{writeFileSync}\NormalTok{(}\StringTok{"students.csv"}\OperatorTok{,}\NormalTok{ sampleData)}\OperatorTok{;}
    \BuiltInTok{console}\OperatorTok{.}\FunctionTok{log}\NormalTok{(}\StringTok{"Created sample students.csv"}\NormalTok{)}\OperatorTok{;}
\NormalTok{  \}}
\NormalTok{\}}
\end{Highlighting}
\end{Shaded}

\begin{center}\rule{0.5\linewidth}{0.5pt}\end{center}

\subsection{Lesson 9.5: Testing with
Jest}\label{lesson-9.5-testing-with-jest}

\subsubsection{Setting Up Jest}\label{setting-up-jest}

\begin{Shaded}
\begin{Highlighting}[]
\ExtensionTok{npm}\NormalTok{ init }\AttributeTok{{-}y}
\ExtensionTok{npm}\NormalTok{ install }\AttributeTok{{-}{-}save{-}dev}\NormalTok{ jest}
\end{Highlighting}
\end{Shaded}

Add to \texttt{package.json}:

\begin{Shaded}
\begin{Highlighting}[]
\FunctionTok{\{}
  \DataTypeTok{"scripts"}\FunctionTok{:} \FunctionTok{\{}
    \DataTypeTok{"test"}\FunctionTok{:} \StringTok{"jest"}
  \FunctionTok{\}}
\FunctionTok{\}}
\end{Highlighting}
\end{Shaded}

\subsubsection{Writing Tests}\label{writing-tests}

\begin{Shaded}
\begin{Highlighting}[]
\CommentTok{// math.js}
\KeywordTok{const}\NormalTok{ add }\OperatorTok{=}\NormalTok{ (a}\OperatorTok{,}\NormalTok{ b) }\KeywordTok{=\textgreater{}}\NormalTok{ a }\OperatorTok{+}\NormalTok{ b}\OperatorTok{;}
\KeywordTok{const}\NormalTok{ getGrade }\OperatorTok{=}\NormalTok{ (score) }\KeywordTok{=\textgreater{}}\NormalTok{ \{}
  \ControlFlowTok{if}\NormalTok{ (score }\OperatorTok{\textgreater{}=} \DecValTok{90}\NormalTok{) }\ControlFlowTok{return} \StringTok{"A"}\OperatorTok{;}
  \ControlFlowTok{if}\NormalTok{ (score }\OperatorTok{\textgreater{}=} \DecValTok{80}\NormalTok{) }\ControlFlowTok{return} \StringTok{"B"}\OperatorTok{;}
  \ControlFlowTok{if}\NormalTok{ (score }\OperatorTok{\textgreater{}=} \DecValTok{70}\NormalTok{) }\ControlFlowTok{return} \StringTok{"C"}\OperatorTok{;}
  \ControlFlowTok{return} \StringTok{"F"}\OperatorTok{;}
\NormalTok{\}}\OperatorTok{;}
\KeywordTok{const}\NormalTok{ circleArea }\OperatorTok{=}\NormalTok{ (radius) }\KeywordTok{=\textgreater{}} \BuiltInTok{Math}\OperatorTok{.}\ConstantTok{PI} \OperatorTok{*}\NormalTok{ radius }\OperatorTok{**} \DecValTok{2}\OperatorTok{;}

\NormalTok{module}\OperatorTok{.}\AttributeTok{exports} \OperatorTok{=}\NormalTok{ \{ add}\OperatorTok{,}\NormalTok{ getGrade}\OperatorTok{,}\NormalTok{ circleArea \}}\OperatorTok{;}
\end{Highlighting}
\end{Shaded}

\begin{Shaded}
\begin{Highlighting}[]
\CommentTok{// math.test.js}
\KeywordTok{const}\NormalTok{ \{ add}\OperatorTok{,}\NormalTok{ getGrade}\OperatorTok{,}\NormalTok{ circleArea \} }\OperatorTok{=} \PreprocessorTok{require}\NormalTok{(}\StringTok{"./math"}\NormalTok{)}\OperatorTok{;}

\FunctionTok{describe}\NormalTok{(}\StringTok{"add function"}\OperatorTok{,}\NormalTok{ () }\KeywordTok{=\textgreater{}}\NormalTok{ \{}
  \FunctionTok{test}\NormalTok{(}\StringTok{"adds two positive numbers"}\OperatorTok{,}\NormalTok{ () }\KeywordTok{=\textgreater{}}\NormalTok{ \{}
    \FunctionTok{expect}\NormalTok{(}\FunctionTok{add}\NormalTok{(}\DecValTok{5}\OperatorTok{,} \DecValTok{3}\NormalTok{))}\OperatorTok{.}\FunctionTok{toBe}\NormalTok{(}\DecValTok{8}\NormalTok{)}\OperatorTok{;}
\NormalTok{  \})}\OperatorTok{;}

  \FunctionTok{test}\NormalTok{(}\StringTok{"adds negative numbers"}\OperatorTok{,}\NormalTok{ () }\KeywordTok{=\textgreater{}}\NormalTok{ \{}
    \FunctionTok{expect}\NormalTok{(}\FunctionTok{add}\NormalTok{(}\OperatorTok{{-}}\DecValTok{5}\OperatorTok{,} \OperatorTok{{-}}\DecValTok{3}\NormalTok{))}\OperatorTok{.}\FunctionTok{toBe}\NormalTok{(}\OperatorTok{{-}}\DecValTok{8}\NormalTok{)}\OperatorTok{;}
\NormalTok{  \})}\OperatorTok{;}

  \FunctionTok{test}\NormalTok{(}\StringTok{"adds zero"}\OperatorTok{,}\NormalTok{ () }\KeywordTok{=\textgreater{}}\NormalTok{ \{}
    \FunctionTok{expect}\NormalTok{(}\FunctionTok{add}\NormalTok{(}\DecValTok{0}\OperatorTok{,} \DecValTok{7}\NormalTok{))}\OperatorTok{.}\FunctionTok{toBe}\NormalTok{(}\DecValTok{7}\NormalTok{)}\OperatorTok{;}
\NormalTok{  \})}\OperatorTok{;}
\NormalTok{\})}\OperatorTok{;}

\FunctionTok{describe}\NormalTok{(}\StringTok{"getGrade function"}\OperatorTok{,}\NormalTok{ () }\KeywordTok{=\textgreater{}}\NormalTok{ \{}
  \FunctionTok{test}\NormalTok{(}\StringTok{"returns A for score \textgreater{}= 90"}\OperatorTok{,}\NormalTok{ () }\KeywordTok{=\textgreater{}}\NormalTok{ \{}
    \FunctionTok{expect}\NormalTok{(}\FunctionTok{getGrade}\NormalTok{(}\DecValTok{95}\NormalTok{))}\OperatorTok{.}\FunctionTok{toBe}\NormalTok{(}\StringTok{"A"}\NormalTok{)}\OperatorTok{;}
    \FunctionTok{expect}\NormalTok{(}\FunctionTok{getGrade}\NormalTok{(}\DecValTok{90}\NormalTok{))}\OperatorTok{.}\FunctionTok{toBe}\NormalTok{(}\StringTok{"A"}\NormalTok{)}\OperatorTok{;}
\NormalTok{  \})}\OperatorTok{;}

  \FunctionTok{test}\NormalTok{(}\StringTok{"returns B for score 80{-}89"}\OperatorTok{,}\NormalTok{ () }\KeywordTok{=\textgreater{}}\NormalTok{ \{}
    \FunctionTok{expect}\NormalTok{(}\FunctionTok{getGrade}\NormalTok{(}\DecValTok{85}\NormalTok{))}\OperatorTok{.}\FunctionTok{toBe}\NormalTok{(}\StringTok{"B"}\NormalTok{)}\OperatorTok{;}
\NormalTok{  \})}\OperatorTok{;}

  \FunctionTok{test}\NormalTok{(}\StringTok{"returns F for score \textless{} 60"}\OperatorTok{,}\NormalTok{ () }\KeywordTok{=\textgreater{}}\NormalTok{ \{}
    \FunctionTok{expect}\NormalTok{(}\FunctionTok{getGrade}\NormalTok{(}\DecValTok{55}\NormalTok{))}\OperatorTok{.}\FunctionTok{toBe}\NormalTok{(}\StringTok{"F"}\NormalTok{)}\OperatorTok{;}
\NormalTok{  \})}\OperatorTok{;}
\NormalTok{\})}\OperatorTok{;}

\FunctionTok{describe}\NormalTok{(}\StringTok{"circleArea function"}\OperatorTok{,}\NormalTok{ () }\KeywordTok{=\textgreater{}}\NormalTok{ \{}
  \FunctionTok{test}\NormalTok{(}\StringTok{"calculates area of circle"}\OperatorTok{,}\NormalTok{ () }\KeywordTok{=\textgreater{}}\NormalTok{ \{}
    \FunctionTok{expect}\NormalTok{(}\FunctionTok{circleArea}\NormalTok{(}\DecValTok{5}\NormalTok{))}\OperatorTok{.}\FunctionTok{toBeCloseTo}\NormalTok{(}\FloatTok{78.54}\OperatorTok{,} \DecValTok{1}\NormalTok{)}\OperatorTok{;}
\NormalTok{  \})}\OperatorTok{;}
\NormalTok{\})}\OperatorTok{;}
\end{Highlighting}
\end{Shaded}

Run tests:

\begin{Shaded}
\begin{Highlighting}[]
\ExtensionTok{npm}\NormalTok{ test}
\end{Highlighting}
\end{Shaded}

\subsubsection{Common Jest Matchers}\label{common-jest-matchers}

{\def\LTcaptype{none} % do not increment counter
\begin{longtable}[]{@{}ll@{}}
\toprule\noalign{}
Matcher & Description \\
\midrule\noalign{}
\endhead
\bottomrule\noalign{}
\endlastfoot
\texttt{toBe(value)} & Strict equality (\texttt{===}) \\
\texttt{toEqual(value)} & Deep equality (for objects/arrays) \\
\texttt{toBeCloseTo(num,\ decimals)} & Floating-point comparison \\
\texttt{toBeTruthy()} & Check if truthy \\
\texttt{toBeFalsy()} & Check if falsy \\
\texttt{toBeNull()} & Check if null \\
\texttt{toBeUndefined()} & Check if undefined \\
\texttt{toContain(item)} & Array/string contains item \\
\texttt{toHaveLength(n)} & Check length \\
\texttt{toThrow(error)} & Check if function throws \\
\end{longtable}
}

\subsubsection{Skills Checklist}\label{skills-checklist-14}

\begin{itemize}
\tightlist
\item[$\square$]
  I can create and import modules using CommonJS.
\item[$\square$]
  I can use \texttt{Map} and \texttt{Set} for collections.
\item[$\square$]
  I can write and run Jest tests.
\item[$\square$]
  I can use common Jest matchers.
\end{itemize}

\begin{center}\rule{0.5\linewidth}{0.5pt}\end{center}

\section{Unit 10: Project-Based
Learning}\label{unit-10-project-based-learning}

\begin{center}\rule{0.5\linewidth}{0.5pt}\end{center}

\subsection{Lesson 10.1: Text-Based Adventure
Game}\label{lesson-10.1-text-based-adventure-game}

\begin{Shaded}
\begin{Highlighting}[]
\KeywordTok{class}\NormalTok{ Room \{}
  \FunctionTok{constructor}\NormalTok{(name}\OperatorTok{,}\NormalTok{ description) \{}
    \KeywordTok{this}\OperatorTok{.}\AttributeTok{name} \OperatorTok{=}\NormalTok{ name}\OperatorTok{;}
    \KeywordTok{this}\OperatorTok{.}\AttributeTok{description} \OperatorTok{=}\NormalTok{ description}\OperatorTok{;}
    \KeywordTok{this}\OperatorTok{.}\AttributeTok{exits} \OperatorTok{=}\NormalTok{ \{\}}\OperatorTok{;}
    \KeywordTok{this}\OperatorTok{.}\AttributeTok{items} \OperatorTok{=}\NormalTok{ []}\OperatorTok{;}
\NormalTok{  \}}

  \FunctionTok{addExit}\NormalTok{(direction}\OperatorTok{,}\NormalTok{ room) \{}
    \KeywordTok{this}\OperatorTok{.}\AttributeTok{exits}\NormalTok{[direction}\OperatorTok{.}\FunctionTok{toLowerCase}\NormalTok{()] }\OperatorTok{=}\NormalTok{ room}\OperatorTok{;}
\NormalTok{  \}}

  \FunctionTok{addItem}\NormalTok{(item) \{}
    \KeywordTok{this}\OperatorTok{.}\AttributeTok{items}\OperatorTok{.}\FunctionTok{push}\NormalTok{(item)}\OperatorTok{;}
\NormalTok{  \}}

  \FunctionTok{describe}\NormalTok{() \{}
    \KeywordTok{let}\NormalTok{ desc }\OperatorTok{=} \VerbatimStringTok{\textasciigrave{}{-}{-}{-} }\SpecialCharTok{$\{}\KeywordTok{this}\OperatorTok{.}\AttributeTok{name}\SpecialCharTok{\}}\VerbatimStringTok{ {-}{-}{-}}\SpecialCharTok{\textbackslash{}n$\{}\KeywordTok{this}\OperatorTok{.}\AttributeTok{description}\SpecialCharTok{\}\textbackslash{}n}\VerbatimStringTok{\textasciigrave{}}\OperatorTok{;}
    \ControlFlowTok{if}\NormalTok{ (}\KeywordTok{this}\OperatorTok{.}\AttributeTok{items}\OperatorTok{.}\AttributeTok{length} \OperatorTok{\textgreater{}} \DecValTok{0}\NormalTok{) \{}
\NormalTok{      desc }\OperatorTok{+=} \VerbatimStringTok{\textasciigrave{}You see: }\SpecialCharTok{$\{}\KeywordTok{this}\OperatorTok{.}\AttributeTok{items}\OperatorTok{.}\FunctionTok{join}\NormalTok{(}\StringTok{", "}\NormalTok{)}\SpecialCharTok{\}\textbackslash{}n}\VerbatimStringTok{\textasciigrave{}}\OperatorTok{;}
\NormalTok{    \}}
    \KeywordTok{const}\NormalTok{ exitList }\OperatorTok{=} \BuiltInTok{Object}\OperatorTok{.}\FunctionTok{keys}\NormalTok{(}\KeywordTok{this}\OperatorTok{.}\AttributeTok{exits}\NormalTok{)}\OperatorTok{.}\FunctionTok{join}\NormalTok{(}\StringTok{", "}\NormalTok{)}\OperatorTok{;}
\NormalTok{    desc }\OperatorTok{+=} \VerbatimStringTok{\textasciigrave{}Exits: }\SpecialCharTok{$\{}\NormalTok{exitList}\SpecialCharTok{\}}\VerbatimStringTok{\textasciigrave{}}\OperatorTok{;}
    \ControlFlowTok{return}\NormalTok{ desc}\OperatorTok{;}
\NormalTok{  \}}
\NormalTok{\}}

\KeywordTok{class}\NormalTok{ Player \{}
  \FunctionTok{constructor}\NormalTok{(name}\OperatorTok{,}\NormalTok{ startRoom) \{}
    \KeywordTok{this}\OperatorTok{.}\AttributeTok{name} \OperatorTok{=}\NormalTok{ name}\OperatorTok{;}
    \KeywordTok{this}\OperatorTok{.}\AttributeTok{currentRoom} \OperatorTok{=}\NormalTok{ startRoom}\OperatorTok{;}
    \KeywordTok{this}\OperatorTok{.}\AttributeTok{inventory} \OperatorTok{=}\NormalTok{ []}\OperatorTok{;}
    \KeywordTok{this}\OperatorTok{.}\AttributeTok{health} \OperatorTok{=} \DecValTok{100}\OperatorTok{;}
\NormalTok{  \}}

  \FunctionTok{move}\NormalTok{(direction) \{}
    \KeywordTok{const}\NormalTok{ nextRoom }\OperatorTok{=} \KeywordTok{this}\OperatorTok{.}\AttributeTok{currentRoom}\OperatorTok{.}\AttributeTok{exits}\NormalTok{[direction}\OperatorTok{.}\FunctionTok{toLowerCase}\NormalTok{()]}\OperatorTok{;}
    \ControlFlowTok{if}\NormalTok{ (nextRoom) \{}
      \KeywordTok{this}\OperatorTok{.}\AttributeTok{currentRoom} \OperatorTok{=}\NormalTok{ nextRoom}\OperatorTok{;}
      \ControlFlowTok{return} \KeywordTok{true}\OperatorTok{;}
\NormalTok{    \}}
    \ControlFlowTok{return} \KeywordTok{false}\OperatorTok{;}
\NormalTok{  \}}

  \FunctionTok{takeItem}\NormalTok{(itemName) \{}
    \KeywordTok{const}\NormalTok{ index }\OperatorTok{=} \KeywordTok{this}\OperatorTok{.}\AttributeTok{currentRoom}\OperatorTok{.}\AttributeTok{items}\OperatorTok{.}\FunctionTok{findIndex}\NormalTok{(}
\NormalTok{      (item) }\KeywordTok{=\textgreater{}}\NormalTok{ item}\OperatorTok{.}\FunctionTok{toLowerCase}\NormalTok{() }\OperatorTok{===}\NormalTok{ itemName}\OperatorTok{.}\FunctionTok{toLowerCase}\NormalTok{()}\OperatorTok{,}
\NormalTok{    )}\OperatorTok{;}
    \ControlFlowTok{if}\NormalTok{ (index }\OperatorTok{!==} \OperatorTok{{-}}\DecValTok{1}\NormalTok{) \{}
      \KeywordTok{const}\NormalTok{ item }\OperatorTok{=} \KeywordTok{this}\OperatorTok{.}\AttributeTok{currentRoom}\OperatorTok{.}\AttributeTok{items}\OperatorTok{.}\FunctionTok{splice}\NormalTok{(index}\OperatorTok{,} \DecValTok{1}\NormalTok{)[}\DecValTok{0}\NormalTok{]}\OperatorTok{;}
      \KeywordTok{this}\OperatorTok{.}\AttributeTok{inventory}\OperatorTok{.}\FunctionTok{push}\NormalTok{(item)}\OperatorTok{;}
      \ControlFlowTok{return} \KeywordTok{true}\OperatorTok{;}
\NormalTok{    \}}
    \ControlFlowTok{return} \KeywordTok{false}\OperatorTok{;}
\NormalTok{  \}}

  \FunctionTok{getStatus}\NormalTok{() \{}
    \KeywordTok{const}\NormalTok{ inv }\OperatorTok{=} \KeywordTok{this}\OperatorTok{.}\AttributeTok{inventory}\OperatorTok{.}\AttributeTok{length} \OperatorTok{\textgreater{}} \DecValTok{0} \OperatorTok{?} \KeywordTok{this}\OperatorTok{.}\AttributeTok{inventory}\OperatorTok{.}\FunctionTok{join}\NormalTok{(}\StringTok{", "}\NormalTok{) }\OperatorTok{:} \StringTok{"empty"}\OperatorTok{;}
    \ControlFlowTok{return} \VerbatimStringTok{\textasciigrave{}}\SpecialCharTok{$\{}\KeywordTok{this}\OperatorTok{.}\AttributeTok{name}\SpecialCharTok{\}}\VerbatimStringTok{ | Health: }\SpecialCharTok{$\{}\KeywordTok{this}\OperatorTok{.}\AttributeTok{health}\SpecialCharTok{\}}\VerbatimStringTok{ | Inventory: }\SpecialCharTok{$\{}\NormalTok{inv}\SpecialCharTok{\}}\VerbatimStringTok{\textasciigrave{}}\OperatorTok{;}
\NormalTok{  \}}
\NormalTok{\}}

\KeywordTok{async} \KeywordTok{function} \FunctionTok{playGame}\NormalTok{() \{}
  \KeywordTok{const}\NormalTok{ \{ askQuestion}\OperatorTok{,}\NormalTok{ close \} }\OperatorTok{=} \PreprocessorTok{require}\NormalTok{(}\StringTok{"./input"}\NormalTok{)}\OperatorTok{;}

  \CommentTok{// Build world}
  \KeywordTok{const}\NormalTok{ cave }\OperatorTok{=} \KeywordTok{new} \FunctionTok{Room}\NormalTok{(}\StringTok{"Dark Cave"}\OperatorTok{,} \StringTok{"You are in a dimly lit cave."}\NormalTok{)}\OperatorTok{;}
  \KeywordTok{const}\NormalTok{ tunnel }\OperatorTok{=} \KeywordTok{new} \FunctionTok{Room}\NormalTok{(}
    \StringTok{"Stone Tunnel"}\OperatorTok{,}
    \StringTok{"A narrow tunnel runs north and south."}\OperatorTok{,}
\NormalTok{  )}\OperatorTok{;}
  \KeywordTok{const}\NormalTok{ clearing }\OperatorTok{=} \KeywordTok{new} \FunctionTok{Room}\NormalTok{(}
    \StringTok{"Sunny Clearing"}\OperatorTok{,}
    \StringTok{"A bright clearing with tall trees."}\OperatorTok{,}
\NormalTok{  )}\OperatorTok{;}

\NormalTok{  cave}\OperatorTok{.}\FunctionTok{addExit}\NormalTok{(}\StringTok{"north"}\OperatorTok{,}\NormalTok{ tunnel)}\OperatorTok{;}
\NormalTok{  cave}\OperatorTok{.}\FunctionTok{addItem}\NormalTok{(}\StringTok{"torch"}\NormalTok{)}\OperatorTok{;}
\NormalTok{  tunnel}\OperatorTok{.}\FunctionTok{addExit}\NormalTok{(}\StringTok{"south"}\OperatorTok{,}\NormalTok{ cave)}\OperatorTok{;}
\NormalTok{  tunnel}\OperatorTok{.}\FunctionTok{addExit}\NormalTok{(}\StringTok{"north"}\OperatorTok{,}\NormalTok{ clearing)}\OperatorTok{;}
\NormalTok{  clearing}\OperatorTok{.}\FunctionTok{addExit}\NormalTok{(}\StringTok{"south"}\OperatorTok{,}\NormalTok{ tunnel)}\OperatorTok{;}
\NormalTok{  clearing}\OperatorTok{.}\FunctionTok{addItem}\NormalTok{(}\StringTok{"magic key"}\NormalTok{)}\OperatorTok{;}

  \KeywordTok{const}\NormalTok{ player }\OperatorTok{=} \KeywordTok{new} \FunctionTok{Player}\NormalTok{(}\StringTok{"Hero"}\OperatorTok{,}\NormalTok{ cave)}\OperatorTok{;}

  \BuiltInTok{console}\OperatorTok{.}\FunctionTok{log}\NormalTok{(}\StringTok{"=== Text Adventure ==="}\NormalTok{)}\OperatorTok{;}
  \BuiltInTok{console}\OperatorTok{.}\FunctionTok{log}\NormalTok{(}\StringTok{"Commands: go [direction], take [item], inventory, look, quit}\SpecialCharTok{\textbackslash{}n}\StringTok{"}\NormalTok{)}\OperatorTok{;}

  \ControlFlowTok{while}\NormalTok{ (}\KeywordTok{true}\NormalTok{) \{}
    \BuiltInTok{console}\OperatorTok{.}\FunctionTok{log}\NormalTok{(player}\OperatorTok{.}\AttributeTok{currentRoom}\OperatorTok{.}\FunctionTok{describe}\NormalTok{())}\OperatorTok{;}
    \BuiltInTok{console}\OperatorTok{.}\FunctionTok{log}\NormalTok{(player}\OperatorTok{.}\FunctionTok{getStatus}\NormalTok{())}\OperatorTok{;}

    \KeywordTok{const}\NormalTok{ input }\OperatorTok{=}\NormalTok{ (}\ControlFlowTok{await} \FunctionTok{askQuestion}\NormalTok{(}\StringTok{"}\SpecialCharTok{\textbackslash{}n}\StringTok{\textgreater{} "}\NormalTok{))}\OperatorTok{.}\FunctionTok{trim}\NormalTok{()}\OperatorTok{.}\FunctionTok{toLowerCase}\NormalTok{()}\OperatorTok{;}
    \KeywordTok{const}\NormalTok{ [command}\OperatorTok{,} \OperatorTok{...}\NormalTok{argParts] }\OperatorTok{=}\NormalTok{ input}\OperatorTok{.}\FunctionTok{split}\NormalTok{(}\SpecialStringTok{/}\SpecialCharTok{\textbackslash{}s+}\SpecialStringTok{/}\NormalTok{)}\OperatorTok{;}
    \KeywordTok{const}\NormalTok{ argument }\OperatorTok{=}\NormalTok{ argParts}\OperatorTok{.}\FunctionTok{join}\NormalTok{(}\StringTok{" "}\NormalTok{)}\OperatorTok{;}

    \BuiltInTok{console}\OperatorTok{.}\FunctionTok{log}\NormalTok{(}\StringTok{""}\NormalTok{)}\OperatorTok{;}

    \ControlFlowTok{switch}\NormalTok{ (command) \{}
      \ControlFlowTok{case} \StringTok{"go"}\OperatorTok{:}
        \ControlFlowTok{if}\NormalTok{ (}\OperatorTok{!}\NormalTok{argument) \{}
          \BuiltInTok{console}\OperatorTok{.}\FunctionTok{log}\NormalTok{(}\StringTok{"Go where?"}\NormalTok{)}\OperatorTok{;}
\NormalTok{        \} }\ControlFlowTok{else} \ControlFlowTok{if}\NormalTok{ (}\OperatorTok{!}\NormalTok{player}\OperatorTok{.}\FunctionTok{move}\NormalTok{(argument)) \{}
          \BuiltInTok{console}\OperatorTok{.}\FunctionTok{log}\NormalTok{(}\StringTok{"You can\textquotesingle{}t go that way."}\NormalTok{)}\OperatorTok{;}
\NormalTok{        \}}
        \ControlFlowTok{break}\OperatorTok{;}
      \ControlFlowTok{case} \StringTok{"take"}\OperatorTok{:}
        \ControlFlowTok{if}\NormalTok{ (player}\OperatorTok{.}\FunctionTok{takeItem}\NormalTok{(argument)) \{}
          \BuiltInTok{console}\OperatorTok{.}\FunctionTok{log}\NormalTok{(}\VerbatimStringTok{\textasciigrave{}You picked up the }\SpecialCharTok{$\{}\NormalTok{argument}\SpecialCharTok{\}}\VerbatimStringTok{.\textasciigrave{}}\NormalTok{)}\OperatorTok{;}
\NormalTok{        \} }\ControlFlowTok{else}\NormalTok{ \{}
          \BuiltInTok{console}\OperatorTok{.}\FunctionTok{log}\NormalTok{(}\VerbatimStringTok{\textasciigrave{}There\textquotesingle{}s no }\SpecialCharTok{$\{}\NormalTok{argument}\SpecialCharTok{\}}\VerbatimStringTok{ here.\textasciigrave{}}\NormalTok{)}\OperatorTok{;}
\NormalTok{        \}}
        \ControlFlowTok{break}\OperatorTok{;}
      \ControlFlowTok{case} \StringTok{"inventory"}\OperatorTok{:}
        \BuiltInTok{console}\OperatorTok{.}\FunctionTok{log}\NormalTok{(}\VerbatimStringTok{\textasciigrave{}Inventory: }\SpecialCharTok{$\{}\NormalTok{player}\OperatorTok{.}\AttributeTok{inventory}\OperatorTok{.}\FunctionTok{join}\NormalTok{(}\StringTok{", "}\NormalTok{) }\OperatorTok{||} \StringTok{"empty"}\SpecialCharTok{\}}\VerbatimStringTok{\textasciigrave{}}\NormalTok{)}\OperatorTok{;}
        \ControlFlowTok{break}\OperatorTok{;}
      \ControlFlowTok{case} \StringTok{"look"}\OperatorTok{:}
        \CommentTok{// Description already shown at top}
        \ControlFlowTok{break}\OperatorTok{;}
      \ControlFlowTok{case} \StringTok{"quit"}\OperatorTok{:}
        \BuiltInTok{console}\OperatorTok{.}\FunctionTok{log}\NormalTok{(}\StringTok{"Thanks for playing!"}\NormalTok{)}\OperatorTok{;}
        \FunctionTok{close}\NormalTok{()}\OperatorTok{;}
        \ControlFlowTok{return}\OperatorTok{;}
      \ControlFlowTok{default}\OperatorTok{:}
        \BuiltInTok{console}\OperatorTok{.}\FunctionTok{log}\NormalTok{(}\VerbatimStringTok{\textasciigrave{}Unknown command: \textquotesingle{}}\SpecialCharTok{$\{}\NormalTok{command}\SpecialCharTok{\}}\VerbatimStringTok{\textquotesingle{}\textasciigrave{}}\NormalTok{)}\OperatorTok{;}
\NormalTok{    \}}
\NormalTok{  \}}
\NormalTok{\}}

\FunctionTok{playGame}\NormalTok{()}\OperatorTok{;}
\end{Highlighting}
\end{Shaded}

\begin{center}\rule{0.5\linewidth}{0.5pt}\end{center}

\subsection{Lesson 10.2: Command-Line To-Do
Manager}\label{lesson-10.2-command-line-to-do-manager}

\begin{Shaded}
\begin{Highlighting}[]
\KeywordTok{const}\NormalTok{ fs }\OperatorTok{=} \PreprocessorTok{require}\NormalTok{(}\StringTok{"fs"}\NormalTok{)}\OperatorTok{;}
\KeywordTok{const}\NormalTok{ \{ askQuestion}\OperatorTok{,}\NormalTok{ close \} }\OperatorTok{=} \PreprocessorTok{require}\NormalTok{(}\StringTok{"./input"}\NormalTok{)}\OperatorTok{;}

\KeywordTok{const}\NormalTok{ FILE\_PATH }\OperatorTok{=} \StringTok{"todos.json"}\OperatorTok{;}

\KeywordTok{function} \FunctionTok{loadTodos}\NormalTok{() \{}
  \ControlFlowTok{try}\NormalTok{ \{}
    \KeywordTok{const}\NormalTok{ data }\OperatorTok{=}\NormalTok{ fs}\OperatorTok{.}\FunctionTok{readFileSync}\NormalTok{(FILE\_PATH}\OperatorTok{,} \StringTok{"utf8"}\NormalTok{)}\OperatorTok{;}
    \ControlFlowTok{return} \BuiltInTok{JSON}\OperatorTok{.}\FunctionTok{parse}\NormalTok{(data)}\OperatorTok{;}
\NormalTok{  \} }\ControlFlowTok{catch}\NormalTok{ (err) \{}
    \ControlFlowTok{if}\NormalTok{ (err}\OperatorTok{.}\AttributeTok{code} \OperatorTok{===} \StringTok{"ENOENT"}\NormalTok{) \{}
      \ControlFlowTok{return}\NormalTok{ []}\OperatorTok{;} \CommentTok{// File doesn\textquotesingle{}t exist yet}
\NormalTok{    \}}
    \ControlFlowTok{if}\NormalTok{ (err }\KeywordTok{instanceof} \BuiltInTok{SyntaxError}\NormalTok{) \{}
      \BuiltInTok{console}\OperatorTok{.}\FunctionTok{log}\NormalTok{(}\StringTok{"Warning: Corrupted todos file. Starting fresh."}\NormalTok{)}\OperatorTok{;}
      \ControlFlowTok{return}\NormalTok{ []}\OperatorTok{;}
\NormalTok{    \}}
    \ControlFlowTok{throw}\NormalTok{ err}\OperatorTok{;}
\NormalTok{  \}}
\NormalTok{\}}

\KeywordTok{function} \FunctionTok{saveTodos}\NormalTok{(todos) \{}
\NormalTok{  fs}\OperatorTok{.}\FunctionTok{writeFileSync}\NormalTok{(FILE\_PATH}\OperatorTok{,} \BuiltInTok{JSON}\OperatorTok{.}\FunctionTok{stringify}\NormalTok{(todos}\OperatorTok{,} \KeywordTok{null}\OperatorTok{,} \DecValTok{2}\NormalTok{))}\OperatorTok{;}
\NormalTok{\}}

\KeywordTok{function} \FunctionTok{addTask}\NormalTok{(todos}\OperatorTok{,}\NormalTok{ description) \{}
\NormalTok{  todos}\OperatorTok{.}\FunctionTok{push}\NormalTok{(\{}
\NormalTok{    description}\OperatorTok{,}
    \DataTypeTok{completed}\OperatorTok{:} \KeywordTok{false}\OperatorTok{,}
    \DataTypeTok{createdAt}\OperatorTok{:} \KeywordTok{new} \BuiltInTok{Date}\NormalTok{()}\OperatorTok{.}\FunctionTok{toISOString}\NormalTok{()}\OperatorTok{,}
\NormalTok{  \})}\OperatorTok{;}
  \BuiltInTok{console}\OperatorTok{.}\FunctionTok{log}\NormalTok{(}\VerbatimStringTok{\textasciigrave{}Task added: "}\SpecialCharTok{$\{}\NormalTok{description}\SpecialCharTok{\}}\VerbatimStringTok{"\textasciigrave{}}\NormalTok{)}\OperatorTok{;}
\NormalTok{\}}

\KeywordTok{function} \FunctionTok{viewTasks}\NormalTok{(todos) \{}
  \ControlFlowTok{if}\NormalTok{ (todos}\OperatorTok{.}\AttributeTok{length} \OperatorTok{===} \DecValTok{0}\NormalTok{) \{}
    \BuiltInTok{console}\OperatorTok{.}\FunctionTok{log}\NormalTok{(}\StringTok{"No tasks in your list."}\NormalTok{)}\OperatorTok{;}
    \ControlFlowTok{return}\OperatorTok{;}
\NormalTok{  \}}
  \BuiltInTok{console}\OperatorTok{.}\FunctionTok{log}\NormalTok{(}\StringTok{"{-}{-}{-} Your To{-}Do List {-}{-}{-}"}\NormalTok{)}\OperatorTok{;}
\NormalTok{  todos}\OperatorTok{.}\FunctionTok{forEach}\NormalTok{((task}\OperatorTok{,}\NormalTok{ index) }\KeywordTok{=\textgreater{}}\NormalTok{ \{}
    \KeywordTok{const}\NormalTok{ status }\OperatorTok{=}\NormalTok{ task}\OperatorTok{.}\AttributeTok{completed} \OperatorTok{?} \StringTok{"[X]"} \OperatorTok{:} \StringTok{"[ ]"}\OperatorTok{;}
    \BuiltInTok{console}\OperatorTok{.}\FunctionTok{log}\NormalTok{(}\VerbatimStringTok{\textasciigrave{}}\SpecialCharTok{$\{}\NormalTok{index }\OperatorTok{+} \DecValTok{1}\SpecialCharTok{\}}\VerbatimStringTok{. }\SpecialCharTok{$\{}\NormalTok{status}\SpecialCharTok{\}}\VerbatimStringTok{ }\SpecialCharTok{$\{}\NormalTok{task}\OperatorTok{.}\AttributeTok{description}\SpecialCharTok{\}}\VerbatimStringTok{\textasciigrave{}}\NormalTok{)}\OperatorTok{;}
\NormalTok{  \})}\OperatorTok{;}
  \BuiltInTok{console}\OperatorTok{.}\FunctionTok{log}\NormalTok{(}\StringTok{"{-}{-}{-}{-}{-}{-}{-}{-}{-}{-}{-}{-}{-}{-}{-}{-}{-}{-}{-}{-}{-}{-}{-}"}\NormalTok{)}\OperatorTok{;}
\NormalTok{\}}

\KeywordTok{function} \FunctionTok{markComplete}\NormalTok{(todos}\OperatorTok{,}\NormalTok{ taskNumber) \{}
  \KeywordTok{const}\NormalTok{ index }\OperatorTok{=} \PreprocessorTok{parseInt}\NormalTok{(taskNumber) }\OperatorTok{{-}} \DecValTok{1}\OperatorTok{;}
  \ControlFlowTok{if}\NormalTok{ (index }\OperatorTok{\textgreater{}=} \DecValTok{0} \OperatorTok{\&\&}\NormalTok{ index }\OperatorTok{\textless{}}\NormalTok{ todos}\OperatorTok{.}\AttributeTok{length}\NormalTok{) \{}
\NormalTok{    todos[index]}\OperatorTok{.}\AttributeTok{completed} \OperatorTok{=} \KeywordTok{true}\OperatorTok{;}
    \BuiltInTok{console}\OperatorTok{.}\FunctionTok{log}\NormalTok{(}\VerbatimStringTok{\textasciigrave{}Completed: "}\SpecialCharTok{$\{}\NormalTok{todos[index]}\OperatorTok{.}\AttributeTok{description}\SpecialCharTok{\}}\VerbatimStringTok{"\textasciigrave{}}\NormalTok{)}\OperatorTok{;}
\NormalTok{  \} }\ControlFlowTok{else}\NormalTok{ \{}
    \BuiltInTok{console}\OperatorTok{.}\FunctionTok{log}\NormalTok{(}\StringTok{"Invalid task number."}\NormalTok{)}\OperatorTok{;}
\NormalTok{  \}}
\NormalTok{\}}

\KeywordTok{function} \FunctionTok{deleteTask}\NormalTok{(todos}\OperatorTok{,}\NormalTok{ taskNumber) \{}
  \KeywordTok{const}\NormalTok{ index }\OperatorTok{=} \PreprocessorTok{parseInt}\NormalTok{(taskNumber) }\OperatorTok{{-}} \DecValTok{1}\OperatorTok{;}
  \ControlFlowTok{if}\NormalTok{ (index }\OperatorTok{\textgreater{}=} \DecValTok{0} \OperatorTok{\&\&}\NormalTok{ index }\OperatorTok{\textless{}}\NormalTok{ todos}\OperatorTok{.}\AttributeTok{length}\NormalTok{) \{}
    \KeywordTok{const}\NormalTok{ removed }\OperatorTok{=}\NormalTok{ todos}\OperatorTok{.}\FunctionTok{splice}\NormalTok{(index}\OperatorTok{,} \DecValTok{1}\NormalTok{)[}\DecValTok{0}\NormalTok{]}\OperatorTok{;}
    \BuiltInTok{console}\OperatorTok{.}\FunctionTok{log}\NormalTok{(}\VerbatimStringTok{\textasciigrave{}Deleted: "}\SpecialCharTok{$\{}\NormalTok{removed}\OperatorTok{.}\AttributeTok{description}\SpecialCharTok{\}}\VerbatimStringTok{"\textasciigrave{}}\NormalTok{)}\OperatorTok{;}
\NormalTok{  \} }\ControlFlowTok{else}\NormalTok{ \{}
    \BuiltInTok{console}\OperatorTok{.}\FunctionTok{log}\NormalTok{(}\StringTok{"Invalid task number."}\NormalTok{)}\OperatorTok{;}
\NormalTok{  \}}
\NormalTok{\}}

\KeywordTok{async} \KeywordTok{function} \FunctionTok{main}\NormalTok{() \{}
  \KeywordTok{let}\NormalTok{ todos }\OperatorTok{=} \FunctionTok{loadTodos}\NormalTok{()}\OperatorTok{;}
  \BuiltInTok{console}\OperatorTok{.}\FunctionTok{log}\NormalTok{(}\StringTok{"=== To{-}Do Manager ==="}\NormalTok{)}\OperatorTok{;}

  \ControlFlowTok{while}\NormalTok{ (}\KeywordTok{true}\NormalTok{) \{}
    \BuiltInTok{console}\OperatorTok{.}\FunctionTok{log}\NormalTok{(}
      \StringTok{"}\SpecialCharTok{\textbackslash{}n}\StringTok{1. Add Task  2. View Tasks  3. Mark Complete  4. Delete  5. Quit"}\OperatorTok{,}
\NormalTok{    )}\OperatorTok{;}
    \KeywordTok{const}\NormalTok{ choice }\OperatorTok{=}\NormalTok{ (}\ControlFlowTok{await} \FunctionTok{askQuestion}\NormalTok{(}\StringTok{"Choice: "}\NormalTok{))}\OperatorTok{.}\FunctionTok{trim}\NormalTok{()}\OperatorTok{;}

    \ControlFlowTok{switch}\NormalTok{ (choice) \{}
      \ControlFlowTok{case} \StringTok{"1"}\OperatorTok{:}
        \KeywordTok{const}\NormalTok{ desc }\OperatorTok{=} \ControlFlowTok{await} \FunctionTok{askQuestion}\NormalTok{(}\StringTok{"Task description: "}\NormalTok{)}\OperatorTok{;}
        \FunctionTok{addTask}\NormalTok{(todos}\OperatorTok{,}\NormalTok{ desc)}\OperatorTok{;}
        \FunctionTok{saveTodos}\NormalTok{(todos)}\OperatorTok{;}
        \ControlFlowTok{break}\OperatorTok{;}
      \ControlFlowTok{case} \StringTok{"2"}\OperatorTok{:}
        \FunctionTok{viewTasks}\NormalTok{(todos)}\OperatorTok{;}
        \ControlFlowTok{break}\OperatorTok{;}
      \ControlFlowTok{case} \StringTok{"3"}\OperatorTok{:}
        \FunctionTok{viewTasks}\NormalTok{(todos)}\OperatorTok{;}
        \KeywordTok{const}\NormalTok{ completeNum }\OperatorTok{=} \ControlFlowTok{await} \FunctionTok{askQuestion}\NormalTok{(}\StringTok{"Task number: "}\NormalTok{)}\OperatorTok{;}
        \FunctionTok{markComplete}\NormalTok{(todos}\OperatorTok{,}\NormalTok{ completeNum)}\OperatorTok{;}
        \FunctionTok{saveTodos}\NormalTok{(todos)}\OperatorTok{;}
        \ControlFlowTok{break}\OperatorTok{;}
      \ControlFlowTok{case} \StringTok{"4"}\OperatorTok{:}
        \FunctionTok{viewTasks}\NormalTok{(todos)}\OperatorTok{;}
        \KeywordTok{const}\NormalTok{ deleteNum }\OperatorTok{=} \ControlFlowTok{await} \FunctionTok{askQuestion}\NormalTok{(}\StringTok{"Task number to delete: "}\NormalTok{)}\OperatorTok{;}
        \FunctionTok{deleteTask}\NormalTok{(todos}\OperatorTok{,}\NormalTok{ deleteNum)}\OperatorTok{;}
        \FunctionTok{saveTodos}\NormalTok{(todos)}\OperatorTok{;}
        \ControlFlowTok{break}\OperatorTok{;}
      \ControlFlowTok{case} \StringTok{"5"}\OperatorTok{:}
        \BuiltInTok{console}\OperatorTok{.}\FunctionTok{log}\NormalTok{(}\StringTok{"Goodbye!"}\NormalTok{)}\OperatorTok{;}
        \FunctionTok{close}\NormalTok{()}\OperatorTok{;}
        \ControlFlowTok{return}\OperatorTok{;}
      \ControlFlowTok{default}\OperatorTok{:}
        \BuiltInTok{console}\OperatorTok{.}\FunctionTok{log}\NormalTok{(}\StringTok{"Invalid choice."}\NormalTok{)}\OperatorTok{;}
\NormalTok{    \}}
\NormalTok{  \}}
\NormalTok{\}}

\FunctionTok{main}\NormalTok{()}\OperatorTok{;}
\end{Highlighting}
\end{Shaded}

\begin{center}\rule{0.5\linewidth}{0.5pt}\end{center}

\section{Unit 11: Capstone Project --- Personal Data Manager in
JavaScript}\label{unit-11-capstone-project-personal-data-manager-in-javascript}

\subsection{Overview}\label{overview-1}

Build a fully functional, accessible command-line personal data manager
using Node.js.

\subsection{Project Options}\label{project-options}

\begin{itemize}
\tightlist
\item
  \textbf{Contact Manager:} name, phone, email, notes
\item
  \textbf{Book Tracker:} title, author, genre, read status
\item
  \textbf{Recipe Collection:} name, ingredients, instructions, cuisine
  type
\item
  \textbf{Budget Tracker:} category, amount, date, description
\end{itemize}

\subsection{Required Features}\label{required-features}

\begin{enumerate}
\def\labelenumi{\arabic{enumi}.}
\tightlist
\item
  \textbf{Add item} --- prompt for all attributes and save.
\item
  \textbf{View all items} --- display clearly with numbers.
\item
  \textbf{Search/filter} --- find items by keyword.
\item
  \textbf{Edit item} --- select and modify specific fields.
\item
  \textbf{Delete item} --- remove by number.
\item
  \textbf{Persistence} --- save/load to JSON file.
\end{enumerate}

\subsection{Project Structure}\label{project-structure}

\begin{verbatim}
my-data-manager/
|--- src/
|   |--- models/
|   |   +--- contact.js      (data model/class)
|   |--- storage/
|   |   +--- dataManager.js  (JSON file I/O)
|   |--- ui/
|   |   +--- consoleUI.js    (menu and display)
|   +--- utils/
|       +--- input.js        (reusable input helper)
|--- data/
|   +--- contacts.json       (data storage)
|--- tests/
|   +--- contact.test.js     (Jest tests)
|--- package.json
|--- README.md
+--- index.js                (entry point)
\end{verbatim}

\subsection{Accessibility Guidelines}\label{accessibility-guidelines}

\begin{enumerate}
\def\labelenumi{\arabic{enumi}.}
\tightlist
\item
  Every prompt must clearly describe what input is expected.
\item
  Print a confirmation message after every action.
\item
  Number all listed items for easy selection.
\item
  Provide clear, specific error messages.
\item
  Include a \texttt{help} command listing all available options.
\item
  Use consistent separators (\texttt{-\/-\/-}) to group output.
\end{enumerate}

\subsection{Assessment Criteria}\label{assessment-criteria}

\subsubsection{Functionality (60\%)}\label{functionality-60}

\begin{itemize}
\tightlist
\item[$\square$]
  All CRUD operations (Create, Read, Update, Delete) work correctly.
\item[$\square$]
  Data persists to JSON between program runs.
\item[$\square$]
  Search works by at least one field.
\item[$\square$]
  Program handles invalid input without crashing.
\end{itemize}

\subsubsection{Code Quality (20\%)}\label{code-quality-20}

\begin{itemize}
\tightlist
\item[$\square$]
  Functions are small with single responsibilities.
\item[$\square$]
  Consistent use of \texttt{const}/\texttt{let}, arrow functions,
  template literals.
\item[$\square$]
  Meaningful names for all variables and functions.
\item[$\square$]
  Code is organized into modules.
\item[$\square$]
  Includes JSDoc comments for key functions.
\end{itemize}

\subsubsection{Accessibility \& UX (15\%)}\label{accessibility-ux-15}

\begin{itemize}
\tightlist
\item[$\square$]
  All prompts are descriptive and clear.
\item[$\square$]
  Confirmation messages after every action.
\item[$\square$]
  Consistent, predictable output format.
\item[$\square$]
  Help system available.
\end{itemize}

\subsubsection{Documentation \& Tests
(5\%)}\label{documentation-tests-5}

\begin{itemize}
\tightlist
\item[$\square$]
  \texttt{README.md} with setup and usage instructions.
\item[$\square$]
  At least 5 Jest unit tests for key functions.
\end{itemize}

\subsection{Deliverables}\label{deliverables}

\begin{enumerate}
\def\labelenumi{\arabic{enumi}.}
\tightlist
\item
  All source code files.
\item
  \texttt{data/} folder with example JSON data.
\item
  \texttt{README.md} including:

  \begin{itemize}
  \tightlist
  \item
    Installation steps: \texttt{npm\ install}
  \item
    Run instructions: \texttt{node\ index.js}
  \item
    Available commands list.
  \item
    Screen reader usage tips.
  \item
    Reflection on challenges faced.
  \end{itemize}
\end{enumerate}

\end{document}